\section*{Computational completeness layer: QAM / MML / MMA}

The present block records the first explicit machine reading of the post-v3.31.0 line.  
In the current v3.34.0-r09 working candidate it should be read together with the local HC/OP interface registry of Definition~\texttt{def:r07-local-language-version-registry} in the active-core PDF, the compression map of Definition~\texttt{def:compression-map-from-inherited-text-to-mml-mma-operations} in the active-core PDF, and the reusable tool normal forms of Definition~\texttt{def:first-reusable-hcop-tool-normal-form} in the active-core PDF. That registry gives the active sublanguage versions actually used by the present OP-entry chain, the compression map gives the active audit reading for inherited HC/OP text, and the r09 tool layer gives the controlled invocation form:
\[
  \mathrm{QAM}_{\mathrm{HCOP}}^{\mathrm{loc}}
  \to
  \mathrm{MML}_{\mathrm{HCOP}}^{\mathrm{loc}}
  \to
  \mathrm{MMA}_{\mathrm{HCOP}}^{\mathrm{loc}}.
\]
The broader computational chapter remains a reserve architecture; the r07 interface is the narrow executable slice used for HC-routed OP-entry.

Its purpose is to separate three layers which were already implicit in the project architecture:
\begin{enumerate}[leftmargin=2em]
  \item QAM as the finite tagged object language;
  \item MML as the finite machine / instruction language;
  \item MMA as the matrix-realized execution layer.
\end{enumerate}

The block proceeds in three steps. First, it isolates finite machine completeness and finite universality (the Z3-style regime). Second, it exhibits the unbounded extension mechanism through extensible address towers and fresh-address allocation. Third, it formulates stage-stable universality and the resulting architectural Turing-style completeness reading.

This computational layer is to be read \emph{subordinate} to the earlier strength hierarchy
\[
\mathbf S_1\to \mathbf S_2\to \mathbf S_3(\mathcal R)\to \mathbf S_4(\mathcal R)\to \mathbf S_5
\]
from Sections~19U--19Z and the subsequent machine-status decoder. Accordingly, the present section does not by itself upgrade the unconditional public status of the line beyond \(\mathbf S_1\). Rather, it records an \emph{internal machine-reading package}: finite machine completeness and finite universality sharpen the internal route toward \(\mathbf S_2\), while explicit fresh-address extension and stage-stable universality formulate the internal architectural ingredients that would be needed later on any conditional route toward \(\mathbf S_4(\mathcal R)\) and \(\mathbf S_5\).

\begin{definition}[Finite machine store]\label{def:finite-machine-store-computational-layer}
A \emph{finite machine store} is a finite tagged family
\[
\mathfrak{S}_{\mathrm{fin}}
=
\{(a_i,x_i)\}_{i=1}^{N}
\]
where each \(a_i\) is an address tag and each \(x_i\) is a stored QAM-object.
\end{definition}

\begin{definition}[Finite machine state]\label{def:finite-machine-state-computational-layer}
A \emph{finite machine state} is a tuple
\[
\Sigma
=
\bigl(
\mathfrak{S}_{\mathrm{fin}},
p,
q,
r
\bigr),
\]
where:
\begin{enumerate}[leftmargin=2em]
    \item \(\mathfrak{S}_{\mathrm{fin}}\) is a finite machine store;
    \item \(p\) is the current program position;
    \item \(q\) is the current control state;
    \item \(r\) is the current route-faithful residual/control register.
\end{enumerate}
\end{definition}

\begin{definition}[QAM-object layer]\label{def:qam-object-layer-computational-layer}
The \emph{QAM-object layer} is the finite tagged object language in which the machine store contents \(x_i\) are represented.

A QAM-object may be:
\begin{enumerate}[leftmargin=2em]
    \item a scalar-like tagged object;
    \item a packet-like tagged object;
    \item a tuple/block object;
    \item a routed or normalized residual object;
    \item any finite object built from the current QAM grammar by the allowed formation rules.
\end{enumerate}
\end{definition}

\begin{remark}[QAM as object language]\label{rem:qam-as-object-language-computational-layer}
QAM is not yet used here as a new infinite set theory. At this stage it is used as the finite object language in which machine-addressable project data are encoded.
\end{remark}

\begin{definition}[MML instruction alphabet]\label{def:mml-instruction-alphabet-computational-layer}
The \emph{MML instruction alphabet} is a finite set
\[
\mathcal{I}_{\mathrm{MML}}
\]
of primitive machine instructions, divided into the following classes:
\begin{enumerate}[leftmargin=2em]
    \item \textbf{read/write instructions:}
    \[
    \mathrm{READ}(a),\qquad \mathrm{WRITE}(a,\xi);
    \]
    \item \textbf{routing/transfer instructions:}
    \[
    \mathrm{ROUTE}(\xi),\qquad \mathrm{PROJECT}(\xi),\qquad \mathrm{NORMALIZE}(\xi);
    \]
    \item \textbf{comparison/test instructions:}
    \[
    \mathrm{TEST}_{=},\qquad \mathrm{TEST}_{<},\qquad \mathrm{TEST}_{\mathrm{adm}};
    \]
    \item \textbf{branching/control instructions:}
    \[
    \mathrm{SELECT}(b_1,\dots,b_k),\qquad \mathrm{BRANCH}(j),\qquad \mathrm{RETURN};
    \]
    \item \textbf{cycle/termination instructions:}
    \[
    \mathrm{STEP},\qquad \mathrm{LOOP},\qquad \mathrm{HALT}.
    \]
\end{enumerate}
\end{definition}

\begin{remark}[MML as finite machine language]\label{rem:mml-as-finite-machine-language-computational-layer}
MML is the first finite machine language layer of the project. Its purpose is not yet unrestricted computation, but finite programmable symbolic control on route-faithful QAM-data.
\end{remark}

\begin{definition}[Finite MML program]\label{def:finite-mml-program-computational-layer}
A \emph{finite MML program} is a finite word
\[
P=(I_1,\dots,I_m),\qquad I_j\in\mathcal{I}_{\mathrm{MML}}.
\]
\end{definition}

\begin{definition}[MMA execution layer]\label{def:mma-execution-layer-computational-layer}
The \emph{MMA execution layer} is the matrix-realized execution semantics of MML on finite machine states.

Concretely, an MMA-step is a map
\[
\mathcal{E}_{\mathrm{MMA}}:(\Sigma,I)\longmapsto \Sigma^{+},
\]
sending a finite machine state \(\Sigma\) and one MML instruction \(I\) to the updated state \(\Sigma^{+}\), in a way compatible with the current Millennium-matrix / Heisenberg-compensator realization of the route-faithful control structure.
\end{definition}

\begin{remark}[Why MMA is the right execution layer]\label{rem:why-mma-is-the-right-execution-layer-computational-layer}
MML gives the symbolic instruction form. MMA gives the executable matrix-side realization. Thus QAM, MML, and MMA correspond respectively to object language, machine language, and execution dynamics.
\end{remark}

\begin{definition}[Finite branch emulation]\label{def:finite-branch-emulation-computational-layer}
A \emph{finite branch emulation rule} is any route-faithful realization of conditional continuation by means of:
\begin{enumerate}[leftmargin=2em]
    \item a test instruction;
    \item a finite selector family;
    \item a branch-resolution instruction;
    \item an updated program position and control state.
\end{enumerate}
\end{definition}

\begin{remark}[Why explicit jump is not required]\label{rem:why-explicit-jump-is-not-required-computational-layer}
For the finite completeness level pursued here, it is not necessary to require a primitive classical jump instruction. It is sufficient that the machine language can emulate conditional continuation by a finite branch-resolution mechanism.
\end{remark}

\begin{definition}[Finite machine completeness]\label{def:finite-machine-completeness-computational-layer}
The triple
\[
(\mathrm{QAM},\mathrm{MML},\mathrm{MMA})
\]
is called \emph{finitely machine-complete} if:
\begin{enumerate}[leftmargin=2em]
    \item every finite project-relevant symbolic state can be encoded as a finite QAM-store;
    \item every finite admissible symbolic transformation on such states can be represented by a finite MML program;
    \item every such finite MML program is executable in the MMA layer;
    \item conditional continuation is realizable by finite branch emulation;
    \item halting and finite return are representable on the same machine level.
\end{enumerate}
\end{definition}

\begin{definition}[Finite address space]\label{def:finite-address-space-computational-layer}
A \emph{finite address space} is a finite nonempty set
\[
\mathfrak{A}_{\mathrm{fin}}=\{a_1,\dots,a_N\}
\]
of address tags.
\end{definition}

\begin{definition}[Finite machine store as address map]\label{def:finite-machine-store-as-address-map-computational-layer}
A \emph{finite machine store} on the address space
\[
\mathfrak{A}_{\mathrm{fin}}
\]
is a map
\[
\mathsf{Mem}:\mathfrak{A}_{\mathrm{fin}}\longrightarrow \mathsf{QAM}_{\mathrm{fin}},
\]
where \(\mathsf{QAM}_{\mathrm{fin}}\) denotes the class of finite QAM-objects admissible at machine level.
\end{definition}

\begin{definition}[Store read operation]\label{def:store-read-operation-computational-layer}
For \(a\in\mathfrak{A}_{\mathrm{fin}}\), the \emph{store read operation} is
\[
\mathsf{READ}_{\mathsf{Mem}}(a):=\mathsf{Mem}(a).
\]
\end{definition}

\begin{definition}[Store write operation]\label{def:store-write-operation-computational-layer}
For \(a\in\mathfrak{A}_{\mathrm{fin}}\) and \(\xi\in\mathsf{QAM}_{\mathrm{fin}}\), the \emph{store write operation} is the updated store
\[
\mathsf{WRITE}_{\mathsf{Mem}}(a,\xi)=\mathsf{Mem}^{+},
\]
defined by
\[
\mathsf{Mem}^{+}(a)=\xi,
\qquad
\mathsf{Mem}^{+}(b)=\mathsf{Mem}(b)
\;\text{for all } b\neq a.
\]
\end{definition}

\begin{definition}[Address register]\label{def:address-register-computational-layer}
An \emph{address register} is a distinguished element \(\alpha\in\mathfrak{A}_{\mathrm{fin}}\) serving as the currently selected address for machine-level read/write access.
\end{definition}

\begin{definition}[Control register]\label{def:control-register-computational-layer}
A \emph{control register} is a finite tuple
\[
\mathcal{R}_{\mathrm{ctrl}}=(p,q,r,\alpha),
\]
where \(p\) is the current program position, \(q\) the current control state, \(r\) the current route-faithful residual/control datum, and \(\alpha\) the current address register.
\end{definition}

\begin{definition}[Finite machine configuration]\label{def:finite-machine-configuration-computational-layer}
A \emph{finite machine configuration} is a pair
\[
\mathcal{C}=(\mathsf{Mem},\mathcal{R}_{\mathrm{ctrl}})
\]
consisting of a finite machine store and a control register.
\end{definition}

\begin{definition}[Finite universality]\label{def:finite-universality-computational-layer}
The stack
\[
(\mathrm{QAM},\mathrm{MML},\mathrm{MMA})
\]
is called \emph{finitely universal} if every finite admissible symbolic task on finite machine configurations admits a finite MML realization executable in MMA.
\end{definition}

\begin{definition}[Extensible address tower]\label{def:extensible-address-tower-computational-layer}
An \emph{extensible address tower} is a nested family
\[
\mathfrak{A}_1\subset\mathfrak{A}_2\subset\mathfrak{A}_3\subset\cdots
\]
of finite address spaces such that each inclusion preserves the previously existing address tags and every finite set of address demands is realized at some stage.
\end{definition}

\begin{definition}[Store extension map]\label{def:store-extension-map-computational-layer}
Given two stages \(\mathfrak{A}_n\subset\mathfrak{A}_{n+1}\), a \emph{store extension map} is a map
\[
\mathrm{Ext}_{n}^{n+1}:\mathsf{Mem}_n\longmapsto \mathsf{Mem}_{n+1}
\]
such that old address contents are preserved and every new address is initialized by a fixed admissible default QAM-object.
\end{definition}

\begin{definition}[Fresh tag]\label{def:fresh-tag-computational-layer}
Let \(\mathfrak{A}_{\mathrm{fin}}\) be the currently active finite address/tag carrier set. A \emph{fresh tag} is a tag
\[
a_{\mathrm{new}}\notin\mathfrak{A}_{\mathrm{fin}}.
\]
\end{definition}

\begin{definition}[Fresh-address realization rule]\label{def:fresh-address-realization-rule-computational-layer}
A \emph{fresh-address realization rule} assigns to every fresh tag \(a_{\mathrm{new}}\) an extended finite address space
\[
\mathfrak{A}_{\mathrm{fin}}^{+}=\mathfrak{A}_{\mathrm{fin}}\cup\{a_{\mathrm{new}}\}.
\]
\end{definition}

\begin{definition}[Default initialization object]\label{def:default-initialization-object-computational-layer}
A \emph{default initialization object} is a distinguished finite QAM-object
\[
\mathbf{0}_{\mathrm{QAM}}\in\mathsf{QAM}_{\mathrm{fin}}
\]
used to initialize newly created addresses before any further write/update takes place.
\end{definition}

\begin{definition}[Fresh-address store extension]\label{def:fresh-address-store-extension-computational-layer}
Let
\[
\mathsf{Mem}:\mathfrak{A}_{\mathrm{fin}}\to\mathsf{QAM}_{\mathrm{fin}}
\]
be the current finite machine store, and let \(a_{\mathrm{new}}\notin\mathfrak{A}_{\mathrm{fin}}\) be a fresh tag.

The \emph{fresh-address store extension} is the extended store
\[
\mathsf{Mem}^{+}:\mathfrak{A}_{\mathrm{fin}}^{+}\to\mathsf{QAM}_{\mathrm{fin}}
\]
defined by
\[
\mathsf{Mem}^{+}(a)=\mathsf{Mem}(a)
\quad
\text{for all } a\in\mathfrak{A}_{\mathrm{fin}},
\]
and
\[
\mathsf{Mem}^{+}(a_{\mathrm{new}})=\mathbf{0}_{\mathrm{QAM}}.
\]
\end{definition}

\begin{definition}[Old-address semantic preservation]\label{def:old-address-semantic-preservation-computational-layer}
A fresh-address extension satisfies \emph{old-address semantic preservation} if every admissible instruction that only refers to addresses in the old domain has exactly the same effect on those old addresses before and after extension.
\end{definition}

\begin{definition}[Stage-stable universality]\label{def:stage-stable-universality-computational-layer}
The stack
\[
(\mathrm{QAM},\mathrm{MML},\mathrm{MMA})
\]
is called \emph{stage-stably universal} if every extension-compatible stage-stable task family admits a stable unbounded realization by one stage-stable finite program.
\end{definition}

\begin{remark}[Computational status reading]\label{rem:computational-status-reading-computational-layer}
At the present stage, the project computational architecture should be read in three levels:
\begin{enumerate}[leftmargin=2em]
    \item finite machine completeness;
    \item finite universality (Z3-style regime);
    \item stage-stable universality together with fresh-address extension, yielding an architectural Turing-style regime.
\end{enumerate}
\end{remark}

\begin{theorem}[First finite machine completeness reading theorem]\label{thm:first-finite-machine-completeness-computational-layer}
Assume:
\begin{enumerate}[leftmargin=2em]
    \item finite machine states are encoded in the QAM-object layer;
    \item finite admissible symbolic transformations are representable in MML;
    \item the MMA execution layer executes finite MML programs route-faithfully;
    \item finite branch emulation is available.
\end{enumerate}
Then the triple
\[
(\mathrm{QAM},\mathrm{MML},\mathrm{MMA})
\]
supports a finitely machine-complete internal reading.
\end{theorem}

\begin{proof}
Finite tagged QAM stores supply the data layer, finite MML programs supply the symbolic control layer, MMA supplies executable realization, and branch emulation supplies finite conditional continuation. Hence the machine architecture is complete at the finite programmable level.
\end{proof}

\begin{theorem}[Finite address/store layer theorem]\label{thm:finite-address-store-completeness-computational-layer}
Assume:
\begin{enumerate}[leftmargin=2em]
    \item the QAM-object layer encodes finite machine stores;
    \item the MML instruction alphabet contains finite read/write, test, branch, loop, and halt instructions;
    \item the MMA execution layer executes address-valid MML instructions on finite machine configurations;
    \item finite branch emulation is available.
\end{enumerate}
Then the project architecture supports an explicit complete finite address/store machine layer.
\end{theorem}

\begin{proof}
Finite stores are explicit address maps, read/write closure holds on the same finite address domain, finite control is preserved under address-valid execution, and finite conditional control is realizable without leaving the finite regime. Thus the address/store layer is complete.
\end{proof}

\begin{theorem}[First finite universality reading theorem]\label{thm:first-finite-universality-computational-layer}
Assume:
\begin{enumerate}[leftmargin=2em]
    \item the strengthened finite machine completeness theorem holds;
    \item every finite symbolic task decomposes into finitely many canonical finite programming primitives;
    \item finite branch emulation is available.
\end{enumerate}
Then the stack
\[
(\mathrm{QAM},\mathrm{MML},\mathrm{MMA})
\]
supports a finite-universality reading on the internal machine layer.
\end{theorem}

\begin{proof}
Every finite symbolic task can be decomposed into finite read/write, routing, projection, normalization, test, selection, branch, loop, and halt primitives; these are representable in MML and executable in MMA. Hence every finite admissible symbolic task admits a finite realization.
\end{proof}

\begin{theorem}[Fresh tag = fresh address reading theorem]\label{thm:fresh-tag-equals-fresh-address-computational-layer}
Assume:
\begin{enumerate}[leftmargin=2em]
    \item the architecture already contains an implicit extensible address-tower pattern;
    \item fresh tags are admissibly generable;
    \item fresh-address realization holds;
    \item default initialization is available;
    \item old-address semantic preservation holds.
\end{enumerate}
Then the current QAM/MML/MMA architecture supports an explicit extensible-address reading.
\end{theorem}

\begin{proof}
Fresh-address realization turns new tags into genuine new addresses, default initialization turns those addresses into usable machine cells, and old-address semantic preservation makes the extension conservative on the previous stage. Iterating this process makes the address tower explicit.
\end{proof}

\begin{theorem}[First stage-stable universality reading theorem]\label{thm:first-stage-stable-universality-computational-layer}
Assume:
\begin{enumerate}[leftmargin=2em]
    \item the first finite universality theorem holds;
    \item the fresh tag = fresh address theorem holds;
    \item uniform instruction persistence holds;
    \item uniform execution persistence holds;
    \item finite branch emulation persists along stage extension.
\end{enumerate}
Then the stack
\[
(\mathrm{QAM},\mathrm{MML},\mathrm{MMA})
\]
supports a stage-stably universal machine reading.
\end{theorem}

\begin{proof}
Finite universality gives programmable realization on each finite stage. Fresh-address extension makes store growth explicit and conservative. Uniform instruction and execution persistence keep the same finite program semantics valid on larger stages, and persistent branch emulation stabilizes conditional control. Therefore extension-compatible stage-stable task families admit stable unbounded realizations.
\end{proof}

\begin{theorem}[Conditional architectural Turing-style completeness reading theorem]\label{thm:architectural-turing-style-completeness-computational-layer}
Assume the first stage-stable universality reading theorem. Then the stack
\[
(\mathrm{QAM},\mathrm{MML},\mathrm{MMA})
\]
supports a conditional architectural Turing-style completeness reading along the explicit stage-stable route singled out by the present computational layer.
\end{theorem}

\begin{proof}
By stage-stable universality, the architecture supports programmable finite symbolic control, explicit unbounded store growth, stable finite program semantics across arbitrarily large stages, and persistent conditional control. These are exactly the architectural ingredients isolated for a Turing-style machine reading. Because the older Sections~19U--19Z still distinguish unconditional status \(\mathbf S_1\) from stronger conditional route-dependent upgrades toward \(\mathbf S_4(\mathcal R)\) and \(\mathbf S_5\), the present conclusion is recorded only as a conditional architectural reading rather than as a new unconditional global machine-status theorem.
\end{proof}

\begin{corollary}[Computational status embedding rule]\label{cor:computational-status-upgrade-rule-computational-layer}
The project computational status should now be embedded into the older strength hierarchy in the following disciplined way:
\begin{enumerate}[leftmargin=2em]
    \item unconditionally, the public machine-status reading of the line remains at \(\mathbf S_1\) unless the external route certificates of Sections~19U--19Z are invoked;
    \item internally, the present block sharpens the machine reading from a bare \(\mathbf S_1\) interpretation toward a finitely machine-complete / finitely universal (Z3-style) package, i.e. toward the primitive-core side of \(\mathbf S_2\);
    \item explicitly extensible fresh-address and stage-stable universality readings isolate the architectural ingredients that would later be needed on a conditional route toward \(\mathbf S_4(\mathcal R)\) and \(\mathbf S_5\), but they do not by themselves replace the earlier route-dependent certification ladder.
\end{enumerate}
\end{corollary}

\begin{proof}
Immediate from the preceding completeness theorems.
\end{proof}

\begin{remark}[No parallel machine-status hierarchy]\label{rem:no-parallel-machine-status-hierarchy-computational-layer}
This computational block should therefore not be read as introducing a second independent machine-status ladder beside Sections~19U--19Z. Its role is narrower and more local: it explicates the QAM/MML/MMA machine package that underlies the already existing hierarchy and sharpens the internal architectural route toward stronger conditional computational claims, while leaving the unconditional public status decoder unchanged.
\end{remark}

\subsection*{Version and DOI record}
The table below records the currently preserved Companion chain as a normalized release and working-state register. Zenodo releases are listed together with a publication-date column; earlier intermediate records and local working witnesses are retained so that the release chronology remains reconstructible even where the register groups recent working states before older archive rows.

The following chronology block is the active normalized release register for the present v\docversion\ release. The concept DOI remains the stable title-page pointer; version-specific DOI claims for public releases and pre-release freeze witnesses are recorded explicitly in this chronology.

{\small
\renewcommand{\arraystretch}{1.12}
\setlength{\LTleft}{0pt}
\setlength{\LTright}{0pt}
\begin{longtable}{|p{0.14\textwidth}|p{0.17\textwidth}|p{0.15\textwidth}|p{0.42\textwidth}|}
\hline
\textbf{Version / record} & \textbf{Status} & \textbf{Publication date} & \textbf{DOI / note} \\
\hline
\endfirsthead
\hline
\textbf{Version / record} & \textbf{Status} & \textbf{Publication date} & \textbf{DOI / note} \\
\hline
\endhead
Concept DOI & concept record & ongoing pointer & \href{https://doi.org/10.5281/zenodo.19210728}{10.5281/zenodo.19210728} \\
\hline
v3.28.0 & public Zenodo release & 2026-04-20 & \href{https://doi.org/10.5281/zenodo.19648743}{10.5281/zenodo.19648743} --- earlier public release; completes the deferred tool/QAM/decoder layer and fixes the route-control infrastructure that the present release inherits as its public baseline \\
\hline
v3.28.1-r39 & pre-release working continuation & 2026-04-20 & local working state above the public v3.28.0 release; assembled the certificate-aware extension/routing/probe layer, local retained-corridor closure of the live OP~1 and OP~5 Step~(B) kernels, QAM/decoder shadow reattachment, certificate-aware Monte-Carlo discipline, and restored full historical preservation; carried forward here as the final pre-release working witness for v3.29.0 \\
\hline
v3.29.0 & public Zenodo release & 2026-04-20 & \href{https://doi.org/10.5281/zenodo.19670829}{10.5281/zenodo.19670829} --- prior public release; promotes the certificate-aware extension/routing/probe layer, finite decision/action/word routing, matrix/crystal/field readout tags, admissibility and routing discipline, QAM/decoder shadow reattachment, certificate-aware Monte-Carlo discipline, and local retained-corridor closure of the live OP~1 and OP~5 Step~(B) kernels while preserving the full historical witness in Part~II \\
\hline
v3.29.1-r32 & release-freeze working continuation & 2026-04-21 & local freeze witness above the public v3.29.0 release; completes the conservative upgrade-plus-probe integration line through probe-compatible packets, work queues, executable conservative programs, compact readout/output/intake surfaces, OP-facing request/fulfillment/aggregation/dispatch/program/cycle/return logic, a completed second downstream loop, and an explicit release-freeze stabilization point; no version-specific DOI assigned \\
\hline
v3.30.0 & public Zenodo release & 2026-04-21 & \href{https://doi.org/10.5281/zenodo.19682951}{10.5281/zenodo.19682951} --- prior public release; promotes the conservative upgrade-plus-probe integration line above the public v3.29.0 release, including probe-compatible conservative upgrade packets, finite probe-ranked work queues, queue-extracted working programs, route-faithful lead tasks, executable conservative packets, compact conservative readout/output/intake surfaces, OP-facing request/fulfillment/aggregation/dispatch/program/cycle/return logic, a completed second downstream loop, and an explicit release-freeze stabilization point while preserving the full historical witness in Part~II \\
\hline
v3.31.0 & public Zenodo release & 2026-04-24 & \href{https://doi.org/10.5281/zenodo.19699827}{10.5281/zenodo.19699827} --- prior public release; freezes the groundwork-complete pre-OP frontier synthesis above the public v3.30.0 baseline and publishes the full monolith with the integrated coupled prime-lattice trigger correction and the QAM/MML/MMA machine-realization layer while preserving the full historical witness in Part~II \\
\hline
v3.32.0 & public Zenodo release & 2026-04-24 & \href{https://doi.org/10.5281/zenodo.19704489}{10.5281/zenodo.19704489} --- prior public release; current release final built above the public v3.31.0 baseline; integrates the first machine-fused coupled OP-entry transition line, the coupled prime-lattice machine realization block, and the strengthened finite-to-stage-stable universality reading while preserving the full historical witness in Part~II \\
v3.33.0 & public Zenodo release & 2026-04-24 & \href{https://doi.org/10.5281/zenodo.19716926}{10.5281/zenodo.19716926} --- prior public release; built above the public v3.32.0 baseline; carries the HC-captured local continuation line through lock-corridor capture, minimal HC-lock step, HC menu scan, plausible HC candidate, micro-lock update, and first local HC-capture chain while preserving the full historical witness in Part~II \\
\hline
v3.34.0 & public Zenodo release & 2026-04-25 & \href{https://doi.org/10.5281/zenodo.19746821}{10.5281/zenodo.19746821} --- prior public split-bundle release; promotes the locally audit-clean v3.34.0-r27 line into a reader-facing active core plus preserved historical witness bundle, adds the QAM/MML/MMA store-array layer, frontmatter reader architecture, and Sphaera array mirror/stripe/parity-residual infrastructure while preserving the non-closure status of OP~1/4 and OP~5(B) \\
\hline
v3.35.0 & public Zenodo release & 2026-04-25 & \href{https://doi.org/10.5281/zenodo.19750674}{10.5281/zenodo.19750674} --- current public pre-v4 QAM/tensor preparation release; adds the short/extended abstract architecture, Hilbert-hotel QAM reindexing plan for later v4, guarded mixed QAM readouts, optional binary mask registry, conservative two-sided structural tensor grammar, role-erasure and mixed-readout normal forms, Boolean selector grammar, and release-final DOI packaging while preserving the non-closure status of OP~1/4 and OP~5(B) \\
\hline
v3.34.0-r17 & local working candidate & 2026-04-24 & no Zenodo DOI assigned --- executes the first dispatched handoff-guard packet \(\mathsf W_{\mathrm{handoff}}^{\mathrm{guard}}\) after a stable lock/capture return, installs the handoff-guard defect register \(\Delta_{\mathrm{handoff}}^{\mathrm{guard}}\), and returns exactly one of handoff-ready source, bounded handoff-guard repair source, or smaller frontier source without manufacturing an OP closure certificate \\
v3.34.0-r18 & local working candidate & 2026-04-24 & no Zenodo DOI assigned --- executes the first dispatched OP-normalization/type packet \(\mathsf W_{\mathrm{norm}}^{\mathrm{type}}\) after a handoff-ready return, installs the normalization/type-coherence defect register \(\Delta_{\mathrm{norm}}^{\mathrm{type}}\), and returns exactly one of route/rank-ready typed OP input, bounded normalization/type repair source, or smaller frontier source without manufacturing an OP closure certificate \\
v3.34.0-r19 & local working candidate & 2026-04-24 & no Zenodo DOI assigned --- executes the first dispatched route/rank packet \(\mathsf W_{\mathrm{route}}^{\mathrm{rank}}\) after a normalization-ready return, installs the route/rank defect register \(\Delta_{\mathrm{route}}^{\mathrm{rank}}\), and returns exactly one of comparison-ready routed OP task, bounded route/rank repair source, or smaller frontier source without manufacturing an OP closure certificate \\
v3.34.0-r20 & local working candidate & 2026-04-24 & no Zenodo DOI assigned --- executes the first dispatched comparison/defect packet \(\mathsf W_{\mathrm{cmp}}^{\Delta}\) after a route/rank-ready return, installs the comparison/defect register \(\Delta_{\mathrm{cmp}}^{\Delta}\), and returns exactly one of guarded bind candidate, guarded split candidate, or smaller comparison frontier without manufacturing an OP closure certificate \\
v3.34.0-r21 & local working candidate & 2026-04-24 & no Zenodo DOI assigned --- executes the first dispatched resolve/return packet \(\mathsf W_{\mathrm{resolve}}^{\mathrm{next}}\) after a guarded comparison return, installs the resolve/return defect register \(\Delta_{\mathrm{resolve}}^{\mathrm{next}}\), and returns exactly one of guarded bind-work continuation, guarded split-work continuation, or smaller comparison-frontier descent without manufacturing an OP closure certificate \\
v3.34.0-r22 & local working candidate & 2026-04-24 & no Zenodo DOI assigned --- executes the first continuation-queue packet after resolve/return, installs the continuation-queue defect register \(\Delta_{\mathrm{queue}}^{\mathrm{cont}}\), and returns exactly one of a guarded common-origin bind queue, guarded ranked split queue, or protected frontier-descent queue without manufacturing an OP closure certificate \\
v3.34.0-r23 & local working candidate & 2026-04-24 & no Zenodo DOI assigned --- executes the first queue-head selection packet after continuation queue, installs the queue-head selection defect register \(\Delta_{\mathrm{head}}^{\mathrm{select}}\), and returns exactly one of a guarded common-origin-check head, guarded branch-rank-check head, or guarded frontier-descent-check head while retaining the finite tail and without manufacturing an OP closure certificate \\
v3.34.0-r24 & local working candidate & 2026-04-24 & no Zenodo DOI assigned --- executes the first selected queue-head execution packet, installs the selected queue-head execution defect register \(\Delta_{\mathrm{head}}^{\mathrm{exec}}\), and returns exactly one of an executed common-origin-check surface, executed branch-rank-check surface, or executed frontier-descent-check surface while retaining the finite queue tail and without manufacturing an OP closure certificate \\
v3.34.0-r25 & local working candidate & 2026-04-24 & no Zenodo DOI assigned --- executes the first executed-head tail-entry packet \(\mathsf W_{\mathrm{tail}}^{\mathrm{entry}}\), installs the tail-entry defect \(\Delta_{\mathrm{tail}}^{\mathrm{entry}}\), and returns exactly one of a guarded bind-tail, split-tail, or frontier-tail entry without manufacturing an OP closure certificate \\
v3.34.0-r26 & local audit-freeze candidate & 2026-04-24 & no Zenodo DOI assigned --- installs the first audit-freeze ledger \(F_{\mathrm{HCOP}}^{\mathrm{audit}}\) with audit-freeze defect \(\Delta_{\mathrm{audit}}^{\mathrm{freeze}}\), recording baseline, DOI state, local spine, tail-entry output, preservation obligations, reference/compilation audit data, terminal bibliography, historical witness, and the no-false-certificate marker \\
v3.34.0-r27 & local audit-freeze candidate & 2026-04-24 & no Zenodo DOI assigned --- installs the audit-freeze stabilization object \(F_{\mathrm{freeze}}^{\mathrm{audit}}\), permitting only audit, mechanical repair, metadata repair, release packaging, and Zenodo-description preparation before further theorem expansion \\
\hline
v3.34.0-r15 & local working candidate & 2026-04-24 & no Zenodo DOI assigned --- turns the r14 tail-normalized HC/OP proof skeleton into the first controlled skeleton-to-workstep dispatch object \(\mathsf D_{\mathrm{HCOP}}^{\mathrm{skel}}\), assigning finite admissible next work packets to lock/capture, handoff, normalization, routing, comparison, and resolution nodes while preserving dependency order, guards, witnesses, route tags, theorematic status, and the no-false-certificate rule \\
v3.34.0-r16 & local working candidate & 2026-04-24 & no Zenodo DOI assigned --- executes the first dispatched lock/capture stabilization packet \(\mathsf W_{\mathrm{lock}}^{\mathrm{stab}}\), installs the lock-stability defect register \(\Delta_{\mathrm{lock}}^{\mathrm{stab}}\), and returns exactly one of handoff-ready stable source, bounded local repair source, or smaller frontier source without manufacturing an OP closure certificate \\
\hline
v3.34.0-r14 & local working candidate & 2026-04-24 & no Zenodo DOI assigned --- extracts the first tail-normalized HC/OP proof skeleton \(\mathsf S_{\mathrm{HCOP}}^{\mathrm{tail}}\), organizing the active r01--r04 spine as LOCK/CAPTURE, HANDOFF, NORMALIZE, ROUTE, COMPARE, and RESOLVE source nodes plus later guarded citation tails while preserving dependency order and preventing false OP-closure certificates \\
\hline
v3.34.0-r13 & local working candidate & 2026-04-24 & no Zenodo DOI assigned --- turns the r12 ledger into the first guarded-reduction pilot normal form \(\mathrm{GRCite}_{\mathrm{HCOP}}(\ell;G,W,S)\), allowing only later repeated HC/OP tails to be rewritten as guarded tool citations while first occurrences, local guards, witnesses, certificate tests, proof-critical inequalities, theorematic status, and historical witness material remain explicit \\
\hline
v3.34.0-r12 & local working candidate & 2026-04-24 & no Zenodo DOI assigned --- applies the r11 guarded-reduction protocol for the first time by installing the finite ledger \(\mathcal L_{\mathrm{HCOP}}^{(1)}\) over the active HC/OP proof spine, allowing later repeated lock/capture, handoff, normalization, route/frontier, paired-comparison, probe-diagnostic, and archive-to-tool tails to be cited in guarded normal form while keeping first occurrences, local guards, certificate tests, proof-critical inequalities, theorematic status, and historical witness material explicit \\
\hline
v3.34.0-r11 & local working candidate & 2026-04-24 & no Zenodo DOI assigned --- keeps the r10 tool-contract application layer and adds the first tool-guarded reduction protocol \(\mathcal R_{\mathrm{HCOP}}^{(1)}\), defining reducible and non-reducible HC/OP occurrences, guarded citation normal form, and the partial simplification map \(\mathsf{Red}_{\mathrm{HCOP}}^{(1)}\) while preserving proof content, first definitions, local guards, certificate tests, theorematic status, and the full historical witness in Part~II \\
\hline
v3.34.0-r10 & local working candidate & 2026-04-24 & no Zenodo DOI assigned --- applies the r09 reusable tool contracts back to the active r01--r04 HC/OP proof spine as a conservative cross-reference layer \(\mathcal P_{\mathrm{HCOP}}^{\mathrm{tool}}\), allowing repeated route, frontier, and no-false-certificate safety prose to be replaced by typed tool-contract citations only where the original hypotheses and guards remain visible, while preserving proof content and the full historical witness in Part~II \\
\hline
v3.34.0-r09 & local working candidate & 2026-04-24 & no Zenodo DOI assigned --- keeps the r08 compression map unchanged and extracts the first reusable HC/OP tool normal forms \(\mathbb T_{\mathrm{HCOP}}^{(1)}\), giving typed contracts for lock/capture, handoff, normalization, routing, comparison, resolution, probe ranking, certificate tests, frontier descent, and archive-to-tool extraction, so later text can cite tool contracts instead of repeating safety prose while preserving proof content and the full historical witness in Part~II \\
\hline
v3.34.0-r08 & local working candidate & 2026-04-24 & no Zenodo DOI assigned --- keeps the r07 QAM/MML/MMA interface unchanged and adds \(\mathfrak C_{\mathrm{HCOP}}\), a conservative MML/MMA compression map for inherited HC/OP text, so repeated capture, handoff, normalization, routing, comparison, probe, certificate-test, route-preservation, and historical-tool paragraphs can be reread as typed operations without deleting proof content or manufacturing OP closure certificates, while preserving the full historical witness in Part~II \\
\hline
v3.34.0-r07 & local working candidate & 2026-04-24 & no Zenodo DOI assigned --- keeps the r04 paired OP-entry comparison execution and r06 bibliography hardening unchanged, assigns local interface versions \(\mathrm{QAM}_{\mathrm{HCOP}}^{\mathrm{loc}}\), \(\mathrm{MML}_{\mathrm{HCOP}}^{\mathrm{loc}}\), and \(\mathrm{MMA}_{\mathrm{HCOP}}^{\mathrm{loc}}\), and proves that the r01--r04 HC-routed OP-entry chain is executable as one finite typed MML/MMA program without manufacturing an OP closure certificate, while preserving the full historical witness in Part~II \\
v3.34.0-r06 & local working candidate & 2026-04-24 & no Zenodo DOI assigned --- keeps the r04/r05 mathematical and bibliography-placement state unchanged, expands the terminal bibliography after a whole-document literature scan, adds compact front/back literature anchor notes for Teichmuller/quasiconformal geometry, Riemann--GUE/RMT, Hopf/fibre-bundle topology, division algebras, gauge/QFT landmarks, CKM/PMNS mixing, Julia dynamics, set-theoretic foundations, and probe/cryptographic/computational diagnostics, while preserving the full historical witness in Part~II \\
v3.34.0-r05 & local working candidate & 2026-04-24 & no Zenodo DOI assigned --- keeps the r04 paired OP~1/4--OP~5(B) comparison execution layer unchanged and normalizes the bibliography/back-matter placement by moving the live bibliography to the terminal Final Bibliography block while preserving the full historical witness in Part~II \\
\hline
v3.34.0-r04 & local working candidate & 2026-04-24 & no Zenodo DOI assigned --- executes the paired OP~1/4--OP~5(B) comparison branch above the r03 task-routing layer; compares joint versus separated HC-compatible readouts and returns either a bound common-origin state, a split ranked single-branch state, or a smaller frontier-certified comparison state while preserving route tags and the full historical witness in Part~II \\
\hline
v3.34.0-r03 & local working candidate & 2026-04-24 & no Zenodo DOI assigned --- extracts the first finite OP-entry task-routing layer from the HC-derived normal form; maps the selected live residual branch into OP~1/4 continuum-limit work, OP~5(B) routed $\Theta$--$R$ balance work, anchor/admissibility refinement, or paired OP~1/4--OP~5(B) comparison while preserving route tags and the full historical witness in Part~II \\
\hline
v3.34.0-r02 & local working candidate & 2026-04-24 & no Zenodo DOI assigned --- extracts the first downstream-readable OP-entry normal form from the transport-stable HC-to-OP handoff packet; routes the terminal live residual triple into OP~1/4, OP~5(B), and anchor/admissibility request families while preserving route tags and the full historical witness in Part~II \\
\hline
v3.34.0-r01 & local working candidate & 2026-04-24 & no Zenodo DOI assigned --- begins the transport-stable HC-capture and OP-entry handoff line above the public v3.33.0 release; packages local HC capture into a route-faithful transport-stability test and an OP-facing handoff packet while preserving the full historical witness in Part~II \\
\hline
v3.25.0 & public Zenodo release & 2026-04-16 & \href{https://doi.org/10.5281/zenodo.19615952}{10.5281/zenodo.19615952} --- earlier public release; consolidates the optical-comparison and machine-theoretic overlay line, including compressed optical diagnostic control, bundled proof-use compression, conditional reference-core closure, and the capped machine-theoretic overlay line up to the present equivalence surface \\ 
\hline
v3.18.0 & public Zenodo release & 2026-04-13 & \href{https://doi.org/10.5281/zenodo.19548590}{10.5281/zenodo.19548590} --- earlier public release; carries the public v3.17.0 tensor/operator baseline forward into a first public OP~3/QAM closure package with shell certificates, shell-defect localization, shell-stable readout factorization, and a bundled OP~3$\to$OP~5/global interface while keeping stronger kernel-first activation internal \\
\hline
v3.17.0 & public Zenodo release & 2026-04-10 & \href{https://doi.org/10.5281/zenodo.19522352}{10.5281/zenodo.19522352} --- earlier public release; promotes the third public-shape QAM tensor/operator line above the public v3.16.0 release and keeps translation/round-trip statements in front-facing public form while leaving support-selection and wider compression internal \\
\hline
v3.16.0 & public Zenodo release & 2026-04-10 & \href{https://doi.org/10.5281/zenodo.19489324}{10.5281/zenodo.19489324} --- earlier public release; promotes the second public-shape QAM branch/shell/routed layer while leaving tensor/operator and support-selection expansion internal, and records a deferred kernel-first compression path with activation criteria, staged stress tests, and migration roadmap \\
\hline
v3.15.0 & public Zenodo release & 2026-04-08 & \href{https://doi.org/10.5281/zenodo.19463389}{10.5281/zenodo.19463389} --- earlier public release; freezes the first compressed public-shape QAM core on top of the public v3.14.0 branch/shell baseline \\
\hline
v3.14.0 & public Zenodo release & 2026-04-07 & \href{https://doi.org/10.5281/zenodo.19460589}{10.5281/zenodo.19460589} --- earlier public release; extends the public v3.13.0 line by freezing the first post-closure branch/shell hardening layer \\
\hline
v3.3.0 & local working state & 2026-04-03 & condensed Appendix-C integration of Julia-boundary, master-residual, Umkehrwalze, and reversal-probe sector \\
\hline
v3.3.1 & local working state & 2026-04-03 & clean condensed Appendix-C line; 16-page block reduced to a 3-page support appendix \\
\hline
v3.3.2 & local working state & 2026-04-03 & clean-mainline promotion: compact 3.3 synthesis theorem moved into main text; Appendix-C support retained \\
\hline
v3.3.3 & local final-polish state & 2026-04-03 & removed embedded editorial fragments in the new 3.3 synthesis line and normalized notation in the compressed summary \\
\hline
v3.3.4 & local editorial-clean state & 2026-04-03 & global editorial cleanup: removed remaining visible drafting fragments, normalized compressed remarks, and cleaned heuristic table notes \\
\hline
v3.3.5 & Zenodo release & 2026-04-03 & \href{https://doi.org/10.5281/zenodo.19396143}{10.5281/zenodo.19396143} --- strict wording cleanup: aligned version metadata, removed residual status-weakening phrasing in the 3.3 line, normalized the current-version record, and finalized the 3.3.5 release package \\
\hline
v3.4.0 & local working extension & 2026-04-03 & first 3.4 pass: added closure-potential globalization of the residual branch law, conserved branch energy, and an energy-form reading of the reversal-probe dichotomy \\
\hline
v3.4.1 & local working extension & 2026-04-03 & constrained the exact branch potential to the reversible single-well class and reduced the global scalar ambiguity to the transported $\mathbb Z_2$ pair \\
\hline
v3.4.2 & local working extension & 2026-04-03 & extracted the canonical quartic one-field family for the visible residual branch law and added turning-value parameterization of regular probe shells \\
\hline
v3.4.3 & local working extension & 2026-04-03 & tied the visible quartic branch coefficient to the transported quartic closure coefficient and replaced the technical turning-amplitude label by turning value \\
\hline
v3.4.4 & local working extension & 2026-04-03 & added the first nonlinear period-drift readout, showing that the timing defect on regular periodic shells is already fixed by the same quartic closure coefficient \\
\hline
v3.4.5 & local working extension & 2026-04-03 & extended the quartic readout to the full even hierarchy, proving that the first higher timing defect appears exactly at the order of the first surviving even closure coefficient \\
\hline
v3.4.6 & local working extension & 2026-04-03 & repackaged the transported even hierarchy into a single analytic generator $\Phi(\eta^2)$ for the visible branch potential and its nonlinear shell readouts \\
\hline
v3.4.7 & local working extension & 2026-04-03 & generator-shell probe classifier: small regular shells reduced to the generator radius $s(E)=\Phi^{-1}(E)$ and the transported sign pair $\pm\sqrt{s(E)}$ \\
\hline
v3.4.8 & local working extension & 2026-04-03 & shell-parity criterion: direct convergence versus spinorial/M\"obius lift determined by the even/odd reflector-crossing parity on sufficiently small regular shells; added first diagnostic Julia-type appendix render \\
\hline
v3.4.9 & local working extension & 2026-04-03 & added cautious frame-shifted detector-readout hypothesis for the diagnostic Julia/probe geometry, together with explicit candidate observables for later comparison work \\
\hline
v3.4.10 & local working extension & 2026-04-03 & converted the detector-readout hypothesis into a quantitative tunnel-observable interface: tunnel width, edge intensification, opening angle, and shell parity are now organized as a minimal diagnostic comparison package \\
\hline
v3.4.11 & local working extension & 2026-04-03 & added a calibrated inverse-readout theorem: tunnel width plus shell parity now recover the local shell radius, while edge intensification and opening angle become shell-consistency checks measured by $\Delta_{\mathrm{diag}}$ \\
\hline
v3.4.12 & local working extension & 2026-04-03 & added a local inverse identifiability/stability theorem: the calibrated tunnel-width law now carries a shell-wise stability margin and inverse condition number, so shell reconstruction is controlled whenever the calibration slope stays away from zero \\
\hline
v3.4.13 & local working extension & 2026-04-03 & organized the inverse readout globally into a shell atlas stratified by calibration-critical shells and parity-switch shells, so stable inversion now proceeds chart by chart away from the obstruction set \\
\hline
v3.4.14 & local working extension & 2026-04-03 & added an obstruction-shell transition law: isolated obstructions now carry only fold exchange, parity flip, or their composition, giving the inverse-readout layer a small $\mathbb Z_2$ transition algebra and local obstruction monodromy \\
\hline
v3.4.15 & local working extension & 2026-04-03 & organized the obstruction transitions into a global Klein-four transition cocycle on the shell atlas and showed that, after fold resolution, the remaining monodromy reduces to the parity local system underlying the spinorial/M\"obius lift \\
\hline
v3.4.16 & local working extension & 2026-04-03 & promoted the parity local system to a global double-cover criterion: direct visible convergence is globally available exactly when the parity class trivializes, while every nontrivial case becomes globally single-valued on the canonical parity double cover \\
\hline
v3.4.17 & local working extension & 2026-04-03 & recast the residual parity sector as a genuine cohomological obstruction class in $H^1(B,\mathbb Z_2)$, identifying direct visible convergence with the vanishing class and the canonical parity double cover with its associated classifying two-sheeted cover \\
\hline
v3.4.18 & local working extension & 2026-04-03 & promoted the parity obstruction to an operational global holonomy/readout character: admissible closed shell paths now evaluate the class $\varepsilon_{\mathrm{par}}$ directly, so direct visible return versus spinorial/M\"obius lift is read as a $\mathbb Z_2$ loop outcome factoring through mod-two homology \\
v3.4.19 & local working extension & 2026-04-03 & turned the parity holonomy character into an empirical loop-readout estimator: chartwise parity-jump counting across admissible obstruction shells now recovers the cohomological $\mathbb Z_2$ holonomy, giving a detector-side loop test for direct visible return versus spinorial/M\"obius lift \\
v3.4.20 & local closing state & 2026-04-03 & clean 3.4 closure: compressed the full extension line into a single generator-atlas-holonomy synthesis theorem, identified the minimal handoff package $(\Phi,\Delta_{\mathrm{diag}},\varepsilon_{\mathrm{par}},\chi_{\mathrm{par}},\widehat\chi_{\mathrm{par}})$, and marked the end of the 3.4 build-up before the next-stage simplification work \\
\hline
v3.5.0 & local working extension & 2026-04-03 & first compression-stage theorem: introduced the residual record $\mathfrak R(B)$ and proved the reconstruction normal form for the stabilized visible sector \\
\hline
v3.5.1 & local working extension & 2026-04-03 & record-level transport step: defined admissible residual-record morphisms and isomorphisms, and formalized the faithful translation principle on the stabilized visible sector \\
\hline
v3.5.2 & local abstract-polish state & 2026-04-03 & sharpened the abstract and the compression-stage framing: clarified the residual-record interface, the conservative-frontend requirement, and the split between structurally closed versus still-open global identification problems \\
\hline
v3.5.3 & local status-cleanup state & 2026-04-04 & reclassified the open-questions layer by actual status, absorbed OP~2(T4) into OP~1, marked OP~3 as no longer a primary open problem, and reduced OP~5 to Step~(B) only \\
\hline
v3.5.4 & local chronology-sync state & 2026-04-04 & synchronized the document head with the current 3.5.x line: added the missing v3.5.0--v3.5.4 chronology entries and replaced outdated roadmap goals by the residual-record/compression-stage agenda \\
\hline
v3.5.5 & local reader-guide alignment state & 2026-04-04 & aligned the Reader's Guide with the actual 3.5.x status: separated live fronts from preserved historical roadmap layers and made the OP~1/2/3/5 status split explicit at the document head \\
\hline
v3.5.6 & local main-text status-gloss state & 2026-04-04 & propagated the 3.5.x status split into the main compression-stage remarks and appendix support transitions: clarified that remaining fronts are downstream residual tasks and not invitations to reopen the visible interface \\
\hline
v3.5.7 & local stage-summary state & 2026-04-04 & added an explicit 3.5.x-versus-3.6 orientation block: recorded what the compression line has already achieved and fixed the next target as conservative language stabilization above the residual-record interface \\
\hline
v3.5.8 & local figure-integration state & 2026-04-04 & rebuilt the missing Julia/probe comparison render, integrated it into the text with a cautious diagnostic caption, and restored a clean PDF build with the image in place \\
\hline
v3.6.0 & local working extension & 2026-04-04 & introduced the conservative-frontend principle: any later language layer must factor through the residual-record interpreter and may not introduce new visible primitives \\
\hline
v3.6.1 & local frontend-hardening state & 2026-04-04 & hardened the conservative frontend line by adding closure under composition and a standalone primitive non-drift criterion; clarified that Section~3.6 fixes admissibility constraints for later language layers rather than introducing a second visible interface \\
\hline
v3.6.2 & local frontend-finalization state & 2026-04-04 & promoted the closure/non-drift hardening step into the main line and stabilized the 3.6 block as a completed admissibility framework for later language chains \\
\hline
v3.7.0 & local OP-first consolidation state & 2026-04-04 & formalized the current OP-first reading: the remaining live fronts are organized as residual downstream tasks over the fixed record interface, with direct language growth deliberately deferred \\
\hline
v3.7.1 & Zenodo release & 2026-04-04 & \href{https://doi.org/10.5281/zenodo.19412478}{10.5281/zenodo.19412478} --- public release of the OP-first 3.7 line; froze the present stabilized-reading claim and fixed 3.8 as an explicit OP-closure phase while keeping QAM deferred \\
\hline
v3.8.0 & local two-core-program state & 2026-04-04 & promoted 3.8 from a planning note to a working OP-closure stage: reduced the foundational historical burden to the OP~1 correction core and the OP~5 step~(B) Theta-balance core, with all other live fronts treated as secondary unless dependency is shown \\
\hline
v3.8.1 & local residual-emergence state & 2026-04-04 & recast the field-equation / controlled-emergence question into a residual-controlled physical interface form, concentrating the unresolved burden into the explicit remainder families \(\mathcal E_1\) and \(\mathcal E_5\) instead of introducing a new visible primitive basis \\
\hline
v3.8.2 & local OP~1 reduction state & 2026-04-04 & reduced the remaining OP~1 burden to a single scalar residual channel \(\mathcal R_1(\mu)\), making OP~1 a one-channel closure problem rather than a diffuse correction family \\
\hline
v3.8.3 & local OP~1 root-isolation state & 2026-04-04 & proved that the OP~1 scalar residual channel has a unique root in the physical window \(1.85\le \mu\le 1.90\), numerically \(\mu_\ast\approx 1.86366225\), thereby turning the remaining OP~1 burden from a broad numerical search into an exact structural-identification problem \\
v3.8.4 & local OP~1 frame-shift diagnostic state & 2026-04-04 & showed that the residual \(\gamma_L-\mu_\ast\) shrinks on the matter-frame offset scale, behaves like a local Newton/readout correction, and does not trigger a parity/horizon jump in the current Julia/probe family at the tested residual scales \\
v3.8.5 & local minimal neutrino-inflow test state & 2026-04-04 & converted the archived \(k(\nu)=4\) inflow hypothesis into a minimal OP~1 downstream correction ansatz \(\mathcal R_{1,\nu}(\mu;\eta_\nu)=\mathcal R_1(\mu)+\eta_\nu/\mu^4\), with calibrated \(\eta_\nu^\ast\approx -1.10\times10^{-3}\) and no reopening of the visible primitive basis \\
v3.8.6 & local semi-structural inflow estimate state & 2026-04-04 & showed that the calibrated coefficient satisfies \(\eta_\nu^\ast\approx -\varepsilon_\nu\gamma_L^4\beta_{\mathrm{sym}}^4\) with \(\varepsilon_\nu\approx 1.0012\times10^{-3}\), so the fitted inflow law reduces to a \(k=4\) matter-share factor times a per-mille fermionic leakage efficiency \\
v3.8.7 & local defect-squared inflow state & 2026-04-04 & identified the per-mille leakage efficiency with the already existing quadratic defect invariant \(\eta_{\mathrm{sym}}^2=(\frac13-\beta_{\mathrm{sym}}^2)^2\), yielding the sharper law \(\eta_\nu^\ast\approx -\eta_{\mathrm{sym}}^2\gamma_L^4\beta_{\mathrm{sym}}^4\) \\
v3.8.8 & local second-order fermionic leakage state & 2026-04-04 & used the \(J\)-antisymmetric fermionic sector \(\Hplusminus\) to justify why the first visible inflow scalar should be quadratic in the transported defect, leading to \(\eta_\nu\approx -\lambda_\nu\eta_{\mathrm{sym}}^2\gamma_L^4\beta_{\mathrm{sym}}^4\) with \(\lambda_\nu\approx 0.9990\) \\
v3.8.9 & local cosmological-offset link state & 2026-04-04 & showed that any cosmological contribution to the OP~1 inflow remainder is naturally mediated by the permanent frame-offset term \(\beta_{\mathrm{sym}}^2\) already used in the dark-offset theorem, not by a direct Hubble-rate coupling, yielding \(\eta_\nu^\ast\approx -(\frac13-\beta_{\mathrm{sym}}^2)^2\gamma_L^4\beta_{\mathrm{sym}}^4\) as the natural weak expansion reading \\
v3.8.10 & local one-background collapse state & 2026-04-04 & showed that the frame/readout, fermionic-leakage, and weak cosmological-offset readings of the smooth OP~1 remainder collapse numerically to the same background law \(-\eta_{\mathrm{sym}}^2\gamma_L^4\beta_{\mathrm{sym}}^4\), leaving the discrete prime-lattice mechanism as the only genuinely independent live burden \\
v3.8.11 & local readout-point cancellation state & 2026-04-04 & showed that the scalar prime/continuum residual \(\mathcal R_1(\gamma_L)\approx 9.1216\times10^{-5}\) and the collapsed smooth background term \(-\eta_{\mathrm{sym}}^2\beta_{\mathrm{sym}}^4\approx -9.1208\times10^{-5}\) cancel to within \(\Delta_{\mathrm{can}}\approx 8.68\times10^{-9}\), shifting the live OP~1 burden from finding a new correction term to explaining the near-identity itself \\
v3.8.12 & local background-corrected bridge state & 2026-04-04 & defined the bridge function \(\Phi_1(\mu)=\mathcal R_1(\mu)-\eta_{\mathrm{sym}}(\mu)^2\beta(\mu)^4\), proved that it has a unique root \(\gamma_{\mathrm{br}}\approx 1.86399996799643\) on the physical window, and showed that the Companion matter-frame value \(1.864\) already coincides with this bridge root to about \(3.2\times10^{-8}\) \\
v3.8.13 & local bridge-shape forcing state & 2026-04-04 & showed that the current Companion architecture forces any minimal admissible visible OP~1 bridge term to have the shape \(\lambda_{\mathrm{br}}\eta_{\mathrm{sym}}(\mu)^2\beta(\mu)^4\), leaving only the canonical normalization \(\lambda_{\mathrm{br}}\approx 1.00010\) as the live remaining freedom \\
v3.8.14 & local multiplicative bridge-renormalization state & 2026-04-04 & rewrote the forced bridge equation as the dressed readout law \(\beta(\mu)=\beta_{\mathrm{CL}}(\mu)\exp(\eta_{\mathrm{sym}}(\mu)^2\beta(\mu)^4)\), showing that the apparent non-canonicity of \(\lambda_{\mathrm{br}}\) at \(1.864\) is only the tiny \(8.68\times10^{-9}\) displacement from the exact bridge root rather than a missing normalization mechanism \\
v3.8.15 & local transport-balanced exponential state & 2026-04-04 & derived the exponential self-background dressing as the finite form of a repeated compensator-balanced transport law on the positive readout amplitude, reducing the remaining OP~1 burden to the path-constancy/normalization of the bridge source rather than to the appearance of the exponential itself \\
v3.8.16 & local bridge-contraction state & 2026-04-04 & recast the bridge law as the fixed-point equation \(T_\mu(\beta)=\beta\) and showed that the bridge map is strongly contractive at the distinguished root, with \(|\partial_b T_{\gamma_{\mathrm{br}}}(\beta_{\mathrm{br}})|\approx 0.00311\), thereby turning the source-freezing step into a local stability statement rather than a free assumption \\
v3.8.17 & local OP~1 closure state & 2026-04-04 & packaged the 3.8.2--3.8.16 line into a local closure theorem at the distinguished matter-frame point \(\gamma_{\mathrm{loc}}=\gamma_{\mathrm{br}}\approx 1.86399996799643\), relocating the remaining burden from local OP~1 readout to the question of global extension versus transition to OP~5 Step~(B) \\
v3.8.18 & final local OP~1 closure release state & 2026-04-04 & finalized the 3.8 release with version DOI \href{https://doi.org/10.5281/zenodo.19414221}{10.5281/zenodo.19414221}; the release message is local closure of OP~1 at the distinguished matter-frame point, with OP~5 Step~(B) promoted to the next primary live front \\
v3.8.19 & local OP~5 one-point reduction state & 2026-04-04 & used local OP~1 closure to fix the matter-frame and photon-frame background data, reducing historical OP~5 Step~(B) to the single scalar finite-cutoff identity \(\Omega_{5,\mathrm{loc}}=0\) at \(T_{\mathrm{loc}}=\gamma_{\mathrm{loc}}^{-3/2}\approx0.3929444410\) with target value \(\mathcal T(T_{\mathrm{loc}})\approx0.893255\) \\
v3.8.20 & local OP~5 exact-tail state & 2026-04-04 & rewrote historical OP~5 Step~(B) as the exact tail identity \(\mathcal E_N(T_{\mathrm{loc}})=0.893255-\mathcal T_N(T_{\mathrm{loc}})\) above an arbitrary truncation level \(N\); in particular, using the archived \(N=50\) diagnostic \(\mathcal T_{50}(T_{\mathrm{loc}})\approx-0.146\), the required tail contribution becomes \(\mathcal E_{50}(T_{\mathrm{loc}})\approx1.039255\) \\
v3.8.21 & local OP~5 counting-tail state & 2026-04-04 & rewrote the exact OP~5 tail as a Stieltjes integral against the photon-frame zero counting tail \(N_{\mathrm{ph}}^{>N}(u)\) and, after integration by parts, as the oscillatory counting-function integral with kernel \(K_{T_{\mathrm{loc}}}(u)=T_{\mathrm{loc}}\cos(T_{\mathrm{loc}}u)/u-\sin(T_{\mathrm{loc}}u)/u^2\), thereby isolating the exact analytic interaction that remains to be controlled \\
v3.8.22 & local OP~5 smooth-density split state & 2026-04-04 & decomposed the exact OP~5 tail into a smooth counting-density contribution \(\mathcal E_N^{\mathrm{sm}}(T_{\mathrm{loc}})\) plus a genuine fluctuation term \(\mathcal E_N^{\mathrm{fl}}(T_{\mathrm{loc}})\), turning Step~(B) into the precise question whether the distinguished-cutoff tail is density-dominated or fluctuation-dominated \\
v3.8.23 & local OP~5 leading-density test state & 2026-04-04 & inserted the leading transported photon-frame Riemann--von Mangoldt density, finding \(\mathcal E_{50,\mathrm{lead}}^{\mathrm{sm}}(T_{\mathrm{loc}})\approx 0.047867069\) versus the exact required tail \(\mathcal E_{50}(T_{\mathrm{loc}})\approx 1.039202153\); thus only about \(4.6\%\) of the tail is smooth and Step~(B) is decisively fluctuation-dominated \\
v3.8.24 & local OP~5 pure-fluctuation state & 2026-04-04 & subtracted the leading smooth counting density completely and rewrote Step~(B) as the pure fluctuation-tail identity \(-2\int_0^\infty S_{\mathrm{ph}}^{>N}(u)K_{T_{\mathrm{loc}}}(u)\,du=(0.893255-\mathcal T_N(T_{\mathrm{loc}}))-\mathcal E_{N,\mathrm{lead}}^{\mathrm{sm}}(T_{\mathrm{loc}})\); in particular, at \(N=50\) the fluctuation tail must contribute about \(0.991335084\) \\
v3.8.25 & local OP~5 pair-block state & 2026-04-04 & reorganized the pure fluctuation tail into adjacent photon-frame pair blocks \(P_{N,j}(T_{\mathrm{loc}})\), yielding the exact pair formula controlled by the pair means \(x_{N,j}\) and half-gaps \(d_{N,j}\); this turns the next OP~5 attack into a mesoscopic resonance test rather than another smooth-density refinement \\
v3.8.26 & local OP~5 quadratic-gap state & 2026-04-04 & showed that the exact adjacent pair-block profile is even in the half-gap \(d\), so the linear gap term vanishes identically and the first nontrivial local-gap correction is quadratic; consequently the first mesoscopic OP~5 layer is mean-phase dominated rather than gap-noise dominated \\
v3.8.27 & local OP~5 pair-mean quasi-lattice state & 2026-04-04 & showed that the adjacent pair means are well described by a slowly drifting density-controlled quasi-lattice, with observed step ratios near the prediction \(2/\rho_{\mathrm{ph}}(x_{N,j})\) having mean 0.997634; thus the first mesoscopic OP~5 burden is slow phase drift rather than linear gap disorder \\
v3.8.28 & local OP~5 four-block symmetry state & 2026-04-04 & grouped adjacent pair blocks into four-blocks and showed that the resulting super-gap dependence is even, so the linear pair-mean drift term vanishes identically; thus OP~5 now exhibits a two-stage mesoscopic suppression in which both the linear local-gap term and the linear pair-mean drift term are removed \\
v3.8.29 & local OP~5 super-mean quasi-lattice state & 2026-04-04 & showed that the four-block super-means themselves form a second-stage slowly drifting density-controlled quasi-lattice, with observed step ratios near the prediction \(4/\rho_{\mathrm{ph}}(X_{N,k})\) having mean 0.996050; thus the current OP~5 burden is best read on a hierarchy of quasi-lattices rather than as a raw oscillatory tail \\
v3.8.30 & local OP~5 quarter-turn super-mean state & 2026-04-04 & showed that the super-mean phases at the benchmark cutoff cycle through the four quarter-turn sectors in strict order for the first \(16\) super-means, while drifting on average by about \(-0.0652531\) relative to the exact \(\pi/2\) increment; thus the current OP~5 super-mean layer exhibits weak quarter-turn locking with slow slip \\
v3.8.31 & local OP~5 quarter-cycle packet state & 2026-04-04 & showed that one full quarter-turn cycle of the super-mean layer, i.e. a \(16\)-zero resonance packet, already exhibits strong cancellation: the first two packets leave only about \(2.21\%\) and \(3.99\%\) of their \(\ell^1\)-mass, so the mesoscopic OP~5 burden is best read in terms of packet drift and inter-packet residue rather than isolated block errors \\
v3.8.32 & local OP~5 packet-defect state & 2026-04-04 & derived the first-order packet-defect law, showing that once quarter-cycle cancellation is taken as the zeroth-order picture, the surviving OP~5 residue is carried only by slow packet amplitude drift and quarter-turn phase-slip, not by the ideal packet itself \\
v3.8.33 & local OP~5 packet-envelope state & 2026-04-04 & compressed the first-order packet-defect variables into a single scalar packet-envelope \(\mathfrak D_{N,m}\), proving that the packet residue is controlled by \(|R_{N,m}(T_{\mathrm{loc}})|\le \sqrt2\,\mathfrak D_{N,m}+O(\delta\varepsilon+\varepsilon^2)\); thus OP~5 is now a packet-envelope problem at first order \\
v3.8.34 & local OP~5 observable-envelope state & 2026-04-04 & introduced the observable lower packet-envelope proxy \(\underline{\mathfrak D}_{N,m}=|R_{N,m}(T_{\mathrm{loc}})|/\sqrt2\), giving benchmark values \(\underline{\mathfrak D}_{50,1}\approx 0.00273964\), \(\underline{\mathfrak D}_{50,2}\approx 0.00434351\), and \(\underline{\mathfrak D}_{50,3}\approx 0.00921439\); this moves Step~(B) onto a directly observable higher-scale scalar defect sequence \\
\hline
v3.8.50 & local merge-reconciliation checkpoint & 2026-04-04 & reconciled the former v3.5.2 abstract-polish checkpoint with the current v3.8.49 boundary-layer line; declared the 3.5.2 interface-normalization results fully absorbed into the active 3.8 working monolith and recorded the merge explicitly at the document head and in the main text \\
\hline
v3.8.51 & local reproducible-figure state & 2026-04-04 & regenerated the missing Julia/probe comparison figure from a local Python renderer, embedded the resulting PDF asset as the primary document figure with PNG fallback, and recorded the rebuild path directly in the appendix discussion \\
\hline
v3.8.52 & local current-scale figure state & 2026-04-04 & replaced the earlier standalone comparison marker by the current residual-scale diagnostic family $\gamma\in\{0,\Delta_\beta,\Delta_\theta,\Delta\mu\}$ and updated the appendix figure/caption accordingly \\
\hline
v3.8.53 & local caption/spacing fix state & 2026-04-04 & repaired caption spacing for the current residual-scale probe figure and stabilized the appendix-page layout for the diagnostic render \\
\hline
v3.8.55 & local deep-boundary diagnostic extension & 2026-04-04 & extended the Appendix-C diagnostics with a focused $\Delta\mu$ deep-zoom comparison, localized difference ladders, and edge-skeleton drift overlays to test whether the strongest current residual remains boundary-local under further magnification \\
v3.8.56 & local continuation diagnostic extension & 2026-04-04 & added a six-step continuation test from $0$ to $\Delta\mu$, including a boundary-belt strip, cumulative motion envelope, and horizon-profile continuation to check whether the active shell deforms continuously across the whole live residual interval \\
 v3.8.57 & local sensitivity-localization extension & 2026-04-04 & added finite-difference sensitivity maps, regional/radial sensitivity profiles, and a horizon-angle heatmap to show that the live residual response stays concentrated on the same horizon-adjacent boundary shell rather than spreading into the core \\
\hline
v3.8.58 & local horizon-sector persistence extension & 2026-04-04 & added coarse angular-sector persistence maps and dominant-peak tracking on the M\"obius horizon, showing that the active response remains locked to a small conjugate sector family instead of diffusing around the full circle \\
v3.8.59 & local OP~5 packet-wave diagnostics extension & 2026-04-05 & added visual diagnostics for the quarter-turn super-mean cycle, packet-level cancellation hierarchy, signed packet-defect wave fits, and the rotating packet-mode phase portrait; these figures make the current OP~5 reading visible as packet resonance plus slow signed wave rather than as an unstructured tail \\
v3.8.60 & local OP~5 breathing-envelope visual extension & 2026-04-05 & added visual diagnostics for the transverse envelope trough, the slow secondary envelope wave, the breathing-helix interpretation of the packet mode, and the carrier/envelope ratio test; these figures make the current OP~5 reading visible as an axially dominated carrier with a slower breathing boundary-layer correction rather than as a rigid fixed-radius helix \\
v3.8.61 & local OP~5 frame-flicker visual integration & 2026-04-05 & added live Appendix-C frame-flicker fit, decomposition, and geometric-readout figures and moved the new visual block into the actually compiled PDF path so that the moving-frame hypothesis becomes directly auditable in the active monolith \\
v3.9.0 & Zenodo release & 2026-04-05 & \href{https://doi.org/10.5281/zenodo.19423689}{10.5281/zenodo.19423689} --- public release preserving the stabilized residual-record interface, the reproducible M\"obius-shell diagnostic suite, and the OP~5 packet / breathing-envelope / frame-flicker visual line in one monolithic text \\
\hline
v3.10.0 & Zenodo release & 2026-04-06 & \href{https://doi.org/10.5281/zenodo.19434163}{10.5281/zenodo.19434163} --- public release based on the v3.9.0 monolith, with version and DOI record updated for the present release and front-matter release language cleaned for Zenodo \\
\hline
v3.10.1 & local working continuation & 2026-04-06 & front-cleaned working continuation of the public v3.10.0 release; historical roadmap and release-note blocks moved behind the nomenclature into an archival appendix position \\
\hline
v3.10.2 & local working tool-extraction state & 2026-04-06 & first explicit working tool layer for v3.11 preparation: transport-first, readout-after-transport, checkpoint-absorption, and epsilon-trap tools recorded together with a synchronized diagnosis of the live OP~1 and OP~5 burdens \\
\hline
v3.10.3 & local working OP~1 epsilon-bridge state & 2026-04-06 & reframed the live OP~1 burden as epsilon-trapped bridge normalization: canonical closure is now read as entry into a contractive bridge tube around the distinguished fixed point rather than as a second branch search \\
\hline
v3.10.4 & local working OP~1 theorem-promotion state & 2026-04-06 & promoted the OP~1 epsilon-bridge tool to theorem level by proving a residual-to-fixed-point estimate and a contractive non-escape theorem for the canonical bridge normalization inside the local bridge tube \\
v3.10.5 & local working OP~1 rounded-entry state & 2026-04-06 & proved that the standard rounded matter-frame readout already enters the local contractive bridge tube via the tiny readout-point displacement $|\gamma_L-\gamma_{\mathrm{br}}|\approx 3.2\times 10^{-8}$, reducing the live OP~1 burden to the structural identification $b_{\mathrm{can}}=b_{\mathrm{MF}}$ rather than to a second local bridge search \\
v3.10.6 & local working OP~1 uniqueness promotion state & 2026-04-06 & promoted admissible-readout uniqueness to theorem level, removing the final structural ambiguity between canonical normalization and the standard matter-frame readout and reducing OP~1 to a single explicit local bridge residue \\
v3.10.7 & local working shared-matrix diagnostic state & 2026-04-06 & added shared HC/MM/J leakage diagnostics and a structure-first translation of linear, differential, and linear-differential cryptanalysis for attacking the common OP~1/OP~5 matrix remainder \\
v3.10.8 & local working shared-phase coupling diagnostic state & 2026-04-06 & added a shared-carrier / quadrature diagnostic for testing whether the remaining OP~1 and OP~5 residues are one in-phase mode or two phase-shifted readouts of the same hidden leakage carrier \\
v3.10.9 & local working shared-leakage probe-family state & 2026-04-06 & replaced the abstract phase-coupling parameter by a concrete HC/MM/J cross-sector probe family $\mathcal M_\lambda=\mathcal M+\lambda\Delta_{\mathrm{shared}}$ and isolated its first-order effect on the visible OP~1/OP~5 channels \\
v3.10.10 & local working probe-invariant classification state & 2026-04-06 & introduced the normalized determinant and dot-product invariants of the HC/MM/J probe response matrix, turning the shared-carrier versus quadrature question into a decision-grade first-order diagnostic \\
v3.10.11 & local working canonical low-rank leakage state & 2026-04-06 & fixed a canonical rank-two HC/MM/J leakage class and proved that its first-order response is necessarily one-mode rather than quadrature, so any later sinus/cosinus picture would have to be emergent rather than primitive \\
v3.10.12 & local working bidirectional leakage-plane state & 2026-04-06 & added the minimal paired rank-four cross-sector leakage plane and proved that it is the first structurally honest location where a quadrature/sinus-cosinus picture can arise, whereas primitive rank-two leakage remains one-mode only \\
  v3.10.13 & local working visible-ellipse state & 2026-04-06 & coupled the minimal paired leakage plane to the OP~1/OP~5 mask layer and proved that the visible response is generically an ellipse, collapsing to a line in rank one and to exact circle/quadrature only under conformal mask geometry on the leakage plane \\
  v3.10.14 & local working shared-mask decomposition state & 2026-04-06 & decomposed the visible OP~1/OP~5 mask layer into a shared channel plus separation defect, proved singular-value bounds in the near-shared regime, and showed that exact quadrature cannot be primitive when the two visible masks differ only by a small downstream correction \\
  v3.10.17 & local working observer-split control state & 2026-04-06 & showed that the visible OP~1/OP~5 separation defect is exactly the restriction of the observer-side mask difference to the active leakage plane, proving that near-shared one-mode behavior is the default whenever the two visible masks differ only by a small downstream readout split \\
  v3.10.18 & local working shared-residual dichotomy state & 2026-04-06 & decomposed the remaining primitive OP~1/OP~5 visible residual into a shared matrix-leakage carrier plus a downstream observer-shell split and proved that exact grouped no-drift forces any surviving primitive residual on the active leakage plane to be common-mode rather than primitive quadrature \\
  v3.10.19 & local working active matrix-fit state & 2026-04-06 & pushed the remaining primitive common-mode carrier back to the active HC/MM/J fit and proved that its size is controlled exactly by the active commutator $[\mathcal M,J]|_{U}$ on the admissible branch \\
  v3.10.20 & local working exact-transport dichotomy state & 2026-04-06 & used the grouped-layer admissible-transport architecture to show that exact grouped Millennium transport kills the active commutator on the admissible branch, yielding a clean dichotomy between exact primitive closure and a purely downstream observer-shell residual \\
  v3.10.22 & local working canonical observer-shell residual state & 2026-04-06 & promoted the remaining visible OP~1/OP~5 term, when primitive matrix leakage is absent, to a canonical downstream observer-shell residual carried by a common shell dressing mode rather than by a primitive HC/MM/J leakage mode \\
  v3.10.23 & local working canonical shell-law modeling state & 2026-04-06 & modeled the surviving shared OP~1/OP~5 shell mode after exact transport as a first-order phase-shifted / expansion-dressed observer-shell law, reducing the visible target to exact vanishing versus canonical post-transport shell identification \\
  v3.10.24 & local working shell-dominance classification state & 2026-04-06 & compressed the remaining visible post-transport question to a comparison of first-order shell loads, classifying the shared shell mode as phase-dominant, expansion-dominant, or genuinely mixed at first order \\
  v3.10.25 & local working shell-cancellation corridor state & 2026-04-06 & reduced the shared shell mode to a signed first-order shell-balance mismatch, isolating the cancellation corridor for exact visible closure and showing that same-sign shell loading forces a canonical nonzero shell remainder unless the linear shell load itself vanishes \\
  v3.10.26 & local working normalized shell-defect state & 2026-04-06 & compressed the shared observer-shell problem to a single normalized shell-cancellation defect, extracted an exact cancellation-depth invariant, and proved that exact visible closure requires the normalized first-order shell mismatch to be no larger than the normalized higher-order shell remainder \\
  v3.10.28 & local working sign-corridor shell state & 2026-04-06 & converted the normalized shell verdict into a sharp sign/corridor dichotomy, proved that sign-coherent first-order shell loads force canonical nonzero visible residual whenever the normalized higher-order rescue budget is below one, and extracted the exact anti-phase corridor required for possible visible closure \\
  v3.10.29 & local working shell-conjugacy gain state & 2026-04-06 & pushed the induced phase and expansion gains back to a shell-conjugate observer geometry, introduced the downstream shell involution and transport-balance defect, and showed that the anti-phase corridor is structurally natural whenever phase and expansion are J-descended conjugate readouts of one shell displacement with nearly balanced transport weights \\
\hline
v3.10.30 & local working shell anti-phase budget state & 2026-04-06 & compressed the anti-phase corridor to a single normalized shell-conjugacy load and showed that a small total shell load forces the induced shell ratio into a natural anti-phase tube around $-1$ \\
\hline
v3.10.31 & local working architecture-driven shell state & 2026-04-06 & pushed the shell-balance driver back to a combined architecture budget given by active HC/MM/J fit plus observer-shell dressing, so that first-order visible cancellation is now read as architecture-driven rather than phenomenological \\
\hline
v3.10.32 & local working architecture-rescue corridor state & 2026-04-06 & reduced the remaining visible question to a single budget comparison between the normalized combined architecture load and the normalized higher-order rescue budget, giving a final first-order corridor for exact visible closure after exact transport \\
\hline
v3.11.0 & Zenodo release & 2026-04-06 & \href{https://doi.org/10.5281/zenodo.19444128}{10.5281/zenodo.19444128} --- public release consolidating the v3.10.32 architecture-rescue corridor state, explicit tool-layer extraction, and release-normalized front matter for Zenodo \\
v3.11.2 & local working state & 2026-04-07 & OP~1/OP~5 closed as a native source-return pair via closure-class preservation; branchwise refinement, admissible routed aggregation, and the final visible-support mismatch decomposition added as the compact main-line hardening beyond the public v3.11.0 release  \\
 v3.12.0 & Zenodo release & 2026-04-07 & \href{https://doi.org/10.5281/zenodo.19446264}{10.5281/zenodo.19446264} --- public release promoting the OP~1/OP~5 closure line into the public-facing main text, with branchwise refinement and routed-readout consequences integrated into the monolithic Companion \\
\hline
v0.1.0 & Zenodo release & 2026-03-24 & \href{https://doi.org/10.5281/zenodo.19210729}{10.5281/zenodo.19210729} \\
\hline
v0.3.0 & Zenodo release & 2026-03-24 & \href{https://doi.org/10.5281/zenodo.19211785}{10.5281/zenodo.19211785} \\
\hline
v0.4.0 & Zenodo release & 2026-03-24 & \href{https://doi.org/10.5281/zenodo.19211798}{10.5281/zenodo.19211798} \\
\hline
v0.5.0 & Zenodo release & 2026-03-24 & \href{https://doi.org/10.5281/zenodo.19211877}{10.5281/zenodo.19211877} \\
\hline
v0.6.0 & Zenodo release & 2026-03-24 & \href{https://doi.org/10.5281/zenodo.19212073}{10.5281/zenodo.19212073} \\
\hline
v0.7.0 & Zenodo release & 2026-03-24 & \href{https://doi.org/10.5281/zenodo.19212212}{10.5281/zenodo.19212212} \\
\hline
v0.8.0 & Zenodo release & 2026-03-24 & \href{https://doi.org/10.5281/zenodo.19212324}{10.5281/zenodo.19212324} \\
\hline
v0.9.0 & Zenodo release & 2026-03-24 & \href{https://doi.org/10.5281/zenodo.19220003}{10.5281/zenodo.19220003} \\
\hline
v1.0.0 & Zenodo release & 2026-03-24 & \href{https://doi.org/10.5281/zenodo.19220118}{10.5281/zenodo.19220118} \\
\hline
v1.6.0 & Zenodo release & 2026-03-25 & \href{https://doi.org/10.5281/zenodo.19221335}{10.5281/zenodo.19221335} \\
\hline
v1.7.0 & Zenodo release & 2026-03-25 & \href{https://doi.org/10.5281/zenodo.19221372}{10.5281/zenodo.19221372} \\
\hline
v1.8.0 & Zenodo release & 2026-03-25 & \href{https://doi.org/10.5281/zenodo.19221431}{10.5281/zenodo.19221431} \\
\hline
v1.9.0 & Zenodo release & 2026-03-25 & \href{https://doi.org/10.5281/zenodo.19221473}{10.5281/zenodo.19221473} \\
\hline
v2.0.0 & Zenodo release & 2026-03-25 & \href{https://doi.org/10.5281/zenodo.19221527}{10.5281/zenodo.19221527} \\
\hline
v2.1.0 & Zenodo release & 2026-03-25 & \href{https://doi.org/10.5281/zenodo.19221618}{10.5281/zenodo.19221618} \\
\hline
v2.3.0 & Zenodo release & 2026-03-25 & \href{https://doi.org/10.5281/zenodo.19221864}{10.5281/zenodo.19221864} \\
\hline
v2.4.0 & Zenodo release & 2026-03-25 & \href{https://doi.org/10.5281/zenodo.19221954}{10.5281/zenodo.19221954} \\
\hline
v2.9.0 & Zenodo release & 2026-03-25 & \href{https://doi.org/10.5281/zenodo.19222638}{10.5281/zenodo.19222638} \\
\hline
v2.10.0 & Zenodo release & 2026-03-25 & \href{https://doi.org/10.5281/zenodo.19222894}{10.5281/zenodo.19222894} \\
\hline
v2.11.0 & Zenodo release & 2026-03-25 & \href{https://doi.org/10.5281/zenodo.19223084}{10.5281/zenodo.19223084} \\
\hline
v2.11.1 & local-only repair state & 2026-03-25 & non-Zenodo version-fix state between v2.11.0 and v2.11.2 \\
\hline
v2.11.2 & Zenodo release & 2026-03-25 & \href{https://doi.org/10.5281/zenodo.19223365}{10.5281/zenodo.19223365} \\
\hline
v2.11.3 & Zenodo release & 2026-03-25 & \href{https://doi.org/10.5281/zenodo.19223515}{10.5281/zenodo.19223515} \\
\hline
v2.12.0 & Zenodo release & 2026-03-25 & \href{https://doi.org/10.5281/zenodo.19223782}{10.5281/zenodo.19223782} \\
\hline
v2.13.0 & Zenodo release & 2026-03-25 & \href{https://doi.org/10.5281/zenodo.19223968}{10.5281/zenodo.19223968} \\
\hline
v2.14.0 & Zenodo release & 2026-03-25 & \href{https://doi.org/10.5281/zenodo.19225191}{10.5281/zenodo.19225191} \\
v2.15.0 & Zenodo release & 2026-03-25 & \href{https://doi.org/10.5281/zenodo.19226408}{10.5281/zenodo.19226408} \\
\hline
v2.16.0 & Zenodo release & 2026-03-25 & \href{https://doi.org/10.5281/zenodo.19226445}{10.5281/zenodo.19226445} \\
\hline
v2.17.0 & Zenodo release & 2026-03-25 & \href{https://doi.org/10.5281/zenodo.19226497}{10.5281/zenodo.19226497} \\
v2.18.0 & Zenodo release & 2026-03-25 & \href{https://doi.org/10.5281/zenodo.19226549}{10.5281/zenodo.19226549} \\
\hline
v2.19.0 & Zenodo release & 2026-03-25 & \href{https://doi.org/10.5281/zenodo.19226773}{10.5281/zenodo.19226773} \\
\hline
v2.20.0 & Zenodo release & 2026-03-26 & \href{https://doi.org/10.5281/zenodo.19227889}{10.5281/zenodo.19227889} \\
\hline
v2.21.0 & Zenodo release & 2026-03-26 & \href{https://doi.org/10.5281/zenodo.19227940}{10.5281/zenodo.19227940} \\
\hline
v2.22.0 & Zenodo release & 2026-03-26 & \href{https://doi.org/10.5281/zenodo.19228061}{10.5281/zenodo.19228061} \\
\hline
v2.23.0 & Zenodo release & 2026-03-26 & \href{https://doi.org/10.5281/zenodo.19239228}{10.5281/zenodo.19239228} \\
\hline
v2.23.1 & Zenodo patch release & 2026-03-26 & \href{https://doi.org/10.5281/zenodo.19239967}{10.5281/zenodo.19239967} \\
\hline
v2.24.0 & Zenodo release & 2026-03-27 & \href{https://doi.org/10.5281/zenodo.19258774}{10.5281/zenodo.19258774} \\
v2.25.0 & Zenodo release & 2026-03-27 & \href{https://doi.org/10.5281/zenodo.19259425}{10.5281/zenodo.19259425} \\
 v2.25.1 & Zenodo patch release & 2026-03-27 & \href{https://doi.org/10.5281/zenodo.19262890}{10.5281/zenodo.19262890} \\
\hline
v2.27.0 & Zenodo release & 2026-03-27 & \href{https://doi.org/10.5281/zenodo.19266666}{10.5281/zenodo.19266666} \\
\hline
v2.27.1 & local-only repair state & 2026-03-28 & baseline-preserving post-release working state \\
\hline
v2.27.2 & local-only repair state & 2026-03-28 & normalization and proof-chain hardening \\
\hline
v2.27.3 & local-only repair state & 2026-03-28 & structural closure pass; last state before v2.27.4 integration \\
\hline
v2.27.4 & local working state & 2026-03-28 & spectral-projector layer ($C=\Pi^{(\mathcal{M})}_{\mathrm{spec}}$), readout-relative lemma methodology, Theorem~7 (HC-induced readout geometry), Hypotheses~7.4--7.6, bidirectional probe lensing, SM spectral realization, regime-test appendix \\
\hline
v2.27.5 & Zenodo release & 2026-03-28 & \href{https://doi.org/10.5281/zenodo.19297422}{10.5281/zenodo.19297422} — full proof-chain audit; closure roadmap; 35 missing labels; five Priority-A gaps identified \\
\hline
v2.27.8 & local working state & 2026-03-28 & B1+B2 closed: \texttt{prop:semantic-forcing-chain} explicit 3-part proof from axioms (S1)--(S8); \newline \texttt{prop:structural-lift-principle} non-circular via \texttt{lem:group-transport} + \texttt{lem:hc-admissibility-lemma}; \newline OP~2 formalized; all 7 Roadmap items closed \\
\hline
v2.27.7 & local working state & 2026-03-28 & A4+A5 closed: \texttt{thm:koide-triolic} with $\theta_K$ formula; \newline 0.111\% match and $115\times$ amplification remark; \newline \texttt{thm:bsym-hopf} with proof and Hopf-gap remark; \newline 4.9\%/$6.8\%$, OP1/OP2 coupling; open: B1, B2 \\
\hline
v2.28.0 & Zenodo release & 2026-03-28 & \href{https://doi.org/10.5281/zenodo.19298586}{10.5281/zenodo.19298586}; consolidation of v2.27.4--v2.27.8; spectral-projector unconditional; bijection lemma; SM realization; readout-relative methodology; Thm~7; Hyp~7.4--7.6; OP~2 stub \\
\hline
v2.29.0 & Zenodo release & 2026-03-28 & \href{https://doi.org/10.5281/zenodo.19299245}{10.5281/zenodo.19299245} --- OP~5 step~(A) closed; step~(B) quantified; renormalization explicit \\
\hline
v2.29.1 & local working state & 2026-03-28 & GOE/GUE split (\texttt{prop:goe-fermionic}); $\gL$ = GOE/GUE ratio; $\gamma_L=1.197$ corrected to $1.864$; $\chi^2$ improvement $75.8\%$; $\delta_{\mathrm{CL}}$ exponent $\gamma_L^2\to\gamma_L^5$; \texttt{thm:dual-origin} corrected \\
\hline
v2.29.2 & Zenodo release & 2026-03-28 & \href{https://doi.org/10.5281/zenodo.19300351}{10.5281/zenodo.19300351} --- Einstein blunder reframed; $\Omega_M=\bsym^2$ determination; three-gap convergence \\
\hline
v2.29.3 & Zenodo release & 2026-03-28 & \href{https://doi.org/10.5281/zenodo.19300813}{10.5281/zenodo.19300813} --- OP~2 (T1)(T2) closed; $D_-=L_0|_{\Hplusminus}$ (corrected candidate); GOE essential self-adjointness; T4=OP1 equivalence; \textbf{Tribonacci theorem} $\gamma_L\to\mathbb{T}$; cascade normalization $1/\mathbb{T}+1/\mathbb{T}^2+1/\mathbb{T}^3=1$ \\
\hline
\end{longtable}
}

\noindent Note. The sequence above records the currently preserved release trail.
The day-level publication dates are normalized from the project publication log so that
Zenodo releases and local repair states remain chronologically comparable inside a single
appendix table.


\subsection*{v2.26.0 computational consolidation note}

Version~\textbf{2.26.0} executes the consolidation roadmap that had been reserved in
v2.25.1. The computational strand is retained but internally reorganized into a more
coherent progression: historical perspective, structural machine reading, Triolen
Algebra / Boolean readout, MML--MMA--MCP stratification, proof templates,
certificate-bearing execution, normalization, addressing, scheduling, host export,
and extraction outlook. The present version still remains exploratory in the sense of
future formalization, but the material is no longer presented only as a loose teaser;
it is now collected as a first internally coherent computational chapter of the
Companion.



\subsection*{Hypothesis-control pass for v2.27.0}

Before the next public milestone, the mathematical layer should distinguish more sharply between statements that are expected to become theorem-grade in v2.27.0 and statements that should remain explicitly hypothesis-grade even after the next hardening cycle. The practical rule is conservative: only statements with stabilized definitions, closed dependency chains, and reference-stable support should move upward from hypothesis status into theorem, proposition, or lemma status. Statements that continue to depend on interpretive transport, physical reading, or still-open structural identifications should remain explicitly marked as hypotheses or outlook items.

For the next pass, the intended control logic is therefore blockwise. In the ground-condition / HC / confinement chain, the target is theorem-readiness whenever the dependence on the collapse picture and confinement mechanism can be formulated without appealing to deferred computational scaffolding. In projector / closure mathematics, the target is theorem or lemma status whenever the relevant admissibility and sector-preservation claims are already internal to the mathematical layer. In the transport / Lorentz layer, statements that are still partly interpretive should remain hypothesis-grade unless they can be rewritten as strictly mathematical transport statements. For the Millennium-matrix and trace/branch layers, the default stance remains more cautious: structural mathematical claims may harden, but computational or machine-reading interpretations remain secondary and should not silently inherit theorem status from the mathematical core.

The resulting editorial discipline for v2.27.0 is simple: theorem promotion requires closure, dependency control, and stable notation; everything else remains visibly hypothesis, deferred claim, or teaser architecture. This avoids theorem-to-hypothesis drift in reverse as well: once a statement is kept at hypothesis level, later references should not read it as already mathematically closed merely because it is computationally suggestive.

\subsection*{v2.26.33 local mathematics-first note}

Version~\textbf{2.27.0} does not yet claim a new public release beyond v2.26.0. Instead, it records the explicit shift of near-term priority back to mathematical completion. The computational chapter remains preserved as a secondary but already visible reserve architecture. The immediate next target is v2.27.0 as a mathematics-first hardening cycle: close the mathematical layer, collect tools as mathematical appendices, and only then resume the proof-engine line as a compression language for already stabilized structure.


\subsection*{Theorem-promotion checklist for v2.27.0}
The next public milestone should not promote statements merely because they have become familiar in the internal workflow. Each promotion from hypothesis-language into lemma-, proposition-, or theorem-language must therefore be explicit, local, and reference-stable. For the mathematical closure pass leading to v2.27.0, the promotion discipline is:
\begin{itemize}[leftmargin=*]
  \item statements about the ground condition to HC to confinement chain may be promoted only when the chain has been written in one direction without hidden interpretive jumps;
  \item projector/closure claims may be promoted only when the admissibility conditions and the closure domain are fixed in stable notation;
  \item transport/Lorentz statements may be promoted only when they are phrased as mathematical transport claims rather than as physical metaphor;
  \item Millennium-matrix statements may be promoted only when the matrix is treated first as a mathematical object with domain, codomain, and structural role, not merely as a computational image;
  \item trace/branch claims may be promoted only when the branch criteria and trace-compatibility conditions are stated in theorem-capable form.
\end{itemize}
In parallel, several statement families are to remain explicitly outside theorem-promotion for v2.27.0 unless a full closure pass is actually completed. These include broad computational-universality claims, strong proof-kernel claims, speculative host-language claims, and any physically interpretive reading that still rests on heuristic transport language rather than closed mathematical formulation.

A practical promotion checklist for the next pass is therefore:
\begin{enumerate}[leftmargin=*]
  \item identify candidate statements that are currently written as heuristic summary sentences;
  \item decide whether each candidate should become a Definition, Lemma, Proposition, or remain a Hypothesis;
  \item verify that every promoted statement has stable notation, explicit assumptions, and stable backward references;
  \item verify that no promoted statement depends silently on teaser/computational language;
  \item record all non-promoted but still important claims explicitly as open, deferred, or interpretive.
\end{enumerate}
This checklist is intended to keep the next public release mathematically harder without forcing premature promotion of material that should still remain hypothesis-level.

\subsection*{Mathematical freeze criterion for v2.27.0}
For the next public release, a mathematical core block should count as \emph{frozen for release} only when it is not merely well phrased, but also stable under the local closure discipline introduced in the present subversion sequence. In particular, a block is to be treated as mathematically release-ready for v2.27.0 only if all of the following are satisfied:
\begin{enumerate}[leftmargin=*]
  \item the principal objects of the block are defined in stable notation and do not drift between theorem-language and hypothesis-language;
  \item the block depends only on explicitly cited prior definitions, lemmas, propositions, or tools, rather than on narrative intuition or roadmap placeholders;
  \item the key admissibility conditions are written inside the statement or proof scaffold itself;
  \item all required cross-references are stable and do not point to provisional teaser/computational material;
  \item the block can be read both locally and in the dependency order recorded for the mathematics-first pass without hidden backward jumps;
  \item any remaining interpretive or physical reading is clearly separated from the mathematical closure claim.
\end{enumerate}
A block should \emph{not} be counted as frozen merely because its wording is polished, because it already appears in summary form elsewhere, or because it has an attractive structural interpretation. For v2.27.0, the intended release criterion is therefore conservative: closure requires definitional stability, dependency stability, reference stability, and status stability at the same time.

A practical freeze test for each A-priority block is:
\begin{itemize}[leftmargin=*]
  \item can the block be pointed to by a stable label and theorem-status marker?
  \item can its assumptions be listed in one compact place?
  \item can its proof or justification be read without consulting deferred computational scaffolds?
  \item can its use in later blocks be stated without introducing new notation drift?
\end{itemize}
If the answer to any of these questions remains negative, the block should remain marked as active work rather than as mathematically frozen for public release. This criterion is meant to prevent premature hardening of material that is already suggestive but not yet sufficiently closed.


\subsection*{Release-readiness matrix for v2.27.0}
To prevent the next public release from being driven by impression rather than closure, the mathematics-first pass should end in a compact readiness matrix for the A-priority blocks. The matrix is not meant to replace proofs; it is meant to summarize, in one release-facing place, whether a core block is genuinely ready to be frozen, or whether a single missing condition still blocks public hardening.

\begin{center}
{\tiny\setlength{\tabcolsep}{2.5pt}\begin{tabular}{@{}p{0.18\textwidth} p{0.14\textwidth} p{0.10\textwidth} p{0.23\textwidth} p{0.16\textwidth}@{}}
\toprule
\textbf{Core block} & \textbf{Current state} & \textbf{Freeze-ready?} & \textbf{Missing closure condition} & \textbf{Primary blocker for v2.27.0}\\
\midrule
Ground condition / collapse picture & structurally central, still partly summary-driven & no & one stable definitional pass with fixed assumptions & notation and statement drift\\
HC emergence / confinement chain & strongly anchored, not yet written in one compact theorem-capable line & partial & explicit one-directional closure chain & hidden interpretive compression\\
Projector / closure mathematics & partially stabilized and reusable & partial & admissibility conditions and closure domain fixed in one notation regime & local definition spread\\
Transport / Lorentz layer & present but mixed with interpretive language & no & transport claims restated as mathematical transport statements & metaphor drift\\
Millennium matrix as mathematical object & increasingly structured, still mixed with computational imagery & no & fixed domain/codomain/operator role & object-status not yet fully hardened\\
Trace / branch mathematics & structurally rich, still partly template-style & partial & stable branch criterion and trace compatibility form & theorem-status not yet fixed\\
Tool / appendix lemmas & abundant support layer & partial & priority ordering and stable dependence tags & appendix support not yet fully normalized\\
\bottomrule
\end{tabular}}
\end{center}

For the next public milestone, the conservative reading of this matrix is intentional: \emph{partial} does not yet count as release-ready. A block should only be treated as ready for v2.27.0 if its missing closure condition has been discharged explicitly and its blocker is no longer active. In this sense, the matrix is a release-control device rather than a rhetorical summary: it translates the mathematics-first agenda into a blockwise freeze decision surface.



\subsection*{v2.27.0 gate checklist}
To prevent the next public release from being triggered by momentum alone, the mathematics-first pass should culminate in a compact global gate checklist. This checklist is stricter than the blockwise readiness matrix: even if several A-priority blocks look individually advanced, the public jump to v2.27.0 should only be made once the release conditions hold \emph{together} at the document level.

\begin{itemize}[leftmargin=*]
  \item the A-priority mathematical blocks admit stable definitions, stable dependence order, and stable theorem-status markers;
  \item no public-facing core claim in the targeted v2.27.0 range still depends on deferred computational scaffolds for its mathematical reading;
  \item the blockwise freeze criteria and the release-readiness matrix no longer contain an active blocker in the A-priority line;
  \item the proof-dependency checklist and the reference-stability pass agree on a single release-safe label structure;
  \item hypothesis-control and theorem-promotion discipline are aligned, so that no statement is promoted merely because its interpretation is attractive;
  \item the transport / projector / Millennium-matrix line can be cited without notation drift or mixed theorem-versus-hypothesis status;
  \item appendix tools used by the A-priority line are stable enough to serve as support, but not so overloaded that they silently replace missing core proofs;
  \item the computational / MML / MMA / MCP chapter remains visibly present as teaser architecture, but no release-critical mathematical closure claim depends on it;
  \item the version record, roadmap, local subversion discipline, and closure-control blocks all point to the same public target without internal contradiction.
\end{itemize}

A practical release rule for the next public jump is therefore conservative: \emph{v2.27.0 should only be cut once the A-line is simultaneously blockwise closure-ready and globally gate-clean}. If one of these conditions still fails, the correct response is another local subversion rather than a premature public release. In this sense, the gate checklist converts the mathematics-first roadmap into a final version-control criterion.

\subsection*{Final pre-v2.27.0 control summary}
As the mathematics-first pass approaches a public release decision, the Companion should end with a compact control summary that can be read in one glance. The purpose of this summary is not to replace the detailed work matrix, target-state matrix, or release-readiness table, but to compress them into a final release-facing snapshot. In particular, the summary should answer five questions simultaneously: (i) which A-priority blocks are currently the decisive release drivers, (ii) whether each of them is already freeze-ready, (iii) what the leading blocker still is, (iv) whether the global gate checklist is currently satisfied, and (v) whether the recommended action is \emph{continue mathematical closure} or \emph{prepare public release}.

\begin{center}
\small\begin{tabular}{p{0.18\textwidth} p{0.11\textwidth} p{0.18\textwidth} p{0.14\textwidth} p{0.27\textwidth}}
\toprule
\textbf{A-priority block} & \textbf{Freeze-ready?} & \textbf{Main blocker} & \textbf{Gate status impact} & \textbf{Release recommendation}\\
\midrule
Ground condition / collapse picture & not yet & dependence and formulation hardening & blocks full gate closure & continue mathematical closure\\
HC emergence / confinement chain & partial & explicit mathematical discharge & blocks full gate closure & continue mathematical closure\\
Projector / closure mathematics & partial & stable theorem-level closure & blocks full gate closure & continue mathematical closure\\
Transport / Lorentz layer & conditional & transport formalization discipline & weakens gate confidence & continue mathematical closure\\
Millennium matrix as mathematical object & not yet & object-level hardening before interpretive use & blocks full gate closure & continue mathematical closure\\
Trace / branch mathematics & partial & stable dependence on prior blocks & blocks full gate closure & continue mathematical closure\\
\bottomrule
\end{tabular}
\end{center}

At the present stage, the control summary therefore recommends \emph{no public jump yet beyond the current release line}: the mathematics-first agenda should remain active until the decisive A-priority blocks no longer fail the freeze criterion and the global gate checklist turns from ``not yet'' to ``satisfied.'' The intended role of this final control summary is practical rather than rhetorical: it gives the next release decision a single consolidated dashboard and makes it harder to confuse local progress with document-level readiness.


\subsection*{One-page executive closure view}
The mathematics-first steering logic introduced in the recent local subversions is now dense enough that it should also be compressible into a single executive page. The purpose of this closure view is therefore not to add another planning layer, but to place roadmap, dependency map, target-state matrix, freeze criterion, gate checklist, and pre-release control summary onto one final compact surface. In practical terms, the page should allow a reader to answer four questions immediately: (i) what the next public version is really trying to close mathematically, (ii) which A-priority blocks still drive the release boundary, (iii) whether the document is globally gate-clean, and (iv) whether the correct next action is another local subversion or a public release jump.

\begin{center}
\begin{tabular}{p{0.22\textwidth} p{0.22\textwidth} p{0.18\textwidth} p{0.24\textwidth}}
\toprule
\textbf{Control layer} & \textbf{Main question} & \textbf{Current status} & \textbf{Operational consequence}\\
\midrule
Roadmap / work matrix & what must be closed first? & mathematics-first line active & keep proof-engine / MML line secondary\\
Dependency map & what must come before what? & ordered core chain available & no closure claims out of dependency order\\
Definition / notation pass & are the objects and labels stable? & partially stabilized & unify notation before theorem-promotion\\
Target-state / closure checklist & what counts as done? & blockwise closure conditions visible & only explicit closure counts as completion\\
Freeze criterion / readiness matrix & is a block release-ready? & mixed, mostly partial & keep A-line blocks active until blocker-free\\
Gate checklist / control summary & is the document globally ready? & not yet gate-clean & continue local subversions, no public jump yet\\
\bottomrule
\end{tabular}
\end{center}

In this condensed reading, the present Companion state remains productive but intentionally conservative: the mathematical architecture is now sufficiently organized to support disciplined closure work, but not yet sufficiently gate-clean to justify confusing local progress with public readiness. The executive closure view is therefore meant to function as a one-page release dashboard. Its decision rule is simple: as long as the A-priority mathematical blocks remain only partially frozen, and as long as the global gate still records an active blocker, the correct action is \emph{another explicit local subversion rather than a public version jump}. Once these conditions disappear simultaneously, the dashboard should flip from \emph{continue closure} to \emph{prepare the next public release}.


\subsection*{Release-vs-local decision rule}
A final practical distinction should also be made explicit before the next public jump: not every coherent local gain justifies a new Zenodo-facing version. The present Companion workflow therefore distinguishes between \emph{local closure iterations} and \emph{public release transitions}. A local subversion is appropriate whenever the document gains planning discipline, dependency clarity, notation stability, or blockwise closure information without yet removing the main A-line blockers. A public release transition is appropriate only when these local gains have accumulated into a document-wide state in which the mathematical core is not merely better organized, but sufficiently frozen, referentially stable, and gate-clean to support a new public boundary.

In operational terms, the decision rule may be summarized as follows. Move only by local subversion if one or more of the following remains true: (i) an A-priority block is still only partially closed, (ii) the theorem/hypothesis boundary is still under active normalization, (iii) the reference-stability pass is not yet quiet, (iv) the release-readiness matrix still lists an active primary blocker, or (v) the global gate checklist is not yet satisfied. Consider a public version jump only when all A-priority blocks satisfy their freeze criterion simultaneously, the dependency chain is stable under rereading, the hypothesis-control and theorem-promotion passes no longer produce status drift, and the release-control summary flips from \emph{continue closure} to \emph{prepare release}. This distinction is intentionally conservative: it protects the mathematical layer from being mistaken for release-ready merely because the planning architecture around it has become more refined.

\subsection*{Mathematics-first release doctrine}
The present local stage adopts a simple doctrine for all remaining pre-v2.27.0 work. First, every local change is to be judged by whether it closes, stabilizes, or clarifies a mathematical dependency in the A-priority chain; additions that merely decorate the proof-engine teaser without tightening the mathematics are therefore secondary by rule. Second, no new public release is justified merely by chapter growth or architectural enthusiasm: the release threshold remains mathematical closure, reference stability, and theorem-safe status control. Third, the computational/MML/MMA/MCP layer is to remain visible and preserved, but it is not permitted to become a hidden prerequisite for mathematical claims that should stand on their own. Fourth, every local subversion should remain append-only and auditable, so that recovery from drift or accidental rollback stays possible without reconstructing the chain from memory.

In operational terms, the doctrine reads as follows:
\begin{enumerate}[label=(\roman*)]
  \item prioritize closure of ground condition $\rightarrow$ HC $\rightarrow$ confinement before expanding secondary architecture;
  \item promote statements only when notation, dependencies, and reference anchors are already stable;
  \item treat tools, appendices, and computational reserves as support layers unless and until they are explicitly absorbed into the mathematical core;
  \item prefer local subversions over public version jumps whenever the work mainly sharpens internal control rather than closing a release-grade mathematical block.
\end{enumerate}

This doctrine is intentionally conservative. Its role is to keep the Companion aligned with the structure-first programme: mathematics first, proof-memory preserved, computational architecture retained as a prepared secondary line, and public releases tied to genuine closure rather than mere accumulation.


\subsection*{Mathematics-first action queue}
The next local working order is to be treated as explicit and sequential unless a later control note states otherwise. The active queue is: (1) close the ground-condition / collapse picture at the definition level; (2) harden the HC-emergence / confinement chain so that its dependency status is stable; (3) close the projector / closure layer with stable admissibility language; (4) stabilize the transport / Lorentz layer only after the closure prerequisites are fixed; (5) treat the Millennium matrix strictly as a mathematical object before any renewed computational extension; (6) move trace / branch mathematics only after the preceding transport layer is stable; and (7) keep appendix tools and the computational teaser chapter in support position unless a direct mathematical blocker requires them.

The negative rule is equally important. The next local subversions are not to prioritize new MML/MMA/MCP growth, new proof-engine macros, or broader computational rhetoric unless such additions are directly needed to preserve an already existing mathematical dependency. In this sense, the queue functions as an anti-drift device: it converts the mathematics-first doctrine into a concrete order of execution and makes explicit what should not be pulled forward in parallel.


\paragraph{Mathematics-first session protocol.}
To prevent local growth from drifting back into attractive but premature architecture work, the next internal work sessions should follow a short fixed protocol. Each session should begin by choosing exactly one A-priority mathematical block from the current action queue. The session should then record (i) the active block, (ii) the minimal closure target for that block, (iii) the exact dependencies that must already be stable, and (iv) the specific theorem/lemma/definition/tool boundary that is being touched. Only after this boundary is explicit should any local rewriting begin.

The execution rule for such sessions is intentionally narrow. First perform a definition pass for the chosen block; then a dependency pass; then a closure pass; then a reference-stability check; and only at the end, if necessary, a supporting appendix-tool pass. If a prospective edit does not directly improve the active mathematical block or its immediate dependencies, it should be deferred. In particular, new computational language material, proof-engine rhetoric, or architecture expansion should not be introduced during a session unless the active mathematical block explicitly depends on it.

At the end of each local session, the result should be classified in one of four ways: \emph{closed for now}, \emph{partially hardened}, \emph{dependency blocked}, or \emph{still open}. This classification should then feed back into the release-vs-local decision rule, so that local subversions remain true mathematical work logs rather than merely growing source snapshots. In this way, the present protocol operationalizes the mathematics-first doctrine at the level of day-to-day document evolution.



\subsection*{Anti-parallelization rule}
The present mathematics-first stage also requires an explicit rule against parallel drift. In the next local subversions, not more than one A-priority mathematical block should be treated as the primary closure target of a given pass. Supporting edits in neighboring blocks are allowed only insofar as they remove an immediate dependency or stabilize a reference chain for that active block. What should be avoided is simultaneous advancement of several major fronts at once---for example projector/closure mathematics, transport/Lorentz normalization, and Millennium-matrix formalization in the same pass---because this makes it too easy to confuse distributed motion with genuine closure.

The same rule applies across the mathematics/computational boundary. While the computational teaser chapter remains preserved, it should not be opened in parallel with a mathematical A-block unless a concrete blocker in that block explicitly requires it. In practical terms, the anti-parallelization rule means: one primary closure front per pass; all secondary edits must justify themselves as dependency support, reference stabilization, or notation repair for that front. This rule is intentionally restrictive, because the present problem is not lack of ideas but risk of dispersion. Its purpose is to convert the mathematics-first doctrine into a narrow execution discipline and thereby make local progress more legible, auditable, and release-relevant.



\subsection{Re-entry rule after each local mathematical session}
After each local mathematics-first session, the document should be re-entered through a short global control loop rather than by immediately opening a new frontier. The purpose of this rule is to prevent hidden drift, duplicated partial closure, or silent reclassification of still-open material. In practical terms, the re-entry pass should answer only four questions: (i) did the session close the targeted A-block for now, (ii) which dependencies were clarified or newly exposed, (iii) which references and labels were changed and therefore need stabilization, and (iv) whether the next session should continue the same block or move to the next dependency in the queue.

The re-entry rule therefore imposes the following order: first update the local block status (
\emph{closed for now}, \emph{partially hardened}, \emph{dependency blocked}, or \emph{still open}); second reconcile the result with the target-state matrix and release-readiness matrix; third update the proof-dependency checklist if new requirements were uncovered; and only then decide whether the next local subversion should continue the same mathematical block or step forward in the mathematics-first action queue. This rule is intentionally conservative: local progress is only counted as durable progress when it is reinserted into the global control architecture without breaking reference stability, hypothesis control, or the release-vs-local decision rule.



\subsection*{Session-close audit rule}
Each local mathematics-first session should end with a short audit pass before the next subversion is declared. This audit is intended to prevent apparent progress from being recorded merely at the level of added text. A local pass counts as a meaningful increment only if the targeted block has been re-checked against its definition status, dependency status, reference stability, and release-readiness classification.

The closing audit for a local session should therefore answer four questions explicitly: (i) was a definition actually stabilized, rather than only rephrased; (ii) was at least one dependency clarified, reduced, or made more explicit; (iii) were the relevant references and labels left in a more stable state than before; and (iv) did the session produce a genuine closure gain for the active block, rather than only a planning gain. If these questions cannot be answered positively, the session may still be useful, but it should be recorded as preparatory rather than as a closure increment.

In this sense, the session-close audit forms the final control layer of the mathematics-first protocol. It ensures that the subversion chain remains readable as a chain of actual closure work, not merely as a sequence of appended managerial notes.


\subsection*{Carry-forward rule between local mathematical sessions}
Not every item touched in one local mathematics-first session is allowed to be carried forward as if it were already stable. The present carry-forward rule therefore distinguishes between (i) content that is eligible to be inherited by the next local subversion as a working baseline, and (ii) content that must be carried forward only as still-open material. This distinction is intended to prevent partial rewrites, provisional notation choices, or merely suggestive closure claims from silently hardening through repetition alone.

A block element may be carried forward as part of the active baseline only if all of the following are already true: its local definition is stable enough not to be rewritten immediately in the next pass; its dependency position is unchanged or clarified; its reference anchors survive rereading without ambiguity; and its status classification does not rely on a teaser, soft interpretation, or unresolved supporting note. By contrast, material should be carried forward only as \emph{open}, \emph{dependency blocked}, or \emph{provisional} whenever one of these conditions still fails. In practical terms, this means that unresolved notation, unstable lemma wording, open dependency edges, or uncertain theorem-promotion status must remain visibly marked rather than being treated as already absorbed into the mathematical baseline.

The purpose of this rule is conservative continuity. It allows the next local subversion to inherit genuine closure gains, while forcing incomplete work to remain visibly incomplete. In this way, the carry-forward rule complements the session-close audit: the audit decides whether a session produced a real gain, and the carry-forward rule decides which parts of that gain are allowed to function as stable starting points for the next mathematics-first pass.


\subsection*{Regression-check rule for local mathematical continuity}
Because the present Companion has already experienced version drift and accidental rollback in earlier phases, the mathematics-first regime also requires an explicit regression-check rule. The purpose of this rule is not to slow down productive local work, but to ensure that no apparent gain in planning, notation, or closure discipline is purchased at the cost of silently losing already integrated mathematical material. In particular, each local subversion should be checked not only for what has been added, but also for whether any previously stabilized block, appendix tool, roadmap control layer, or restored support material has shrunk, disappeared, or reverted in status.

Operationally, the regression check is conservative and document-internal. After each local session, the current source should be reread against the immediately preceding local subversion with at least four questions in mind: (i) has the active A-priority block gained real closure content without losing prior dependencies, (ii) have restored support structures such as the table of contents, version chronology, or appendix anchors remained intact, (iii) have any local control blocks been duplicated, displaced, or pushed outside the effective document body, and (iv) has any internal version or bibliographic marker fallen behind the actual subversion state. If one of these checks fails, the correct response is not to treat the current file as progressive merely because it is newer, but to repair the regression before carrying the pass forward.

In this sense, the regression-check rule complements the carry-forward rule and the session-close audit. The audit asks whether the current session produced a real local gain; the carry-forward rule asks what may legitimately be inherited as stable baseline material; and the regression-check rule asks whether that inheritance still contains everything that was already secured before. The combined effect is intentionally protective: it turns local subversions into auditable closure steps rather than a chain of potentially drifting snapshots. A local mathematics-first version should therefore be regarded as acceptable only if it is both \emph{additively improved} and \emph{non-regressive} relative to the immediately preceding baseline.

\subsection*{Baseline-lock rule for local continuation}
To prevent the mathematics-first subversion chain from drifting across multiple informal starting points, the present Companion adopts an explicit baseline-lock rule. For each local pass, exactly one immediately preceding local version is treated as the binding comparison basis. All claims of progress, non-regression, and closure gain are measured against this locked baseline rather than against a vague memory of the recent editing sequence. The purpose of the rule is simple: a local file becomes meaningful as a closure step only if it can be shown to improve upon one clearly identified predecessor without sacrificing already secured material.

Operationally, a baseline may be replaced only under conservative conditions. A newly produced local subversion may become the next locked baseline if, and only if, it satisfies the regression check, preserves the restored support structures, carries forward only legitimately stable material, and moves the active mathematics block measurably closer to its stated closure target. If one of these conditions fails, the correct response is not to redefine the newer file as the new norm, but to keep the earlier baseline locked until the defect has been repaired. In particular, mere growth in page count, roadmap language, or appended managerial notes is not sufficient to justify baseline replacement.

This rule complements the anti-parallelization, re-entry, session-close, carry-forward, and regression rules by providing a stable anchor for comparison. The anti-parallelization rule limits what may be opened at once; the re-entry rule controls how local gains are fed back into the larger roadmap; the session-close audit judges whether a session produced a real gain; the carry-forward rule determines what may be inherited as stable; the regression check guards against silent loss; and the baseline-lock rule finally determines when the current subversion has become stable enough to function as the new working ground for the next mathematics-first step.


\subsection*{Change-scope rule for local mathematics-first sessions}
The mathematics-first regime also requires an explicit rule about \emph{scope}. In any given local session, not everything that appears editable should actually be edited. The active scope is restricted to the currently selected A-priority mathematical block, its immediate dependency edges, the notation and references required to stabilize that block, and only those appendix tools that are strictly necessary to support the same closure pass. This means that many tempting modifications---however reasonable in isolation---should remain out of scope if they do not directly improve the currently active closure target.

The practical purpose of the change-scope rule is to protect mathematical closure from diffuse accretion. A local session may therefore adjust wording, dependencies, notation, labels, or supporting appendices only when these changes can be justified as direct support for the active block. What remains explicitly out of scope includes broad computational expansion, speculative theorem promotion outside the active dependency cone, large stylistic rewrites of unrelated sections, and opportunistic edits that merely make the document longer without bringing the selected block closer to a frozen mathematical state. In this way, the change-scope rule complements the anti-parallelization rule: anti-parallelization limits how many fronts may be open at once, while scope control limits how wide a single front is allowed to spread.





\subsection*{2H. Global pre-release reconciliation for v2.27.0}
With the first A-block now locally frozen through Lemma~C--G and Proposition~H--J, the remaining question is no longer whether the dependency spine exists, but whether this local gain has been fed back consistently into the document-wide release logic. The present pre-release reconciliation therefore records the first A-block not merely as locally strengthened, but as positively reconciled with the mathematics-first gate architecture for the next public release. In particular, the block is now to be read as structurally closed enough that its remaining limitation is no longer missing architecture, but only moderate proof density and future optional unfolding.

Operationally this means three things. First, the block may already count as a positive contribution to the global v2.27.0 gate, because it now satisfies the internal chain ground condition $	o$ HC $	o$ closure $	o$ projector stabilization $	o$ admissible transport $	o$ matrix realization $	o$ trace/branch persistence. Second, no computational teaser layer is needed in order to maintain its current mathematical status. Third, any later expansion of this block should be treated as proof-density enrichment rather than as a repair of a still-broken logical spine. The correct release reading is therefore conservative but positive: the first A-block is globally reconciled with the current mathematics-first release doctrine and should no longer be counted among the primary blockers for v2.27.0.

\paragraph{Global release consequence.}
For the next release decision, the first A-block is to be classified as \emph{mathematically release-supporting}: not yet maximally unfolded, but no longer structurally open. The remaining pre-release work should therefore shift away from re-opening this block and toward checking whether the remaining A-priority material can be brought into the same state without introducing theorem drift, reference instability, or dependence on the computational teaser architecture.

\paragraph{Local closure target.}
After this reconciliation pass, the active goal for the local subversion chain is no longer to add more managerial logic around the first A-block, but to preserve it as a frozen support block while the remaining release blockers are tested against the same gate discipline.



\subsection*{C.14 Fifteenth 3.4 extension: transition cocycle support}
\addcontentsline{toc}{subsection}{C.14 Fifteenth 3.4 extension: transition cocycle support}

\begin{theorem}[Appendix support for Theorem~\ref{thm:mainline-3415}]
\label{thm:appendix-3415}
Let $\mathfrak U=\{U_\alpha\}$ be the stable shell atlas and suppose each isolated obstruction shell carries one of the four local transition types from Theorem~\ref{thm:mainline-3414}. Then these local transitions assemble into a locally constant $V_4$-valued cocycle on admissible overlaps. After adjoining a sheet label that resolves fold exchange, the induced cocycle factors through the quotient $V_4\twoheadrightarrow \mathbb Z_2^{\mathrm{par}}$.
\end{theorem}

\begin{proof}
Every admissible overlap arises from continuation across finitely many isolated obstruction shells. By Theorem~\ref{thm:mainline-3414}, each such shell contributes one of the four involutive transition types. Their compositions therefore stay inside $V_4$, because the fold and parity generators commute and square to the identity. This gives a locally constant overlap label.

On triple overlaps, composing continuation maps in two stages or one stage leads to the same inverse shell identification, because both are obtained by following the same continuation class through the same finite obstruction data. Hence the cocycle relation holds. Resolving fold exchange by adjoining the local inverse-sheet label kills the fold generator on the lifted data, leaving only the parity component, which is exactly the quotient to $\mathbb Z_2^{\mathrm{par}}$.
\end{proof}

\paragraph{3.4.15 working interpretation.}
The inverse-readout atlas is therefore not just stratified; it is glued by a tiny and completely explicit transition law. Before fold resolution, the global bookkeeping lives in a Klein-four transition cocycle. After fold resolution, the only surviving topological memory is the parity lift. This is the most economical global form of the spinorial/M"obius mechanism obtained so far.




\subsection*{3.4 Sixteenth extension: transition cocycle and parity-resolved local system}
\addcontentsline{toc}{subsection}{3.4 Sixteenth extension: transition cocycle and parity-resolved local system}
The local obstruction law of Section~3.4.14 gives the primitive transition moves, but the next structural step is to organize those moves over the whole shell atlas rather than shell by shell. The natural object here is a tiny transition cocycle. It records how inverse shell data are re-identified when one passes between adjacent stable charts and then separates the fold ambiguity from the genuine parity lift.

\begin{theorem}[3.4 transition cocycle and parity-resolved local system]
\label{thm:mainline-3415}
Let $\mathfrak U=\{U_\alpha\}$ be the shell atlas from Theorem~\ref{thm:mainline-3413}, where each $U_\alpha$ is a connected component of $I\setminus\Sigma$ carrying stable inverse readout. Then the following hold.
\begin{enumerate}[label=(\roman*)]
  \item For every nonempty admissible overlap obtained by continuation across a single isolated obstruction shell, there is a locally constant transition element
  \[
    g_{\alpha\beta}\in V_4:=\{\mathrm{id},\,T^{\mathrm{fold}},\,T^{\mathrm{par}},\,T^{\mathrm{par}}\circ T^{\mathrm{fold}}\}
  \]
  such that inverse shell data satisfy
  \[
    (s,\pi)_\beta = g_{\alpha\beta}\,(s,\pi)_\alpha.
  \]
  \item On every admissible triple overlap one has the cocycle relation
  \[
    g_{\alpha\gamma}=g_{\beta\gamma}\circ g_{\alpha\beta},
  \]
  so the inverse-readout layer defines a $V_4$-valued local transition system over the shell atlas.
  \item After passing to the fold-resolved two-sheet inverse cover, the fold generator is absorbed into the choice of sheet, and the remaining transition system reduces to a parity local system valued in
  \[
    \mathbb Z_2^{\mathrm{par}}=\{\mathrm{id},T^{\mathrm{par}}\}.
  \]
\end{enumerate}
\end{theorem}

\begin{proof}
Statement (i) is exactly the content of Theorem~\ref{thm:mainline-3414}: every isolated obstruction contributes only fold exchange, parity flip, or their composition. Since those operations act discretely and the obstruction type cannot vary inside a sufficiently small continuation corridor, the corresponding transition element is locally constant on the overlap.

For (ii), the continuation from $U_\alpha$ to $U_\gamma$ through $U_\beta$ is represented by successive application of the two local transition maps. Because the local transitions are honest involutive re-identifications of the same inverse shell data, their composition agrees with direct continuation, which is the cocycle relation.

For (iii), resolving the fold ambiguity means that one fixes one of the two local inverse sheets as part of the enlarged state. The fold involution then acts trivially on the lifted data, because it has been absorbed into the sheet label itself. What remains is only whether the transported parity class is preserved or flipped, hence a residual $\mathbb Z_2$-valued local system.
\end{proof}

\begin{corollary}[3.4 parity monodromy class]
\label{cor:mainline-3415}
Under the hypotheses of Theorem~\ref{thm:mainline-3415}, every closed admissible continuation of inverse shell data determines a well-defined class in the parity local system $\mathbb Z_2^{\mathrm{par}}$. The visible direct-return case corresponds to the trivial class, while the spinorial/M"obius continuation is exactly the nontrivial parity class.
\end{corollary}

\begin{proof}
By Theorem~\ref{thm:mainline-3415}(iii), all fold ambiguity is removed on the lifted inverse cover, so a closed continuation can leave only the trivial or nontrivial parity action. The trivial action returns the same visible branch, whereas the nontrivial parity action returns only after the $\mathbb Z_2$ lift. This is precisely the spinorial/M"obius case already identified locally in Corollary~\ref{cor:mainline-3414}.
\end{proof}

\paragraph{Interpretive consequence.}
The shell atlas now carries a genuine global readout organization. Stable charts are not merely patched together ad hoc: they are glued by a tiny $V_4$-valued transition cocycle, and after resolving the fold choice the remaining global obstruction is purely parity-theoretic. This is the cleanest form so far of the claim that the visible spinorial/M"obius branch is not decorative language, but the nontrivial global parity sector of the inverse-readout geometry.




\subsection*{3.4 Seventeenth extension: parity double cover and lift-triviality criterion}
\addcontentsline{toc}{subsection}{3.4 Seventeenth extension: parity double cover and lift-triviality criterion}
The parity-resolved local system of Section~3.4.15 already identifies the spinorial/M"obius sector as a genuine global class. The next structural step is to say exactly when that class disappears and, when it does not, what the canonical repair is. The natural answer is a parity double cover: it trivializes precisely the residual $\mathbb Z_2$ obstruction and therefore converts the parity-resolved inverse readout into an honest global section.

\begin{theorem}[3.4 parity double cover and lift-triviality criterion]
\label{thm:mainline-3416}
Let $B$ be a connected admissible shell base carrying the parity local system $\mathcal L_{\mathrm{par}}$ from Theorem~\ref{thm:mainline-3415}. Then the following are equivalent.
\begin{enumerate}[label=(\roman*)]
  \item The parity local system $\mathcal L_{\mathrm{par}}$ is globally trivial on $B$.
  \item There exists a globally continuous parity-resolved branch choice
  \[
    \sigma:B\to\{\pm 1\}
  \]
  such that, after fold resolution, inverse shell data admit a globally single-valued visible representative
  \[
    (s,\sigma\sqrt{s},\pi)
  \]
  on $B$.
  \item Every closed admissible continuation in $B$ has trivial parity monodromy.
\end{enumerate}
If these equivalent conditions fail, then there exists a canonical connected double cover
\[
  p:\widetilde B\to B
\]
trivializing $\mathcal L_{\mathrm{par}}$. On $\widetilde B$, the parity-resolved inverse readout becomes globally single-valued, and every admissible closed continuation lifts to direct visible convergence.
\end{theorem}

\begin{proof}
The equivalence of (i) and (iii) is the standard characterization of a $\mathbb Z_2$ local system: triviality is exactly the vanishing of its monodromy on all closed admissible loops. If $\mathcal L_{\mathrm{par}}$ is trivial, one may choose a global parity frame, which is precisely the global sign selection in (ii). Conversely, any such global sign choice identifies the two parity fibers coherently over all of $B$, so the local system is trivial.

If the parity class is nontrivial, the associated principal $\mathbb Z_2$-bundle defines the canonical connected double cover $p:\widetilde B\to B$. Pulling back $\mathcal L_{\mathrm{par}}$ to $\widetilde B$ trivializes it by construction. Hence the parity-resolved inverse shell data admit a global branch choice upstairs, and every lifted closed continuation returns with trivial parity action, i.e. as direct visible convergence on the cover.
\end{proof}

\begin{corollary}[3.4 direct-return versus lift criterion]
\label{cor:mainline-3416}
On a connected admissible shell base, the direct visible-return regime is globally available if and only if the parity local system is trivial. The spinorial/M"obius regime is therefore not an extra local ambiguity, but exactly the nontrivial global parity class; it disappears on the canonical parity double cover.
\end{corollary}

\begin{proof}
The direct visible-return regime is exactly the existence of a globally coherent parity choice after fold resolution, which is condition~(ii) of Theorem~\ref{thm:mainline-3416}. The second claim is the nontrivial case of the same theorem.
\end{proof}

\paragraph{Interpretive consequence.}
The inverse-readout geometry now has a clean global dichotomy. Either the parity local system is already trivial, in which case the visible branch closes globally without further enlargement, or the residual obstruction is genuinely $\mathbb Z_2$ and the correct global object is the canonical parity double cover. In that sense, the spinorial/M"obius sector is no longer only a language for local continuation failure; it is exactly the statement that the visible readout trivializes only after passing to the natural two-sheeted cover.


\subsection*{3.4 Eighteenth extension: cohomological parity obstruction class}
\addcontentsline{toc}{subsection}{3.4 Eighteenth extension: cohomological parity obstruction class}
The parity double cover already gives the correct global repair. The next and cleanest reformulation is to record the remaining parity sector as a genuine cohomological obstruction. This does not add a new mechanism; it simply packages the already extracted $\mathbb Z_2$ data in the natural global invariant attached to the shell atlas.

\begin{theorem}[3.4 cohomological parity obstruction class]
\label{thm:mainline-3417}
Let $B$ be a connected admissible shell base equipped with the parity-resolved inverse-readout atlas from Sections~3.4.15--3.4.16. After fold resolution, let
\[
  g^{\mathrm{par}}_{\alpha\beta}\in \mathbb Z_2
\]
be the parity transition functions on an admissible shell atlas $\{U_\alpha\}$. Then:
\begin{enumerate}[label=(\roman*)]
  \item the family $\{g^{\mathrm{par}}_{\alpha\beta}\}$ is a \v{C}ech $1$-cocycle with coefficients in $\mathbb Z_2$;
  \item its cohomology class
  \item one has
  \[
  \begin{aligned}
    &\varepsilon_{\mathrm{par}}=0
    \quad\Longleftrightarrow\quad
    \text{the parity local system is trivial}\\[0.2em]
    &\quad\Longleftrightarrow\quad
    \text{direct visible convergence is globally available};
  \end{aligned}
  \]
  \item if $\varepsilon_{\mathrm{par}}\neq 0$, then the canonical parity double cover from Theorem~\ref{thm:mainline-3416} is the two-sheeted cover classified by $\varepsilon_{\mathrm{par}}$.
\end{enumerate}
\end{theorem}

\begin{proof}
Theorem~\ref{thm:mainline-3415} already reduces the shell-atlas transition data to the parity sector after fold resolution. Hence the remaining transition functions take values in the constant group $\mathbb Z_2$, and the cocycle relation on triple overlaps is exactly the cocycle identity for a \v{C}ech $1$-cochain. Changing the admissible atlas or parity frame modifies $g^{\mathrm{par}}_{\alpha\beta}$ by a \v{C}ech coboundary, so the cohomology class $\varepsilon_{\mathrm{par}}$ is well defined.

The equivalence in (iii) is the standard identification of a $\mathbb Z_2$ local system with its class in $H^1(B,\mathbb Z_2)$: vanishing means that the cocycle is cohomologous to zero, hence a global parity frame exists; this is exactly the direct visible-return regime from Theorem~\ref{thm:mainline-3416}. If the class does not vanish, the same theorem provides the canonical parity double cover trivializing the local system, and that cover is precisely the one classified by $\varepsilon_{\mathrm{par}}$.
\end{proof}

\begin{corollary}[3.4 parity obstruction as the global lift defect]
\label{cor:mainline-3417}
The spinorial/M\"obius regime is completely encoded by the single class
\[
  \varepsilon_{\mathrm{par}}\in H^1(B,\mathbb Z_2).
\]
It vanishes exactly in the direct visible-return regime and is nonzero exactly when the inverse readout trivializes only after passage to the canonical parity double cover.
\end{corollary}

\begin{proof}
This is the reformulation of Theorem~\ref{thm:mainline-3417}(iii)--(iv) in terms of the unique cohomological obstruction carried by the parity local system.
\end{proof}

\paragraph{Interpretive consequence.}
The parity side of the 3.4 program is now fully compressed. Local parity flips, shell-wise monodromy, the residual $\mathbb Z_2$ local system, the canonical double cover, and the distinction between direct visible closure and spinorial/M\"obius lift are no longer separate stories. They are all expressions of the single class
\[
  \varepsilon_{\mathrm{par}}\in H^1(B,\mathbb Z_2).
\]
In that sense, the remaining lift defect is now a genuine cohomological obstruction rather than merely a descriptive continuation phenomenon.


\subsection*{3.4 Nineteenth extension: parity holonomy character and loop-readout criterion}
\addcontentsline{toc}{subsection}{3.4 Nineteenth extension: parity holonomy character and loop-readout criterion}
The cohomological obstruction class from Section~3.4.17 is already the smallest global invariant. The next step is to make it directly readable along admissible closed shell paths. The natural object is the parity holonomy character: it evaluates the class on a loop and therefore tells us, in the smallest possible way, whether visible closure occurs already on the base or only after passage to the parity double cover.

\begin{theorem}[3.4 parity holonomy character and loop-readout criterion]
\label{thm:mainline-3418}
Let $B$ be a connected admissible shell base with parity obstruction class
\[
  \varepsilon_{\mathrm{par}}\in H^1(B,\mathbb Z_2)
\]
from Theorem~\ref{thm:mainline-3417}. Fix a base point $b_0\in B$. Then:
\begin{enumerate}[label=(\roman*)]
  \item there is a canonical parity holonomy character
  \[
    \chi_{\mathrm{par}}:\pi_1(B,b_0)\to \mathbb Z_2,
    \qquad
    \chi_{\mathrm{par}}([\gamma])=\langle \varepsilon_{\mathrm{par}},[\gamma]_{2}\rangle,
  \]
  where $[\gamma]_2\in H_1(B,\mathbb Z_2)$ denotes the mod-two homology class of $\gamma$;
  \item $\chi_{\mathrm{par}}([\gamma])=0$ if and only if every parity-resolved inverse continuation along $\gamma$ returns with direct visible closure on $B$;
  \item $\chi_{\mathrm{par}}([\gamma])=1$ if and only if continuation along $\gamma$ closes only after the nontrivial deck transformation of the canonical parity double cover from Theorem~\ref{thm:mainline-3416};
  \item the character factors through mod-two homology,
  \[
    \pi_1(B,b_0)\twoheadrightarrow H_1(B,\mathbb Z_2)\xrightarrow{\ \langle \varepsilon_{\mathrm{par}},-\rangle\ }\mathbb Z_2,
  \]
  so two admissible loops with the same class in $H_1(B,\mathbb Z_2)$ have the same parity readout.
\end{enumerate}
\end{theorem}

\begin{proof}
A $\mathbb Z_2$ local system determines a monodromy representation into $\mathbb Z_2$, and for the parity local system this is exactly the parity action picked up by admissible closed continuation. By Theorem~\ref{thm:mainline-3417}, the local system is classified by the cohomology class $\varepsilon_{\mathrm{par}}$. Standard evaluation of $H^1(B,\mathbb Z_2)$ on mod-two $1$-cycles gives the homological formula
\[
  \chi_{\mathrm{par}}([\gamma])=\langle \varepsilon_{\mathrm{par}},[\gamma]_2\rangle.
\]
The value $0$ means that the parity transport around $\gamma$ is trivial, which is exactly direct visible return on the base. The value $1$ means that the transported branch returns with the nontrivial parity action, so visible closure occurs only after passing to the parity double cover. Because the target group is abelian, the character factors through the abelianization of $\pi_1(B,b_0)$, and with mod-two coefficients this is precisely the factorization through $H_1(B,\mathbb Z_2)$.
\end{proof}

\begin{corollary}[3.4 loop-equivalence principle for parity readout]
\label{cor:mainline-3418}
On a connected admissible shell base, the global direct-return versus spinorial/M\"obius decision for any admissible closed continuation class depends only on the value of
\[
  \langle \varepsilon_{\mathrm{par}},[\gamma]_2\rangle\in\mathbb Z_2.
\]
In particular, homologous admissible loops mod $2$ are parity-indistinguishable for the inverse readout.
\end{corollary}

\begin{proof}
This is exactly statement~(iv) of Theorem~\ref{thm:mainline-3418}: the parity readout factors through the mod-two homology class of the loop.
\end{proof}

\paragraph{Interpretive consequence.}
The parity obstruction is now not only a global class but an operational readout. A closed admissible shell path carries a tiny $\mathbb Z_2$ outcome obtained by evaluating $\varepsilon_{\mathrm{par}}$ on its mod-two class. The visible direct-return regime corresponds to holonomy value $0$, while the spinorial/M\"obius regime corresponds to holonomy value $1$. This is the first fully global and directly readable form of the parity defect.



\subsection*{3.4 Twentieth extension: empirical parity-holonomy estimator and loop calibration}
\addcontentsline{toc}{subsection}{3.4 Twentieth extension: empirical parity-holonomy estimator and loop calibration}
The parity holonomy character from Section~3.4.18 is already the smallest global obstruction. The next step is to make it measurable from the readout side itself. Since the shell atlas already resolves admissible continuations into stable chart arcs separated by obstruction-shell crossings, the only possible source of a nontrivial parity outcome is the parity-changing part of those crossings. This yields a direct empirical estimator for the global $\mathbb Z_2$ holonomy.

\begin{theorem}[3.4 empirical parity-holonomy estimator]
\label{thm:mainline-3419}
Let $\gamma$ be an admissible closed shell path on a connected shell base $B$, decomposed into stable chart arcs whose endpoints meet obstruction shells transversely. For each encountered obstruction-shell crossing, let
\[
  \delta_k\in\mathbb Z_2
\]
be defined by
\[
  \delta_k=
  \begin{cases}
    0,&\text{for a pure fold transition},\\
    1,&\text{for a parity transition},\\
    1,&\text{for a mixed fold-plus-parity transition},
  \end{cases}
\]
so that $\delta_k$ records only the parity-changing component of the local transition law from Theorem~\ref{thm:mainline-3414}. Define the empirical parity estimator
\[
  \widehat\chi_{\mathrm{par}}(\gamma)
  :=
  \sum_{k=1}^{N_\gamma}\delta_k \pmod 2.
\]
Then:
\begin{enumerate}[label=(\roman*)]
  \item $\widehat\chi_{\mathrm{par}}(\gamma)$ depends only on the parity labels of the stable chart arcs and is independent of the chosen admissible chart decomposition;
  \item the empirical estimator agrees with the cohomological holonomy character,
  \[
    \widehat\chi_{\mathrm{par}}(\gamma)
    =
    \chi_{\mathrm{par}}([\gamma])
    =
    \langle \varepsilon_{\mathrm{par}},[\gamma]_2\rangle;
  \]
  \item consequently, $\widehat\chi_{\mathrm{par}}(\gamma)=0$ if and only if the calibrated inverse readout closes directly on the base, while $\widehat\chi_{\mathrm{par}}(\gamma)=1$ if and only if closure occurs only after the nontrivial deck action on the parity double cover.
\end{enumerate}
\end{theorem}

\begin{proof}
By Theorem~\ref{thm:mainline-3414}, only the parity component of an obstruction-shell crossing contributes to the post-crossing branch label; pure folds only exchange inverse sheets and therefore do not affect the residual $\mathbb Z_2$ parity class. Summing the parity-changing contributions modulo $2$ along the closed path therefore computes the total parity transport accumulated by the continuation. Refining the chart decomposition inserts pairs of complementary local transitions whose parity contributions cancel modulo $2$, so the quantity $\widehat\chi_{\mathrm{par}}(\gamma)$ is well defined.

After fold resolution, the remaining transition data is exactly the parity local system from Theorem~\ref{thm:mainline-3415}. Hence the mod-two sum of parity changes along $\gamma$ is the monodromy of that local system around the loop, which by Theorem~\ref{thm:mainline-3418} is precisely $\chi_{\mathrm{par}}([\gamma])=\langle \varepsilon_{\mathrm{par}},[\gamma]_2\rangle$. The final visible-return statements are the two cases of that holonomy value.
\end{proof}

\begin{corollary}[3.4 detector-side loop test for the parity obstruction]
\label{cor:mainline-3419}
Suppose a calibrated inverse readout assigns stable chart arcs and parity labels to a detector-side closed shell trace. Then the parity obstruction can be tested directly by the mod-two count of parity-changing crossings:
\[
  \widehat\chi_{\mathrm{par}}(\gamma)=0
  \Longleftrightarrow
  \text{detector-side direct visible return},
\]
\[
  \widehat\chi_{\mathrm{par}}(\gamma)=1
  \Longleftrightarrow
  \text{detector-side spinorial/M\"obius lift signal}.
\]
In particular, the global $\mathbb Z_2$ lift defect becomes empirically checkable once the shell trace has been chartwise calibrated.
\end{corollary}

\begin{proof}
This is exactly statement~(iii) of Theorem~\ref{thm:mainline-3419}, read through the calibrated inverse readout from Sections~3.4.11--3.4.13.
\end{proof}

\paragraph{Interpretive consequence.}
The parity side now closes the loop between geometry and readout. The cohomological class is no longer only a classification object and the holonomy character is no longer only an abstract loop functional. Once a detector-side or probe-side shell trace has been calibrated chartwise, the global lift defect is measured by a single mod-two crossing count. This is the sharpest operational form so far of the statement that the spinorial/M\"obius regime is a real readout phenomenon rather than a decorative analogy.


\subsection*{C.15 Sixteenth 3.4 extension: parity double-cover support}
\addcontentsline{toc}{subsection}{C.15 Sixteenth 3.4 extension: parity double-cover support}

\begin{theorem}[Appendix support for Theorem~\ref{thm:mainline-3416}]
\label{thm:appendix-3416}
Let $B$ be a connected admissible shell base and let $\mathcal L_{\mathrm{par}}$ be the parity local system obtained after fold resolution. Then the parity monodromy determines a unique connected two-sheeted cover $p:\widetilde B\to B$ up to isomorphism. Pulling back $\mathcal L_{\mathrm{par}}$ along $p$ trivializes the parity transport, and any other connected cover with the same property factors through $\widetilde B$.
\end{theorem}

\begin{proof}
The parity local system is a locally constant $\mathbb Z_2$-bundle over $B$. Its monodromy representation sends every admissible closed continuation class to either the trivial or nontrivial parity action. The associated kernel subgroup therefore defines a connected index-two cover $p:\widetilde B\to B$, unique up to canonical isomorphism. By construction, loops lifted to $\widetilde B$ lie in the kernel of the parity monodromy, so the pulled-back local system has trivial transport. If another connected cover trivializes the same monodromy, its deck group must also kill that kernel and therefore factor through the index-two quotient defining $\widetilde B$.
\end{proof}

\paragraph{3.4.16 working interpretation.}
The parity sector has now reached its cleanest possible global form. The shell atlas first produced local obstruction types, then a Klein-four transition cocycle, then a reduced parity local system, and now finally the canonical object that removes the remaining obstruction: the parity double cover. Once lifted to that cover, the inverse-readout branch is globally single-valued again. This is the sharpest mathematical version so far of the statement that the visible spinorial/M"obius branch is a double-cover phenomenon rather than a decorative metaphor.


\subsection*{C.16 Seventeenth 3.4 extension: cohomological obstruction support}
\addcontentsline{toc}{subsection}{C.16 Seventeenth 3.4 extension: cohomological obstruction support}

\begin{theorem}[Appendix support for Theorem~\ref{thm:mainline-3417}]
\label{thm:appendix-3417}
Let $\{U_\alpha\}$ be an admissible shell atlas on $B$, and let $g^{\mathrm{par}}_{\alpha\beta}\in\mathbb Z_2$ be the parity transition cocycle after fold resolution. Then the cohomology class
\[
  \varepsilon_{\mathrm{par}}=[g^{\mathrm{par}}_{\alpha\beta}]\in H^1(B,\mathbb Z_2)
\]
classifies the parity local system and its associated connected two-sheeted cover up to isomorphism.
\end{theorem}

\begin{proof}
Because the post-fold transition functions take values in the constant group $\mathbb Z_2$, they define a \v{C}ech $1$-cocycle. Two admissible atlases related by refinement or change of parity frame produce cohomologous cocycles, so the resulting class in $H^1(B,\mathbb Z_2)$ is intrinsic. Standard classification of principal $\mathbb Z_2$-bundles over $B$ identifies this cohomology class with the isomorphism class of the parity local system. The associated connected two-sheeted cover is the cover classified by the same class and is exactly the parity double cover constructed in Theorem~\ref{thm:appendix-3416}.
\end{proof}

\paragraph{3.4.17 working interpretation.}
The remaining global defect has now been compressed to its sharpest invariant. What previously appeared as parity flips, spinorial lift behavior, and double-cover repair is encoded by one class in $H^1(B,\mathbb Z_2)$. This does not change the mechanism; it shows that the mechanism is already as small as it can be.


\subsection*{C.17 Eighteenth 3.4 extension: parity holonomy support}
\addcontentsline{toc}{subsection}{C.17 Eighteenth 3.4 extension: parity holonomy support}

\begin{theorem}[Appendix support for Theorem~\ref{thm:mainline-3418}]
\label{thm:appendix-3418}
Let $\mathcal L_{\mathrm{par}}$ be the parity local system on a connected admissible shell base $B$, classified by
\[
  \varepsilon_{\mathrm{par}}\in H^1(B,\mathbb Z_2).
\]
Then the monodromy representation of $\mathcal L_{\mathrm{par}}$ is exactly the homomorphism obtained by evaluating $\varepsilon_{\mathrm{par}}$ on mod-two loop classes. Its kernel consists precisely of those admissible closed continuation classes that lift to direct visible closure on the canonical parity double cover.
\end{theorem}

\begin{proof}
A locally constant $\mathbb Z_2$-bundle on $B$ has monodromy in the discrete group $\mathbb Z_2$, hence a homomorphism from $\pi_1(B,b_0)$ to $\mathbb Z_2$. By Theorem~\ref{thm:appendix-3417}, the same local system is classified by the cocycle class $\varepsilon_{\mathrm{par}}\in H^1(B,\mathbb Z_2)$. Under the standard identification of $H^1(B,\mathbb Z_2)$ with homomorphisms out of mod-two $1$-cycles, the monodromy around a loop is given by evaluation of $\varepsilon_{\mathrm{par}}$ on its mod-two homology class. The lifted double cover from Theorem~\ref{thm:appendix-3416} trivializes the parity transport, so the loops in the kernel are exactly those whose lifted continuation closes directly on the cover.
\end{proof}

\paragraph{3.4.18 working interpretation.}
The cohomological class now has a direct operational meaning. It is not only something that classifies the parity defect abstractly; it can be evaluated on a closed admissible shell path and returns a two-valued global outcome. This is the cleanest bridge so far between the atlas/cohomology side and the actual readout side.



\subsection*{C.18 Nineteenth 3.4 extension: empirical parity-holonomy support}
\addcontentsline{toc}{subsection}{C.18 Nineteenth 3.4 extension: empirical parity-holonomy support}

\begin{theorem}[Appendix support for Theorem~\ref{thm:mainline-3419}]
\label{thm:appendix-3419}
Let $\gamma$ be an admissible closed shell path decomposed into stable chart arcs meeting obstruction shells transversely. Then the mod-two count of parity-changing crossings is equal to the monodromy of the parity local system around $\gamma$, and hence equal to the evaluation of the parity obstruction class on the mod-two loop class of $\gamma$.
\end{theorem}

\begin{proof}
Each local obstruction crossing carries one of the four transition types from Theorem~\ref{thm:mainline-3414}. Passing to the post-fold parity local system collapses the fold component and retains only whether the crossing changes the residual parity label. Therefore the transported parity label after traversing the full loop is obtained by multiplying the local $\mathbb Z_2$ transition factors, which in additive notation is exactly the mod-two sum of the local parity-change indicators $\delta_k$. This product is the monodromy of the parity local system along the loop. By Theorem~\ref{thm:appendix-3418}, that monodromy is also the evaluation of $\varepsilon_{\mathrm{par}}$ on the mod-two homology class of $\gamma$.
\end{proof}

\paragraph{3.4.19 working interpretation.}
The parity obstruction can now be read directly from a calibrated shell trace. Once the trace has been decomposed into admissible chart arcs, the full global lift defect collapses to a mod-two count of parity-changing transitions. This is the first point at which the global cohomological obstruction becomes an explicit detector-side or probe-side measurement rule.


\subsection*{C.19 Twentieth 3.4 extension: closing synthesis support}
\addcontentsline{toc}{subsection}{C.19 Twentieth 3.4 extension: closing synthesis support}

\begin{theorem}[Appendix support for Theorem~\ref{thm:mainline-3420}]
\label{thm:appendix-3420}
The stabilized 3.4 line factors through a single visible package
\[
  \bigl(\Phi,\Delta_{\mathrm{diag}},\varepsilon_{\mathrm{par}},\chi_{\mathrm{par}},\widehat\chi_{\mathrm{par}}\bigr)
\]
with branch variable $\eta$, in the sense that every local regular shell observable and every global parity-readout statement used in the mainline 3.4 section is a consequence of these data together with the shell atlas away from the obstruction set.
\end{theorem}

\begin{proof}
Theorem~\ref{thm:mainline-346} reduces the visible branch dynamics to the analytic generator $\Phi(\eta^2)$, and Theorems~\ref{thm:appendix-3411}--\ref{thm:appendix-3413} show that calibrated shell inversion away from the obstruction set is entirely controlled by the shell coordinate and parity label, with $\Delta_{\mathrm{diag}}$ measuring chartwise mismatch. Theorems~\ref{thm:appendix-3415}--\ref{thm:appendix-3418} then compress the remaining global ambiguity to the parity local system, its cohomological class, and its holonomy character. Finally, Theorem~\ref{thm:appendix-3419} identifies the empirical estimator with this same holonomy class along admissible closed traces. Hence the listed package captures the entire visible 3.4 layer.
\end{proof}

\paragraph{3.4.20 working interpretation.}
The 3.4 material is now closed as a single generator-atlas-holonomy block. Local residual dynamics is carried by the one-field generator, local diagnostic mismatch is measured by $\Delta_{\mathrm{diag}}$, and the only global leftover is the parity obstruction class with its loop readout. This is the natural handoff point for the next stage.


\subsection*{C.20 First 3.5 extension: residual-record support}
\addcontentsline{toc}{subsection}{C.20 First 3.5 extension: residual-record support}

\begin{theorem}[Appendix support for Theorem~\ref{thm:mainline-3500}]
\label{thm:appendix-3500}
Let $B$ be a connected admissible shell base satisfying the stabilized 3.4 hypotheses. Then the visible branch layer admits a faithful record description by
\[
  \mathfrak R(B)=\bigl(\Phi,\Delta_{\mathrm{diag}},\varepsilon_{\mathrm{par}},\chi_{\mathrm{par}},\widehat\chi_{\mathrm{par}}\bigr),
\]
in the sense that every chartwise regular-shell observable and every global parity-readout statement in the stabilized 3.4 sector is recoverable from this record together with the chartwise shell coordinate and parity label.
\end{theorem}

\begin{proof}
Theorem~\ref{thm:appendix-3420} already proves that the full stabilized 3.4 line factors through the same five visible data together with the shell atlas away from the obstruction set. The generator $\Phi$ determines the local branch law, the inverse shell charts determine the regular-shell coordinate and parity label, $\Delta_{\mathrm{diag}}$ measures chartwise diagnostic mismatch, and the parity package records the only remaining global ambiguity together with its operational loop estimator. Therefore the visible branch layer is faithfully encoded by the record $\mathfrak R(B)$, and the only extra ingredients needed for local reconstruction are the chartwise shell coordinate and parity label already present in the inverse atlas.
\end{proof}

\paragraph{3.5.0 working interpretation.}
The first simplification step does not yet introduce a new formal language. It only observes that the visible branch layer has already closed enough to be handled as a single admissible record. This is exactly the right level of compression before any later language-level or proof-economy layer is introduced. In the current status reading, this also marks the point after which later progress should be judged mainly by record-level sharpening, not by adding fresh visible ingredients.



\subsection*{C.21 Second 3.5 extension: residual-record morphism support}
\addcontentsline{toc}{subsection}{C.21 Second 3.5 extension: residual-record morphism support}

\begin{theorem}[Appendix support for Theorem~\ref{thm:mainline-3501}]
\label{thm:appendix-3501}
Let $F:B_1\to B_2$ be a residual-record isomorphism between connected admissible shell bases. Then the entire stabilized visible generator-atlas-holonomy sector is preserved under transport along $F$: the transported generator law, calibrated inverse shell data, shell-consistency diagnostics, parity obstruction class, parity holonomy, and empirical loop estimator all remain faithful to the same residual record.
\end{theorem}

\begin{proof}
By Theorem~\ref{thm:appendix-3500}, the stabilized visible sector factors through the admissible record together with chartwise shell coordinates. The definition of residual-record isomorphism requires preservation of exactly these ingredients at the level at which they enter the reconstruction normal form: the generator data, the shell-consistency defect, and the parity package. Therefore every chartwise visible formula and every global parity-readout statement already proved in the stabilized sector is transported faithfully along $F$, and no new visible branch datum is opened by this translation.
\end{proof}

\paragraph{3.5.1 working interpretation.}
The first move beyond the residual record is not to enlarge the visible layer, but to declare when two residual presentations should count as the same. This makes later compression work more disciplined: any future formal front-end must preserve the record-level morphism class if it is to remain faithful to the present mathematics. Historically open blocks that survive beyond this point therefore survive as residual identification or quantitative tasks, not as unresolved questions about what the visible data are.


\subsection*{C.22 Third 3.5 extension: residual-record interpreter support}
\addcontentsline{toc}{subsection}{C.22 Third 3.5 extension: residual-record interpreter support}

\begin{theorem}[Appendix support for Theorem~\ref{thm:mainline-3502}]
\label{thm:appendix-3502}
Let $B$ be a connected admissible shell base in the stabilized generator-atlas-holonomy regime. Then every visible statement in this regime is representable by a record-level formula built from
\[
  \mathfrak R(B)=\bigl(\Phi,\Delta_{\mathrm{diag}},\varepsilon_{\mathrm{par}},\chi_{\mathrm{par}},\widehat\chi_{\mathrm{par}}\bigr)
\]
together with chartwise shell coordinate and parity label, and this representation is unique up to residual-record isomorphism.
\end{theorem}

\begin{proof}
Theorem~\ref{thm:appendix-3500} already establishes that the entire stabilized visible branch layer factors through the admissible residual record together with shell-chart data. Theorem~\ref{thm:appendix-3501} shows that these data are preserved exactly under residual-record isomorphisms. Hence every visible theorem in the stabilized sector may be rewritten as a formula over the record package and shell-chart coordinates, and any two such rewritings can differ only by a residual-record isomorphism. This proves exact representability at record level.
\end{proof}

\paragraph{3.5.2 working interpretation.}
The meaning of the interpreter step is simple: from this point onward, the stabilized visible theory has a fixed admissible translation interface. Any future compression layer may become more economical or more set-theoretically explicit, but it should still be judged by whether it preserves the same record-level content. This is also the right point to separate the later live fronts from the historical OP labels: the surviving questions concern residual corrections, controlled emergence, spectral identification, Dixon compression, and the Theta-sum remainder, while OP~2 as a separate operator problem and OP~3 as a primary live front are no longer active blockers.



\subsection*{C.23 First 3.6 extension: conservative frontend support}
\addcontentsline{toc}{subsection}{C.23 First 3.6 extension: conservative frontend support}

\begin{theorem}[Appendix support for Theorem~\ref{thm:mainline-3600}]
\label{thm:appendix-3600}
Let $(\mathrm{Fr}(B),\mathcal F_B,\mathcal J_B)$ be a conservative frontend over the stabilized visible sector of a connected admissible shell base $B$. Then every frontend statement coming from a visible statement is fixed, up to unique record-level translation, by the same residual-record content as the original visible statement.
\end{theorem}

\begin{proof}
Theorem~\ref{thm:appendix-3502} already identifies the exact record-level representative of every visible theorem in the stabilized sector. The defining factorization rule for a conservative frontend requires frontend presentation to pass through that same representative, while transport-compatibility and primitive non-drift forbid any deviation on the visible image. Hence the frontend contributes only a conservative presentation layer and no second visible basis.
\end{proof}

\paragraph{3.6.0 working interpretation.}
The gain of the first 3.6 step is methodological but real: later language work is no longer a vague future promise. It now has a theorem-level admissibility test. This is precisely what keeps the planned ZFM-/frontend line subordinate to the mathematics-first strategy of the Companion.


\clearpage
\subsection*{Appendix C.124. Visual frame-flicker diagnostics for OP~5}
\addcontentsline{toc}{subsection}{Appendix C.124. Visual frame-flicker diagnostics for OP~5}
The following figures make the Section~3.8.48 hypothesis auditable on the current benchmark window. They do not prove the frame-flicker ansatz, but they show that a single smooth packet-scale moving-frame mode can already account for most of the mismatch beyond the rigid carrier and can do so in the two correction channels predicted by the current first-order frame-flicker decomposition.

\begin{remark}[Systems-level reading of the OP~5 visual block]\label{rem:appendixC-op5-systems-reading}
In the systems language summarized in Figure~\texttt{fig:qam-vm-os-zfs} in the active-core PDF, the present OP~5 figures should be read as visible-support views rather than as independent kernel identities. The rigid carrier, the packet breathing, and the localized readout slip are therefore shown here precisely because they belong to the auditable view layer of one transported channel. Their role is to expose where the visible readout departs from the rigid carrier picture, while the structural closure judgment itself still remains delegated to the deeper branch/shell/transport architecture.
\end{remark}


\begin{remark}[Visual-to-proof workflow for the OP~5 block]\label{rem:appendixC-op5-visual-workflow}
The present OP~5 figures should therefore be used in the following order. First, they identify a visible defect pattern at the readout layer. Second, they motivate a packet/carrier/slip decomposition that can be audited channelwise. Third, that decomposition must be checked against the active transport and closure-preservation architecture rather than treated as a proof by itself. Only after such a check has been made may the same pattern be recorded as a later kernel-first compression candidate. In particular, the figures are admissible as \emph{diagnostic entry points} and as \emph{audit surfaces}, but not as autonomous closure theorems.
\end{remark}

\begin{remark}[What the OP~5 figures are allowed to certify]\label{rem:appendixC-op5-allowed-certification}
At the present stage, the OP~5 visual block may certify only that a single packet-scale moving-frame correction remains empirically competitive as a visible-support description on the chosen benchmark window, and that the induced residual can be organized in the two currently tracked correction channels. It may not certify uniqueness of the correction mechanism, global closure of the OP~5 hypothesis, or activation of a new transport rule beyond the active theorem chain.
\end{remark}


\begin{figure}[H]
\centering
\IfFileExists{op5_frame_flicker_fit_v3862.pdf}{\includegraphics[width=0.98\textwidth]{op5_frame_flicker_fit_v3862.pdf}}{\safegraphic[width=0.98\textwidth]{op5_frame_flicker_fit_v3862.png}}
\caption{\textbf{OP~5 packet/readout fit on the current benchmark window.} Upper panel: the observed OP~5 packet/readout trace, the rigid carrier baseline, and the first-order packet/readout correction inside the visible packet-support window. Lower panel: the residual beyond the rigid carrier split into the breathing-envelope channel \(\alpha\psi_{N,m}\sin\Theta_{N,m}\) and the readout-slip channel \(\beta A_{N,m}\psi_{N,m}\cos\Theta_{N,m}\).}
\label{fig:appendixC-op5-frame-flicker-fit-live}
\end{figure}

\begin{figure}[H]
\centering
\IfFileExists{op5_frame_flicker_components_v3862.pdf}{\includegraphics[width=0.98\textwidth]{op5_frame_flicker_components_v3862.pdf}}{\safegraphic[width=0.98\textwidth]{op5_frame_flicker_components_v3862.png}}
\caption{\textbf{Packet-support mode, visible channels, and incremental OP~5 build-up.} The four panels show the packet-support window \(\psi_{N,m}\), the rigid carrier/readout quadratures, the two visible channels fed by one packet-scale moving-frame mode, and the stepwise build-up from rigid carrier to packet/breathing correction and then to the full OP~5 readout trace.}
\label{fig:appendixC-op5-frame-flicker-components-live}
\end{figure}

\begin{figure}[H]
\centering
\IfFileExists{op5_frame_flicker_geometry_v3862.pdf}{\includegraphics[width=0.98\textwidth]{op5_frame_flicker_geometry_v3862.pdf}}{\safegraphic[width=0.98\textwidth]{op5_frame_flicker_geometry_v3862.png}}
\caption{\textbf{Geometric OP~5 reading: carrier helix, breathing packet, and readout slip.} The panels compare the rigid carrier helix with the packet/readout-distorted helix, show the top-view deviation, display the packet breathing envelope along the helix, and isolate the localized readout phase slip. The natural reading is therefore not a second bulk spiral, but the same carrier seen through one softened visible frame.}
\label{fig:appendixC-op5-frame-flicker-geometry-live}
\end{figure}

\noindent\textit{Reading guide.} Figures~\ref{fig:appendixC-op5-frame-flicker-fit-live}--\ref{fig:appendixC-op5-frame-flicker-geometry-live} should be read in order as \emph{trace $\rightarrow$ channel decomposition $\rightarrow$ geometry}. In that sequence, the packet/readout defect moves from the visible time trace, to the two correction channels, to the helical geometric picture.

\paragraph{Appendix conclusion.}
On the present benchmark window, the frame-flicker ansatz remains hypothesis-grade, but it is no longer only verbal. A single smooth moving-frame mode reproduces most of the residual correction beyond the rigid carrier and does so in exactly the two channels predicted by the first-order decomposition. This strengthens the current reading that OP~5 is better treated as a carrier-plus-readout-defect problem than as a missing second rigid transport wave.

\begin{remark}[Packet/flicker diagnostics as readout evidence]
Within the current OP transport architecture, the packet-wave, breathing-envelope, and frame-flicker figures should be read as visible-support diagnostics rather than as primary closure criteria. Their mathematical role is to show that a large part of the observed OP~5 residue is compatible with a carrier-plus-readout-defect picture, which is exactly the kind of situation filtered later by the visible/structural mismatch separation tool.
\end{remark}

\paragraph{Bibliography placement note.}
The bibliography that formerly appeared at this historical appendix position has been moved to the final back matter so that the document closes with the references. The version-record and historical-release material continues below unchanged as archival proof-memory.






\subsection*{Working common-mode carrier is downstream from the active HC/MM/J fit}

The previous subsection reduced the primitive OP~1/OP~5 obstruction to a common-mode carrier plus a downstream observer-shell split. The next structural step is to push that common-mode carrier itself back to the active HC/MM/J fit. On the active admissible branch, the shared carrier is not an independent visible object: it is simply the shared visible mask applied to the complementary output generated by the active matrix-leakage operator. In particular, if the active branch is exactly HC/MM/J-compatible, then the primitive common-mode carrier must vanish identically.

\begin{definition}[Working active matrix-fit defect and shared carrier readout]
Let $U:=\mathrm{span}\{u_1,u_2\}$ be the active admissible plane and $W:=\mathrm{span}\{w_1,w_2\}$ the active complementary leakage plane. Define the active matrix-fit defect by
\[
  E_{\mathrm{fit}}^{(U)} := \bigl((\mathrm{id}-C)\mathcal M C\bigr)\big|_{U}: U\to W.
\]
For a unit active probe state
\[
  x(\theta):=\cos\theta\,u_1+\sin\theta\,u_2\in U,
\]
define the active hidden leakage response and the shared visible carrier by
\[
  y_{\mathrm{fit}}(\theta):=E_{\mathrm{fit}}^{(U)}x(\theta)\in W,
  \qquad
  c_{\mathrm{sh}}^{\mathrm{fit}}(\theta):=L_{\mathrm{sh}}^{(W)}\bigl(y_{\mathrm{fit}}(\theta)\bigr).
\]
\end{definition}

\begin{proposition}[Working exact commutator identity on the admissible branch]
For every $x\in U\subseteq \mathrm{im}(C)$ one has
\[
  (\mathrm{id}-C)\mathcal M x
  =
  \frac12\,[\mathcal M,J]x.
\]
Equivalently,
\[
  E_{\mathrm{fit}}^{(U)}
  =
  \frac12\,[\mathcal M,J]\big|_{U}.
\]
In particular,
\[
  \|E_{\mathrm{fit}}^{(U)}\|
  \le
  \frac12\,\|[\mathcal M,J]\big|_{U}\|.
\]
\end{proposition}

\begin{proof}
Since $U\subseteq \mathrm{im}(C)$, every $x\in U$ satisfies $Cx=x$ and $Jx=x$. Therefore
\[
  (\mathrm{id}-C)\mathcal M x
  =
  \frac12(\mathrm{id}-J)\mathcal M x
  =
  \frac12\bigl(\mathcal Mx-J\mathcal Mx\bigr).
\]
Using $Jx=x$, this becomes
\[
  \frac12\bigl(\mathcal MJx-J\mathcal Mx\bigr)
  =
  \frac12\,[\mathcal M,J]x.
\]
Restricting to $U$ gives the operator identity, and the norm bound is immediate.
\end{proof}

\begin{theorem}[Working shared carrier is controlled by the active HC/MM/J fit]
For every active probe angle $\theta$ one has
\[
  c_{\mathrm{sh}}^{\mathrm{fit}}(\theta)
  =
  L_{\mathrm{sh}}^{(W)}\bigl(E_{\mathrm{fit}}^{(U)}x(\theta)\bigr)
  =
  \frac12\,L_{\mathrm{sh}}^{(W)}\bigl([\mathcal M,J]x(\theta)\bigr).
\]
Hence
\[
  |c_{\mathrm{sh}}^{\mathrm{fit}}(\theta)|
  \le
  \|L_{\mathrm{sh}}^{(W)}\|_{W^*}\,\|E_{\mathrm{fit}}^{(U)}\|\,\|x(\theta)\|
  \le
  \frac12\,\|L_{\mathrm{sh}}^{(W)}\|_{W^*}\,\|[\mathcal M,J]\big|_{U}\|.
\]
In particular, if
\[
  [\mathcal M,J]\big|_{U}=0,
\]
then
\[
  c_{\mathrm{sh}}^{\mathrm{fit}}(\theta)=0
\]
for all $\theta$, so the primitive common-mode OP~1/OP~5 carrier vanishes on the active branch.
\end{theorem}

\begin{proof}
The first displayed identity is the definition of $c_{\mathrm{sh}}^{\mathrm{fit}}$. The second follows from the previous proposition. Since $\|x(\theta)\|=1$ on the active unit circle in $U$, the norm estimate follows immediately. If the restricted commutator vanishes, then $E_{\mathrm{fit}}^{(U)}=0$, hence $y_{\mathrm{fit}}(\theta)=0$ and therefore $c_{\mathrm{sh}}^{\mathrm{fit}}(\theta)=0$ for all $\theta$.
\end{proof}

\begin{corollary}[Working primitive common-mode residual is a pure matrix-fit obstruction]
Assume exact grouped no-drift on the active leakage plane, so that $\Delta L^{(W)}=0$. Then every surviving primitive OP~1/OP~5 residual on the active branch is completely controlled by the active HC/MM/J fit defect:
\[
  d(\theta)=c_{\mathrm{sh}}^{\mathrm{fit}}(\theta)\binom{1}{1},
  \qquad
  |c_{\mathrm{sh}}^{\mathrm{fit}}(\theta)|
  \le
  \frac12\,\|L_{\mathrm{sh}}^{(W)}\|_{W^*}\,\|[\mathcal M,J]\big|_{U}\|.
\]
In particular, under exact grouped no-drift and exact active commutation with the umkehrwalze, the primitive OP~1/OP~5 channel closes identically on the active branch.
\end{corollary}

\begin{proof}
Under exact grouped no-drift one has $c_{\mathrm{obs}}\equiv 0$ by the previous shared-residual dichotomy, so only the common-mode component remains. The displayed bound is the previous theorem. If the restricted commutator also vanishes, then the common-mode carrier vanishes identically, hence the primitive visible channel closes.
\end{proof}

\paragraph{Working interpretation.}
This is the clean matrix-side endpoint of the current reduction. Once the observer-shell split is removed on the active leakage plane, the remaining primitive OP~1/OP~5 carrier is not a mysterious visible residue anymore: it is exactly the shared mask applied to the active HC/MM/J commutator defect. So the primitive common-mode channel survives only to the extent that the Millennium matrix still fails to commute with the umkehrwalze on the active admissible branch.



\subsection*{Working architectural promotion of the active commutator defect}

The active commutator defect was introduced as the primitive common-mode obstruction on the active branch. However, the older grouped-layer architecture already contains a stronger structural principle: admissible Millennium transport is only allowed when it intertwines with the involution. Therefore the active commutator should not be read as a free new defect by default, but as a quantity that must vanish whenever the transport really is the admissible grouped-layer converter claimed by the document architecture. What then remains, if anything, must be downstream dressing rather than primitive HC/MM/J leakage.

\begin{theorem}[Working admissible grouped transport kills the active commutator]
Assume the Millennium matrix $\mathcal M$ acts on the active admissible plane $U$ as an admissible grouped-layer converter in the sense of Definition~\texttt{def:millennium-transport-grouped-layers} in the active-core PDF, and assume its transport role is structure preserving with respect to the involution/umkehrwalze in the sense of Lemma~\texttt{lem:millennium-intertwining-lemma} in the active-core PDF. Then
\[
  [\mathcal M,J]\big|_{U}=0.
\]
Consequently,
\[
  E_{\mathrm{fit}}^{(U)}=0,
  \qquad
  c_{\mathrm{sh}}^{\mathrm{fit}}(\theta)=0
  \quad\text{for all }\theta.
\]
\end{theorem}

\begin{proof}
By Lemma~\texttt{lem:millennium-intertwining-lemma} in the active-core PDF, every Millennium transport that claims to preserve the normalized involutive kernel must satisfy
\[
  \mathcal M\mathcal J=\mathcal J\mathcal M.
\]
On the active branch the visible involution is realized by the operator $J$, so the same intertwining statement on the active branch reads
\[
  \mathcal M Jx=J\mathcal M x
  \qquad (x\in U).
\]
Hence
\[
  [\mathcal M,J]x=0
  \qquad (x\in U),
\]
which is exactly the claimed vanishing of the restricted commutator. The identities
\[
  E_{\mathrm{fit}}^{(U)}=\tfrac12[\mathcal M,J]\big|_{U}
\]
and
\[
  c_{\mathrm{sh}}^{\mathrm{fit}}(\theta)=L_{\mathrm{sh}}^{(W)}\bigl(E_{\mathrm{fit}}^{(U)}x(\theta)\bigr)
\]
were established in the previous subsection, so both vanish immediately.
\end{proof}

\begin{corollary}[Working primitive OP~1/OP~5 common mode cannot survive exact admissible transport]
Assume exact grouped no-drift on the active leakage plane and admissible grouped Millennium transport on the active branch. Then the primitive visible OP~1/OP~5 channel closes identically on that branch:
\[
  d(\theta)=0
  \qquad\text{for all }\theta.
\]
In particular, any still visible residual must be interpreted as downstream observer-shell dressing, phase-shifted readout, or some other post-transport effect rather than as a primitive HC/MM/J leakage mode.
\end{corollary}

\begin{proof}
Exact grouped no-drift gives $c_{\mathrm{obs}}(\theta)=0$ by the shared-residual dichotomy. The previous theorem gives $c_{\mathrm{sh}}^{\mathrm{fit}}(\theta)=0$. Since the primitive visible channel decomposes as the sum of the common-mode part and the observer-split part, both pieces vanish and therefore $d(\theta)=0$ for all $\theta$.
\end{proof}

\begin{remark}[Working interpretation: exact closure versus canonical observer-shell residual]
The current architecture therefore supports a sharp dichotomy.
\begin{enumerate}[leftmargin=2em,label=(\roman*)]
\item If the Millennium matrix really acts as the admissible grouped-layer converter claimed by the structural architecture, then the active commutator vanishes on the active branch and no primitive common-mode OP~1/OP~5 leakage survives there.
\item If a residual nevertheless remains visible, then it should no longer be hunted as a primitive HC/MM/J closure defect. It must instead be read as a canonical downstream observer-shell residual: phase-shifted readout, shell dressing, expansion dressing, or another post-transport effect living after the exact admissible transport stage.
\end{enumerate}
This is the strongest structural narrowing achieved so far: primitive matrix leakage and canonical observer-shell residual are no longer being mixed. The architecture itself says that only one of them may remain once admissible grouped transport is enforced.
\end{remark}


\subsection*{Canonical downstream observer-shell residual after exact admissible transport}
The previous dichotomy showed that once grouped no-drift holds on the active leakage plane and admissible grouped Millennium transport is exact on the active branch, no primitive HC/MM/J leakage can remain visible. Any still visible OP~1/OP~5 remainder must therefore live strictly downstream from the exact transport stage. The next structural step is to package that remainder as a single canonical observer-shell mode rather than as two independent unfinished closures.

\begin{definition}[Working downstream observer-shell dressing mode]
Assume exact grouped no-drift on the active leakage plane $W$ and exact admissible grouped Millennium transport on the active branch $U$. A downstream observer-shell dressing mode is a scalar family
\[
  r_{\mathrm{obs}}(\theta)
\]
and a common visible shell profile
\[
  e_{\mathrm{obs}}:=\binom{1}{1}
\]
such that the residual visible channel can be written in the form
\[
  d(\theta)=r_{\mathrm{obs}}(\theta)\,e_{\mathrm{obs}}+d_{\mathrm{split}}(\theta),
\]
where $d_{\mathrm{split}}(\theta)$ is a pure observer-side split term living after the exact transport stage.
\end{definition}

\begin{theorem}[Working canonical observer-shell residual after exact admissible transport]
Assume exact grouped no-drift on the active leakage plane $W$ and exact admissible grouped Millennium transport on the active branch $U$. Then any remaining visible OP~1/OP~5 residual admits the canonical decomposition
\[
  d(\theta)=r_{\mathrm{obs}}(\theta)\binom{1}{1}+c_{\mathrm{obs}}(\theta)\binom{1}{-1}.
\]
Moreover, the primitive matrix-leakage contribution vanishes identically, so the common-mode coefficient $r_{\mathrm{obs}}(\theta)$ is a purely downstream observer-shell residual rather than a primitive HC/MM/J leakage carrier.
\end{theorem}

\begin{proof}
By the shared-residual dichotomy, the primitive visible channel decomposes as a common-mode part and an observer-split part,
\[
  d(\theta)=c_{\mathrm{sh}}(\theta)\binom{1}{1}+c_{\mathrm{obs}}(\theta)\binom{1}{-1}.
\]
The previous active-matrix-fit analysis identifies the primitive common-mode part with the active HC/MM/J fit carrier. Under exact admissible grouped Millennium transport on $U$, the active commutator vanishes and therefore the primitive common-mode carrier vanishes as well. Hence any surviving common-mode coefficient can only appear after the exact transport stage, i.e. as downstream observer-shell dressing. Renaming this remaining downstream common-mode coefficient by $r_{\mathrm{obs}}(\theta)$ yields the stated decomposition.
\end{proof}

\begin{corollary}[Working no-drift collapses to one observer-shell mode]
Under the hypotheses of the theorem, if the observer-side split also vanishes on the active leakage plane,
\[
  c_{\mathrm{obs}}(\theta)=0\qquad\text{for all }\theta,
\]
then the full visible OP~1/OP~5 remainder collapses to a single canonical observer-shell residual:
\[
  d(\theta)=r_{\mathrm{obs}}(\theta)\binom{1}{1}.
\]
In particular, OP~1 and OP~5 then do not carry two independent unfinished closures on the active branch; they are merely two reads of the same downstream shell residual.
\end{corollary}

\begin{proof}
The theorem already gives the decomposition
\[
  d(\theta)=r_{\mathrm{obs}}(\theta)\binom{1}{1}+c_{\mathrm{obs}}(\theta)\binom{1}{-1}.
\]
If $c_{\mathrm{obs}}(\theta)=0$, only the common shell term remains, which is exactly the stated formula.
\end{proof}

\begin{remark}[Working interpretation of the surviving residual]
This is the intended structural endpoint of the current OP~1/OP~5 reduction. Once exact admissible transport kills the primitive HC/MM/J leakage and grouped no-drift removes the active observer split, any remaining visible residue should not be interpreted as evidence for two separate unfinished closures. It is a single canonical downstream shell residual, shared by OP~1 and OP~5 and only seen through two observer labels.
\end{remark}


\subsection*{From two unfinished closures to one shell-mode law}
The previous reduction reached a decisive structural endpoint: once exact admissible transport removes primitive HC/MM/J leakage and grouped no-drift removes the active observer split, OP~1 and OP~5 no longer describe two independent unfinished closures on the active branch. What remains is a single downstream observer-shell mode. The next step is therefore not to force two separate zero-conditions, but to state explicitly that the visible target has collapsed to one scalar shell law.

\begin{definition}[Working shared shell-mode amplitude]
Assume the hypotheses of the previous corollary, so that
\[
  d(\theta)=r_{\mathrm{obs}}(\theta)\binom{1}{1}.
\]
The scalar family
\[
  r_{\mathrm{obs}}(\theta)
\]
is called the shared shell-mode amplitude of the visible OP~1/OP~5 residual.
\end{definition}

\begin{theorem}[Working one-law reduction of OP~1 and OP~5 after exact transport]
Assume exact grouped no-drift on the active leakage plane $W$, exact admissible grouped Millennium transport on the active branch $U$, and vanishing active observer split on $W$. Then the visible OP~1/OP~5 channel is governed by a single scalar law:
\[
  d_1(\theta)=d_5(\theta)=r_{\mathrm{obs}}(\theta).
\]
In particular, the visible target is no longer a pair of independent closure equations, but the single shell-mode equation
\[
  r_{\mathrm{obs}}(\theta)=0
\]
if exact visible closure is required, or the single shell-mode identification problem
\[
  r_{\mathrm{obs}}(\theta)=r_{\mathrm{can}}(\theta)
\]
if a canonical nonzero observer-shell residual is allowed.
\end{theorem}

\begin{proof}
Under the stated hypotheses, the previous corollary gives
\[
  d(\theta)=r_{\mathrm{obs}}(\theta)\binom{1}{1}.
\]
Reading off the two components immediately yields
\[
  d_1(\theta)=r_{\mathrm{obs}}(\theta),
  \qquad
  d_5(\theta)=r_{\mathrm{obs}}(\theta).
\]
Hence the visible OP~1/OP~5 target collapses from two residual equations to one scalar shell-mode law. If one demands exact visible closure, that law is $r_{\mathrm{obs}}(\theta)=0$. If one instead admits a canonical downstream observer-shell remainder, the problem is to identify the canonical shell law $r_{\mathrm{can}}(\theta)$ and match the visible residual against it.
\end{proof}

\begin{corollary}[Working exact closure and canonical residual are one-mode alternatives]
Under the same hypotheses, one cannot simultaneously interpret the surviving visible channel as two independent unfinished closures and as a single canonical observer-shell mode. Exactly one of the following descriptions may hold:
\begin{enumerate}[leftmargin=2em,label=(\roman*)]
\item exact visible closure, i.e.
\[
  r_{\mathrm{obs}}(\theta)=0\qquad\text{for all }\theta;
\]
\item canonical observer-shell remainder, i.e.
\[
  r_{\mathrm{obs}}(\theta)=r_{\mathrm{can}}(\theta)\not\equiv 0.
\]
\end{enumerate}
In either case, OP~1 and OP~5 remain locked to the same scalar mode.
\end{corollary}

\begin{proof}
The theorem shows that both visible channels are equal to the same scalar amplitude $r_{\mathrm{obs}}(\theta)$. Therefore any interpretation of the surviving visible remainder must be an interpretation of this single scalar family. It cannot simultaneously serve as two independent residual equations and as one shared shell mode. The two listed alternatives are exactly the two structurally admissible readings.
\end{proof}

\begin{remark}[Working interpretation: closure target versus shell-law target]
This is the point where the present architecture changes the question. The old reading asked whether OP~1 and OP~5 can both be closed independently. The new reading says: after exact admissible transport and active no-drift, there is only one visible target left. The remaining decision is whether that target should vanish or whether it should be identified with a canonical post-transport shell law, for example expansion dressing or phase-shifted observer-shell readout.
\end{remark}


\subsection*{Canonical phase-shifted / expansion-dressed shell law}
The previous reduction showed that after exact admissible transport and active no-drift, the visible OP~1/OP~5 question collapses to a single scalar shell-mode amplitude
\[
  r_{\mathrm{obs}}(\theta).
\]
The next structural step is to stop treating this quantity as a bare leftover and instead give a first explicit law for it. The guiding hypothesis is that the observer shell is not a second ontology but a small post-transport dressing of the same exact ground closure. In that case the surviving visible mode should arise from a phase-shifted and possibly expansion-dressed readout of the same underlying shell carrier.

\begin{definition}[Working observer-shell phase and expansion defects]
A post-transport observer-shell defect state is described by two scalar families
\[
  \delta_{\mathrm{obs}}(\theta),
  \qquad
  \eta_{\mathrm{exp}}(\theta),
\]
where
\begin{itemize}[leftmargin=2em]
\item $\delta_{\mathrm{obs}}(\theta)$ measures the local phase shift between the exact transport shell and the visible observer shell,
\item $\eta_{\mathrm{exp}}(\theta)$ measures the local expansion / shell-spacing drift of the visible shell relative to the exact transport shell.
\end{itemize}
Both quantities are downstream observer-shell parameters: they vanish on the exact transport shell itself.
\end{definition}

\begin{definition}[Working first-order canonical shell law]
A first-order canonical observer-shell law is an identity of the form
\[
  r_{\mathrm{can}}(\theta)
  =A_{\mathrm{sh}}(\theta)\,\delta_{\mathrm{obs}}(\theta)
  +B_{\mathrm{sh}}(\theta)\,\eta_{\mathrm{exp}}(\theta)
  +R_{\ge 2}(\theta),
\]
where $A_{\mathrm{sh}}(\theta)$ and $B_{\mathrm{sh}}(\theta)$ are shell-response coefficients and the remainder satisfies
\[
  |R_{\ge 2}(\theta)|
  \le C_{\mathrm{sh}}(\theta)
  \bigl(|\delta_{\mathrm{obs}}(\theta)|+|\eta_{\mathrm{exp}}(\theta)|\bigr)^2.
\]
In the strictly phase-shifted case one has $B_{\mathrm{sh}}\equiv 0$, while in the pure expansion-dressing case one has $A_{\mathrm{sh}}\equiv 0$.
\end{definition}

\begin{theorem}[Working first-order phase-shift / expansion law for the shared shell mode]
Assume exact admissible grouped Millennium transport on the active branch, exact grouped no-drift on the active leakage plane, and vanishing active observer split on that plane. Assume moreover that the visible observer shell is a smooth downstream perturbation of the exact transport shell, parameterized by phase defect $\delta_{\mathrm{obs}}(\theta)$ and expansion defect $\eta_{\mathrm{exp}}(\theta)$. Then the shared visible mode admits a first-order shell law
\[
  r_{\mathrm{obs}}(\theta)
  =A_{\mathrm{sh}}(\theta)\,\delta_{\mathrm{obs}}(\theta)
  +B_{\mathrm{sh}}(\theta)\,\eta_{\mathrm{exp}}(\theta)
  +R_{\ge 2}(\theta),
\]
with a quadratic remainder in the total shell defect magnitude. In particular, exact visible closure is equivalent to the vanishing of this canonical shell law, while a nonzero visible remainder is compatible with exact transport provided it is entirely generated by downstream shell dressing.
\end{theorem}

\begin{proof}
Under the stated hypotheses the visible channel has already collapsed to a single scalar amplitude,
\[
  d(\theta)=r_{\mathrm{obs}}(\theta)\binom{1}{1}.
\]
Hence every surviving visible effect must be read as a downstream perturbation of the same exact transport shell rather than as a primitive matrix-leakage channel. By the assumed smooth dependence of the visible observer shell on the two downstream shell parameters $\delta_{\mathrm{obs}}(\theta)$ and $\eta_{\mathrm{exp}}(\theta)$, the scalar amplitude $r_{\mathrm{obs}}(\theta)$ is a smooth scalar function of these two arguments. A first-order Taylor expansion around the exact shell point
\[
  (\delta_{\mathrm{obs}},\eta_{\mathrm{exp}})=(0,0)
\]
yields
\[
  r_{\mathrm{obs}}(\theta)
  =A_{\mathrm{sh}}(\theta)\,\delta_{\mathrm{obs}}(\theta)
  +B_{\mathrm{sh}}(\theta)\,\eta_{\mathrm{exp}}(\theta)
  +R_{\ge 2}(\theta),
\]
where the remainder is quadratic in the total shell defect magnitude. This is exactly the stated first-order canonical shell law.
\end{proof}

\begin{corollary}[Working exact closure versus canonical nonzero shell law]
Under the same hypotheses, the remaining visible OP~1/OP~5 target is reduced to the dichotomy
\begin{enumerate}[leftmargin=2em,label=(\roman*)]
\item exact visible closure:
\[
  A_{\mathrm{sh}}(\theta)\,\delta_{\mathrm{obs}}(\theta)
  +B_{\mathrm{sh}}(\theta)\,\eta_{\mathrm{exp}}(\theta)
  +R_{\ge 2}(\theta)=0;
\]
\item canonical nonzero observer-shell remainder:
\[
  r_{\mathrm{obs}}(\theta)=r_{\mathrm{can}}(\theta)
  \not\equiv 0
\]
with $r_{\mathrm{can}}$ given by the same first-order shell law.
\end{enumerate}
In both cases the visible question is one scalar shell law rather than two independent OP~1 / OP~5 closures.
\end{corollary}

\begin{proof}
The previous theorem identifies the surviving scalar mode with a first-order observer-shell law in the two downstream shell parameters. Therefore the visible target is exhausted by deciding whether that law vanishes identically or whether it should be retained as a canonical nonzero shell remainder.
\end{proof}

\begin{remark}[Working interpretation of the shell-law coefficients]
The coefficients $A_{\mathrm{sh}}(\theta)$ and $B_{\mathrm{sh}}(\theta)$ are the first visible response coefficients of the exact transport shell to observer-side phase shift and observer-side expansion dressing. They are not new primitive leakage parameters. Their role is to measure how a single exact transport closure is seen from a displaced observer shell. This is why OP~1 and OP~5 remain locked to the same scalar law even after exact transport has already removed the primitive HC/MM/J leakage.
\end{remark}


\subsection*{Phase-dominant, expansion-dominant,\ and mixed shell-law regimes}
The canonical shell law now reduces the visible post-transport question to one scalar combination of two downstream shell effects. The next structural step is therefore not to search for two further closures, but to decide which of the two first-order shell loads is actually primary.

\begin{definition}[Working shell loads and dominance ratio]
Define the first-order shell loads
\[
  \Lambda_{\mathrm{ph}}(\theta):=|A_{\mathrm{sh}}(\theta)\,\delta_{\mathrm{obs}}(\theta)|,
  \qquad
  \Lambda_{\mathrm{exp}}(\theta):=|B_{\mathrm{sh}}(\theta)\,\eta_{\mathrm{exp}}(\theta)|,
\]
and the total first-order shell load
\[
  \Lambda_{\mathrm{lin}}(\theta):=\Lambda_{\mathrm{ph}}(\theta)+\Lambda_{\mathrm{exp}}(\theta).
\]
Whenever $\Lambda_{\mathrm{ph}}(\theta)\neq 0$, define the shell-load ratio
\[
  \rho_{\mathrm{sh}}(\theta):=\frac{\Lambda_{\mathrm{exp}}(\theta)}{\Lambda_{\mathrm{ph}}(\theta)}.
\]
Thus $\rho_{\mathrm{sh}}\ll 1$ indicates phase-dominant shell dressing, while $\rho_{\mathrm{sh}}\gg 1$ indicates expansion-dominant shell dressing.
\end{definition}

\begin{definition}[Working phase-dominant, expansion-dominant, and mixed regimes]
Fix $\varepsilon_{\mathrm{dom}},\varepsilon_{\mathrm{rem}}\in[0,1)$ with
\[
  \varepsilon_{\mathrm{dom}}+\varepsilon_{\mathrm{rem}}<1.
\]
At a given parameter value $\theta$:
\begin{enumerate}[leftmargin=2em,label=(\roman*)]
\item the shell law is called \emph{phase-dominant} if
\[
  \Lambda_{\mathrm{exp}}(\theta)\le \varepsilon_{\mathrm{dom}}\,\Lambda_{\mathrm{ph}}(\theta)
  \quad\text{and}\quad
  |R_{\ge 2}(\theta)|\le \varepsilon_{\mathrm{rem}}\,\Lambda_{\mathrm{ph}}(\theta);
\]
\item it is called \emph{expansion-dominant} if
\[
  \Lambda_{\mathrm{ph}}(\theta)\le \varepsilon_{\mathrm{dom}}\,\Lambda_{\mathrm{exp}}(\theta)
  \quad\text{and}\quad
  |R_{\ge 2}(\theta)|\le \varepsilon_{\mathrm{rem}}\,\Lambda_{\mathrm{exp}}(\theta);
\]
\item it is called \emph{genuinely mixed at first order} if neither of the previous two dominance conditions holds.
\end{enumerate}
\end{definition}

\begin{theorem}[Working shell-law dominance classification]
Assume the first-order canonical shell law of the previous subsection. Then:
\begin{enumerate}[leftmargin=2em,label=(\roman*)]
\item in the phase-dominant regime one has
\[
  |r_{\mathrm{obs}}(\theta)-A_{\mathrm{sh}}(\theta)\,\delta_{\mathrm{obs}}(\theta)|
  \le (\varepsilon_{\mathrm{dom}}+\varepsilon_{\mathrm{rem}})\,\Lambda_{\mathrm{ph}}(\theta),
\]
so the shared shell mode is controlled by the observer-side phase-shift term up to a strictly smaller relative error;
\item in the expansion-dominant regime one has
\[
  |r_{\mathrm{obs}}(\theta)-B_{\mathrm{sh}}(\theta)\,\eta_{\mathrm{exp}}(\theta)|
  \le (\varepsilon_{\mathrm{dom}}+\varepsilon_{\mathrm{rem}})\,\Lambda_{\mathrm{exp}}(\theta),
\]
so the shared shell mode is controlled by the observer-shell expansion term up to a strictly smaller relative error;
\item in the genuinely mixed regime, the visible shell mode cannot be reduced at first order to only one of the two shell defects.
\end{enumerate}
\end{theorem}

\begin{proof}
By the first-order shell law,
\[
  r_{\mathrm{obs}}(\theta)
  =A_{\mathrm{sh}}(\theta)\,\delta_{\mathrm{obs}}(\theta)
  +B_{\mathrm{sh}}(\theta)\,\eta_{\mathrm{exp}}(\theta)
  +R_{\ge 2}(\theta).
\]
Hence
\[
  |r_{\mathrm{obs}}(\theta)-A_{\mathrm{sh}}(\theta)\,\delta_{\mathrm{obs}}(\theta)|
  \le \Lambda_{\mathrm{exp}}(\theta)+|R_{\ge 2}(\theta)|,
\]
and similarly
\[
  |r_{\mathrm{obs}}(\theta)-B_{\mathrm{sh}}(\theta)\,\eta_{\mathrm{exp}}(\theta)|
  \le \Lambda_{\mathrm{ph}}(\theta)+|R_{\ge 2}(\theta)|.
\]
In the phase-dominant regime this gives
\[
  |r_{\mathrm{obs}}(\theta)-A_{\mathrm{sh}}(\theta)\,\delta_{\mathrm{obs}}(\theta)|
  \le (\varepsilon_{\mathrm{dom}}+\varepsilon_{\mathrm{rem}})\,\Lambda_{\mathrm{ph}}(\theta),
\]
which is strictly smaller than the dominant phase load because $\varepsilon_{\mathrm{dom}}+\varepsilon_{\mathrm{rem}}<1$. The expansion-dominant case is identical after swapping the two shell loads. If neither dominance condition holds, then neither first-order term controls the shell mode up to a strictly smaller remainder, so the shell mode is genuinely mixed at first order.
\end{proof}

\begin{corollary}[Working next decision for the shared shell mode]
After exact admissible transport, active no-drift, and collapse to a single visible shell mode, the next structural decision is reduced to the comparison of the two first-order shell loads
\[
  \Lambda_{\mathrm{ph}}(\theta)
  \quad\text{and}\quad
  \Lambda_{\mathrm{exp}}(\theta),
\]
not to two independent closure equations. In particular, the visible question is now whether the common shell mode is primarily
\begin{enumerate}[leftmargin=2em,label=(\roman*)]
\item phase-shift driven,
\item expansion-drift driven,
\item or genuinely mixed at first order.
\end{enumerate}
\end{corollary}

\begin{proof}
This is the direct content of the dominance classification theorem.
\end{proof}

\begin{remark}[Working interpretation of the present status]
The current document does not yet prove globally which regime is realized. What it does prove is that the decision has been compressed to a single shell-law comparison problem. The remaining visible burden is therefore much narrower than the earlier picture of two unfinished OP~1 / OP~5 closures: it is now a one-mode question about which downstream observer-shell load is structurally primary.
\end{remark}



\subsection*{Signed shell-balance mismatch and the cancellation corridor}
The shell-load classification identifies which first-order observer-shell contribution is dominant, but exact visible closure depends on something sharper: whether the two signed first-order shell contributions actually cancel. The next structural step is therefore to replace the unsigned load comparison by a signed shell-balance mismatch.

\begin{definition}[Working signed shell coordinates and shell-balance mismatch]
Define the signed first-order shell coordinates
\[
  p_{\mathrm{sh}}(\theta):=A_{\mathrm{sh}}(\theta)\,\delta_{\mathrm{obs}}(\theta),
  \qquad
  e_{\mathrm{sh}}(\theta):=B_{\mathrm{sh}}(\theta)\,\eta_{\mathrm{exp}}(\theta),
\]
and the signed shell-balance mismatch
\[
  \chi_{\mathrm{sh}}(\theta):=p_{\mathrm{sh}}(\theta)+e_{\mathrm{sh}}(\theta).
\]
Whenever $p_{\mathrm{sh}}(\theta)\neq 0$, define the signed shell-balance ratio
\[
  \zeta_{\mathrm{sh}}(\theta):=\frac{e_{\mathrm{sh}}(\theta)}{p_{\mathrm{sh}}(\theta)}.
\]
Thus exact first-order cancellation corresponds to $\chi_{\mathrm{sh}}(\theta)=0$, equivalently $\zeta_{\mathrm{sh}}(\theta)=-1$.
\end{definition}

\begin{definition}[Working shell-cancellation corridor]
Fix $\varepsilon_{\mathrm{bal}},\varepsilon_{\mathrm{rem}}\in[0,1)$ with
\[
  \varepsilon_{\mathrm{bal}}+\varepsilon_{\mathrm{rem}}<1.
\]
At a given parameter value $\theta$, the shell law is said to lie in the \emph{cancellation corridor} if
\[
  |\chi_{\mathrm{sh}}(\theta)|\le \varepsilon_{\mathrm{bal}}\,\Lambda_{\mathrm{lin}}(\theta)
  \qquad\text{and}\qquad
  |R_{\ge 2}(\theta)|\le \varepsilon_{\mathrm{rem}}\,\Lambda_{\mathrm{lin}}(\theta).
\]
Equivalently, the first-order signed mismatch and the higher-order shell remainder are both smaller than the total first-order shell load.
\end{definition}

\begin{theorem}[Working shell-balance reduction and exact cancellation corridor]
Assume the first-order canonical shell law. Then:
\begin{enumerate}[leftmargin=2em,label=(\roman*)]
\item one has the exact reduction
\[
  r_{\mathrm{obs}}(\theta)=\chi_{\mathrm{sh}}(\theta)+R_{\ge 2}(\theta);
\]
\item in the cancellation corridor one has
\[
  |r_{\mathrm{obs}}(\theta)|
  \le (\varepsilon_{\mathrm{bal}}+\varepsilon_{\mathrm{rem}})\,\Lambda_{\mathrm{lin}}(\theta),
\]
so exact visible closure is reduced to a small signed shell-balance mismatch rather than to two separate closures;
\item if $p_{\mathrm{sh}}(\theta)$ and $e_{\mathrm{sh}}(\theta)$ have the same sign, then
\[
  |\chi_{\mathrm{sh}}(\theta)|=|p_{\mathrm{sh}}(\theta)|+|e_{\mathrm{sh}}(\theta)|=\Lambda_{\mathrm{lin}}(\theta),
\]
so exact visible closure at first order is impossible unless the total first-order shell load itself vanishes.
\end{enumerate}
\end{theorem}

\begin{proof}
By definition of the signed shell coordinates,
\[
  p_{\mathrm{sh}}(\theta)+e_{\mathrm{sh}}(\theta)=A_{\mathrm{sh}}(\theta)\,\delta_{\mathrm{obs}}(\theta)+B_{\mathrm{sh}}(\theta)\,\eta_{\mathrm{exp}}(\theta)=\chi_{\mathrm{sh}}(\theta),
\]
so the first-order shell law becomes
\[
  r_{\mathrm{obs}}(\theta)=\chi_{\mathrm{sh}}(\theta)+R_{\ge 2}(\theta).
\]
This proves (i). If the shell law lies in the cancellation corridor, then
\[
  |r_{\mathrm{obs}}(\theta)|\le |\chi_{\mathrm{sh}}(\theta)|+|R_{\ge 2}(\theta)|
  \le (\varepsilon_{\mathrm{bal}}+\varepsilon_{\mathrm{rem}})\,\Lambda_{\mathrm{lin}}(\theta),
\]
which proves (ii). If $p_{\mathrm{sh}}(\theta)$ and $e_{\mathrm{sh}}(\theta)$ have the same sign, then they add without cancellation, hence
\[
  |\chi_{\mathrm{sh}}(\theta)|=|p_{\mathrm{sh}}(\theta)+e_{\mathrm{sh}}(\theta)|=|p_{\mathrm{sh}}(\theta)|+|e_{\mathrm{sh}}(\theta)|=\Lambda_{\mathrm{lin}}(\theta),
\]
which proves (iii).
\end{proof}

\begin{corollary}[Working next decision after shell-load dominance]
After the shell-load dominance classification, the remaining visible decision is whether the signed ratio $\zeta_{\mathrm{sh}}(\theta)$ enters a neighborhood of $-1$. In particular:
\begin{enumerate}[leftmargin=2em,label=(\roman*)]
\item if $\zeta_{\mathrm{sh}}(\theta)\approx -1$, the visible shell mode lies near the cancellation corridor and exact visible closure remains possible;
\item if $\zeta_{\mathrm{sh}}(\theta)>0$, then the first-order shell law is sign-coherent and forces a canonical nonzero observer-shell remainder unless the total first-order shell load itself vanishes;
\item if $\zeta_{\mathrm{sh}}(\theta)$ stays bounded away from $-1$, then the shared shell mode is canonically nonzero already at first order.
\end{enumerate}
\end{corollary}

\begin{proof}
The previous theorem shows that the visible shell mode is controlled by the signed mismatch $\chi_{\mathrm{sh}}(\theta)=p_{\mathrm{sh}}(\theta)(1+\zeta_{\mathrm{sh}}(\theta))$ together with the higher-order remainder. Thus entry into the cancellation corridor is equivalent to placing the signed ratio near $-1$, while positive ratio excludes first-order cancellation.
\end{proof}

\begin{remark}[Working interpretation of the present shell-law reduction]
This is the first place where the two-shell-defect picture collapses all the way to a single signed comparison. The visible question is no longer which of two unfinished closures dominates, but whether the phase-shift and expansion-drift contributions cancel each other or reinforce each other on the observer shell. If they reinforce, the visible residual is canonically nonzero; if they nearly cancel, exact visible closure remains structurally possible.
\end{remark}


\subsection*{Normalized shell-cancellation defect and the exact first-order verdict}
The signed shell-balance mismatch already reduces the visible problem to a single first-order comparison. The next structural compression is to normalize that mismatch by the total first-order shell load itself. This removes the overall shell scale and leaves only the genuinely visible cancellation defect.

\begin{definition}[Working normalized shell defect and cancellation depth]
Whenever $\Lambda_{\mathrm{lin}}(\theta)>0$, define the normalized shell-cancellation defect by
\[
  \rho_{\mathrm{sh}}(\theta):=\frac{\chi_{\mathrm{sh}}(\theta)}{\Lambda_{\mathrm{lin}}(\theta)}
  =\frac{p_{\mathrm{sh}}(\theta)+e_{\mathrm{sh}}(\theta)}{|p_{\mathrm{sh}}(\theta)|+|e_{\mathrm{sh}}(\theta)|}.
\]
Set $\rho_{\mathrm{sh}}(\theta):=0$ when $\Lambda_{\mathrm{lin}}(\theta)=0$. Define the cancellation depth by
\[
  \kappa_{\mathrm{sh}}(\theta):=1-|\rho_{\mathrm{sh}}(\theta)|.
\]
Thus $|\rho_{\mathrm{sh}}|$ measures the normalized first-order noncancellation, while $\kappa_{\mathrm{sh}}$ measures how deeply the shell law enters the first-order cancellation corridor.
\end{definition}

\begin{theorem}[Working normalized shell defect, cancellation depth, and exact first-order criterion]
For every parameter value $\theta$ one has:
\begin{enumerate}[leftmargin=2em,label=(\roman*)]
\item
\[
  |\rho_{\mathrm{sh}}(\theta)|\le 1,
  \qquad
  0\le \kappa_{\mathrm{sh}}(\theta)\le 1;
\]
\item if $p_{\mathrm{sh}}(\theta)$ and $e_{\mathrm{sh}}(\theta)$ have the same sign or one of them vanishes, then
\[
  |\rho_{\mathrm{sh}}(\theta)|=1,
  \qquad
  \kappa_{\mathrm{sh}}(\theta)=0;
\]
\item if $p_{\mathrm{sh}}(\theta)$ and $e_{\mathrm{sh}}(\theta)$ have opposite sign, then
\[
  |\rho_{\mathrm{sh}}(\theta)|
  =\frac{\bigl||p_{\mathrm{sh}}(\theta)|-|e_{\mathrm{sh}}(\theta)|\bigr|}{|p_{\mathrm{sh}}(\theta)|+|e_{\mathrm{sh}}(\theta)|},
\]
\[
  \kappa_{\mathrm{sh}}(\theta)
  =\frac{2\min\{|p_{\mathrm{sh}}(\theta)|,|e_{\mathrm{sh}}(\theta)|\}}{|p_{\mathrm{sh}}(\theta)|+|e_{\mathrm{sh}}(\theta)|};
\]
\item exact first-order shell cancellation holds if and only if
\[
  \rho_{\mathrm{sh}}(\theta)=0,
\]
which is equivalent to opposite sign and equal first-order shell magnitudes,
\[
  |p_{\mathrm{sh}}(\theta)|=|e_{\mathrm{sh}}(\theta)|.
\]
\end{enumerate}
\end{theorem}

\begin{proof}
By definition,
\[
  |\chi_{\mathrm{sh}}(\theta)|=|p_{\mathrm{sh}}(\theta)+e_{\mathrm{sh}}(\theta)|
  \le |p_{\mathrm{sh}}(\theta)|+|e_{\mathrm{sh}}(\theta)|
  =\Lambda_{\mathrm{lin}}(\theta),
\]
so $|\rho_{\mathrm{sh}}(\theta)|\le 1$, and the cancellation depth automatically lies in $[0,1]$. This proves (i). If $p_{\mathrm{sh}}$ and $e_{\mathrm{sh}}$ have the same sign, or one vanishes, then
\[
  |p_{\mathrm{sh}}+e_{\mathrm{sh}}|=|p_{\mathrm{sh}}|+|e_{\mathrm{sh}}|,
\]
so $|\rho_{\mathrm{sh}}|=1$ and therefore $\kappa_{\mathrm{sh}}=0$, proving (ii). If the signs are opposite, then
\[
  |p_{\mathrm{sh}}+e_{\mathrm{sh}}|=\bigl||p_{\mathrm{sh}}|-|e_{\mathrm{sh}}|\bigr|,
\]
which yields the displayed formula for $|\rho_{\mathrm{sh}}|$ after division by $|p_{\mathrm{sh}}|+|e_{\mathrm{sh}}|$. The formula for $\kappa_{\mathrm{sh}}$ follows because
\[
  1-\frac{\bigl||p|-|e|\bigr|}{|p|+|e|}=\frac{|p|+|e|-\bigl||p|-|e|\bigr|}{|p|+|e|}=\frac{2\min\{|p|,|e|\}}{|p|+|e|}.
\]
This proves (iii). Finally, $\rho_{\mathrm{sh}}(\theta)=0$ is equivalent to $\chi_{\mathrm{sh}}(\theta)=0$, which is exactly first-order shell cancellation. Under opposite sign this is equivalent to equal magnitudes, proving (iv).
\end{proof}

\begin{definition}[Working normalized higher-order shell remainder]
Whenever $\Lambda_{\mathrm{lin}}(\theta)>0$, define the normalized higher-order shell remainder by
\[
  \rho_{\mathrm{rem}}(\theta):=\frac{|R_{\ge 2}(\theta)|}{\Lambda_{\mathrm{lin}}(\theta)}.
\]
Set $\rho_{\mathrm{rem}}(\theta):=0$ when $\Lambda_{\mathrm{lin}}(\theta)=0$.
\end{definition}

\begin{corollary}[Working exact visible closure requires normalized first-order cancellation]
Assume the first-order canonical shell law and let $\Lambda_{\mathrm{lin}}(\theta)>0$. Then:
\begin{enumerate}[leftmargin=2em,label=(\roman*)]
\item
\[
  \frac{|r_{\mathrm{obs}}(\theta)|}{\Lambda_{\mathrm{lin}}(\theta)}
  \le |\rho_{\mathrm{sh}}(\theta)|+\rho_{\mathrm{rem}}(\theta);
\]
\item a necessary condition for exact visible closure at $\theta$ is
\[
  |\rho_{\mathrm{sh}}(\theta)|\le \rho_{\mathrm{rem}}(\theta);
\]
\item if
\[
  |\rho_{\mathrm{sh}}(\theta)|>\rho_{\mathrm{rem}}(\theta),
\]
then the shared visible shell mode is canonically nonzero at $\theta$.
\end{enumerate}
\end{corollary}

\begin{proof}
By the shell-balance reduction theorem,
\[
  r_{\mathrm{obs}}(\theta)=\chi_{\mathrm{sh}}(\theta)+R_{\ge 2}(\theta),
\]
so after division by $\Lambda_{\mathrm{lin}}(\theta)$ one gets
\[
  \frac{|r_{\mathrm{obs}}(\theta)|}{\Lambda_{\mathrm{lin}}(\theta)}
  \le \frac{|\chi_{\mathrm{sh}}(\theta)|}{\Lambda_{\mathrm{lin}}(\theta)}+\frac{|R_{\ge 2}(\theta)|}{\Lambda_{\mathrm{lin}}(\theta)}
  =|\rho_{\mathrm{sh}}(\theta)|+\rho_{\mathrm{rem}}(\theta),
\]
which proves (i). If exact visible closure holds, then
\[
  0=r_{\mathrm{obs}}(\theta)=\chi_{\mathrm{sh}}(\theta)+R_{\ge 2}(\theta),
\]
so $|\chi_{\mathrm{sh}}(\theta)|=|R_{\ge 2}(\theta)|$. Division by $\Lambda_{\mathrm{lin}}(\theta)$ gives (ii). Statement (iii) is the contrapositive of (ii).
\end{proof}

\begin{remark}[Working interpretation of the normalized shell verdict]
The visible shell problem has now collapsed to a single normalized comparison. The first-order part is measured by $|\rho_{\mathrm{sh}}|$, the higher-order rescue budget by $\rho_{\mathrm{rem}}$, and exact visible closure is possible only if the normalized first-order noncancellation does not exceed the normalized higher-order shell remainder. In particular, once $|\rho_{\mathrm{sh}}|$ is visibly larger than $\rho_{\mathrm{rem}}$, the shared shell mode is forced to remain canonically nonzero.
\end{remark}


\subsection*{Sign-coherent obstruction and the anti-phase closure corridor}
The normalized shell verdict can be compressed one step further. The visible question is not merely whether the normalized first-order mismatch is larger than the higher-order rescue budget, but also whether the two first-order shell loads are sign-coherent or anti-phase. In the sign-coherent regime, first-order cancellation is impossible; in the anti-phase regime, exact visible closure can survive only inside a sharp corridor around the ratio value $-1$.

\begin{theorem}[Working sign-coherent obstruction and anti-phase corridor for visible closure]
Assume the first-order canonical shell law and let $\Lambda_{\mathrm{lin}}(\theta)>0$.
\begin{enumerate}[leftmargin=2em,label=(\roman*)]
\item If $p_{\mathrm{sh}}(\theta)$ and $e_{\mathrm{sh}}(\theta)$ have the same sign or one of them vanishes, then
\[
  |\rho_{\mathrm{sh}}(\theta)|=1.
\]
Consequently, if
\[
  \rho_{\mathrm{rem}}(\theta)<1,
\]
then exact visible closure at $\theta$ is impossible and the shared shell mode is canonically nonzero.
\item Suppose $p_{\mathrm{sh}}(\theta)$ and $e_{\mathrm{sh}}(\theta)$ have opposite sign and $p_{\mathrm{sh}}(\theta)\neq 0$. Set
\[
  \zeta_{\mathrm{sh}}(\theta):=\frac{e_{\mathrm{sh}}(\theta)}{p_{\mathrm{sh}}(\theta)}.
\]
Then necessarily $\zeta_{\mathrm{sh}}(\theta)<0$, and one has
\[
  |\rho_{\mathrm{sh}}(\theta)|
  =\frac{|1+\zeta_{\mathrm{sh}}(\theta)|}{1+|\zeta_{\mathrm{sh}}(\theta)|}.
\]
Hence a necessary condition for exact visible closure is
\[
  \frac{|1+\zeta_{\mathrm{sh}}(\theta)|}{1+|\zeta_{\mathrm{sh}}(\theta)|}
  \le
  \rho_{\mathrm{rem}}(\theta),
\]
that is,
\[
  |1+\zeta_{\mathrm{sh}}(\theta)|
  \le
  \rho_{\mathrm{rem}}(\theta)\bigl(1+|\zeta_{\mathrm{sh}}(\theta)|\bigr).
\]
\item In particular, if
\[
  |1+\zeta_{\mathrm{sh}}(\theta)|
  >
  \rho_{\mathrm{rem}}(\theta)\bigl(1+|\zeta_{\mathrm{sh}}(\theta)|\bigr),
\]
then the shared visible shell mode is canonically nonzero at $\theta$.
\end{enumerate}
\end{theorem}

\begin{proof}
Statement (i) is immediate from the normalized shell-defect theorem, which gives $|\rho_{\mathrm{sh}}|=1$ in the sign-coherent regime. The exact visible closure criterion then yields the necessary inequality
\[
  1=|\rho_{\mathrm{sh}}(\theta)|\le \rho_{\mathrm{rem}}(\theta).
\]
Therefore, if $\rho_{\mathrm{rem}}(\theta)<1$, exact visible closure is impossible and the shell mode remains canonically nonzero.

For (ii), opposite sign implies $\zeta_{\mathrm{sh}}(\theta)<0$. Since
\[
  \chi_{\mathrm{sh}}(\theta)=p_{\mathrm{sh}}(\theta)+e_{\mathrm{sh}}(\theta)=p_{\mathrm{sh}}(\theta)(1+\zeta_{\mathrm{sh}}(\theta)),
\]
and
\[
  \Lambda_{\mathrm{lin}}(\theta)=|p_{\mathrm{sh}}(\theta)|+|e_{\mathrm{sh}}(\theta)|=|p_{\mathrm{sh}}(\theta)|\bigl(1+|\zeta_{\mathrm{sh}}(\theta)|\bigr),
\]
one gets
\[
  |\rho_{\mathrm{sh}}(\theta)|
  =\frac{|\chi_{\mathrm{sh}}(\theta)|}{\Lambda_{\mathrm{lin}}(\theta)}
  =\frac{|1+\zeta_{\mathrm{sh}}(\theta)|}{1+|\zeta_{\mathrm{sh}}(\theta)|}.
\]
The displayed corridor condition then follows from the exact visible closure criterion $|\rho_{\mathrm{sh}}|\le \rho_{\mathrm{rem}}$.

Statement (iii) is the contrapositive of that necessary condition.
\end{proof}

\begin{remark}[Working interpretation of the sign/corridor shell verdict]
The visible shell problem has now become extremely sharp. If the phase and expansion shell loads are sign-coherent, then first-order cancellation is impossible and higher order would have to rescue the full shell scale to produce exact visible closure. If they are anti-phase, then exact visible closure can survive only inside a narrow corridor where the signed shell ratio $\zeta_{\mathrm{sh}}$ sits sufficiently close to $-1$. Outside that corridor, the shared shell mode is forced to remain canonically nonzero.
\end{remark}



\subsection*{Single-shell-coordinate reduction and naturality of the anti-phase corridor}
The anti-phase corridor around $\zeta_{\mathrm{sh}}=-1$ is sharp, but it is not yet clear whether that corridor is structurally natural or only a fine-tuned accident. The next structural step is therefore to ask whether the two downstream shell defects are actually independent, or whether they are both induced by one and the same observer-shell displacement. If the latter is true, then the signed shell ratio collapses at first order to one structural gain ratio, and the question ``why near $-1$?'' becomes a question about one induced shell geometry rather than about two unrelated corrections.

\begin{definition}[Working single-shell-coordinate observer law]
A downstream observer shell is said to be controlled by a \emph{single shell coordinate} if there exists a scalar family $\xi_{\mathrm{sh}}(\theta)$ and coefficient families $a_{\mathrm{sh}}(\theta), b_{\mathrm{sh}}(\theta)$ such that
\[
  \delta_{\mathrm{obs}}(\theta)
  =a_{\mathrm{sh}}(\theta)\,\xi_{\mathrm{sh}}(\theta)+r_{\delta}(\theta),
  \qquad
  \eta_{\mathrm{exp}}(\theta)
  =b_{\mathrm{sh}}(\theta)\,\xi_{\mathrm{sh}}(\theta)+r_{\eta}(\theta),
\]
with quadratic remainder bounds
\[
  |r_{\delta}(\theta)|+|r_{\eta}(\theta)|
  \le C_{\xi}(\theta)\,|\xi_{\mathrm{sh}}(\theta)|^2.
\]
Thus both downstream shell defects are first-order images of the same observer-shell displacement variable.
\end{definition}

\begin{definition}[Working shell gains induced by one shell coordinate]
Under the single-shell-coordinate observer law, define the induced first-order shell gains
\[
  \Gamma_{\mathrm{ph}}(\theta):=A_{\mathrm{sh}}(\theta)\,a_{\mathrm{sh}}(\theta),
  \qquad
  \Gamma_{\mathrm{exp}}(\theta):=B_{\mathrm{sh}}(\theta)\,b_{\mathrm{sh}}(\theta),
\]
and, whenever $\Gamma_{\mathrm{ph}}(\theta)\neq 0$, the induced shell gain ratio
\[
  \zeta_{\Gamma}(\theta):=\frac{\Gamma_{\mathrm{exp}}(\theta)}{\Gamma_{\mathrm{ph}}(\theta)}.
\]
\end{definition}

\begin{theorem}[Working first-order shell ratio collapses to an induced gain ratio]
Assume the first-order canonical shell law and the single-shell-coordinate observer law. Then the signed first-order shell coordinates satisfy
\[
  p_{\mathrm{sh}}(\theta)=\Gamma_{\mathrm{ph}}(\theta)\,\xi_{\mathrm{sh}}(\theta)+O\!\left(|\xi_{\mathrm{sh}}(\theta)|^2\right),
\]
\[
  e_{\mathrm{sh}}(\theta)=\Gamma_{\mathrm{exp}}(\theta)\,\xi_{\mathrm{sh}}(\theta)+O\!\left(|\xi_{\mathrm{sh}}(\theta)|^2\right),
\]
and hence
\[
  \chi_{\mathrm{sh}}(\theta)
  =\bigl(\Gamma_{\mathrm{ph}}(\theta)+\Gamma_{\mathrm{exp}}(\theta)\bigr)\,\xi_{\mathrm{sh}}(\theta)
  +O\!\left(|\xi_{\mathrm{sh}}(\theta)|^2\right).
\]
If moreover $\Gamma_{\mathrm{ph}}(\theta)\neq 0$ and $\xi_{\mathrm{sh}}(\theta)\neq 0$, then
\[
  \zeta_{\mathrm{sh}}(\theta)
  =\zeta_{\Gamma}(\theta)+O\!\left(|\xi_{\mathrm{sh}}(\theta)|\right).
\]
In particular, the anti-phase corridor is first-order natural exactly when the induced gain ratio satisfies
\[
  \zeta_{\Gamma}(\theta)\approx -1.
\]
\end{theorem}

\begin{proof}
Insert the single-shell-coordinate expansions into the definitions
\[
  p_{\mathrm{sh}}(\theta)=A_{\mathrm{sh}}(\theta)\,\delta_{\mathrm{obs}}(\theta),
  \qquad
  e_{\mathrm{sh}}(\theta)=B_{\mathrm{sh}}(\theta)\,\eta_{\mathrm{exp}}(\theta).
\]
This gives
\[
  p_{\mathrm{sh}}(\theta)
  =A_{\mathrm{sh}}(\theta)a_{\mathrm{sh}}(\theta)\,\xi_{\mathrm{sh}}(\theta)
   +A_{\mathrm{sh}}(\theta)r_{\delta}(\theta)
  =\Gamma_{\mathrm{ph}}(\theta)\,\xi_{\mathrm{sh}}(\theta)+O\!\left(|\xi_{\mathrm{sh}}(\theta)|^2\right),
\]
with the analogous formula for $e_{\mathrm{sh}}(\theta)$. Summing yields the formula for $\chi_{\mathrm{sh}}(\theta)$. If $\Gamma_{\mathrm{ph}}(\theta)\neq 0$ and $\xi_{\mathrm{sh}}(\theta)\neq 0$, then division of the two first-order expansions gives
\[
  \zeta_{\mathrm{sh}}(\theta)
  =\frac{\Gamma_{\mathrm{exp}}(\theta)+O(|\xi_{\mathrm{sh}}(\theta)|)}{\Gamma_{\mathrm{ph}}(\theta)+O(|\xi_{\mathrm{sh}}(\theta)|)}
  =\frac{\Gamma_{\mathrm{exp}}(\theta)}{\Gamma_{\mathrm{ph}}(\theta)}+O\!\left(|\xi_{\mathrm{sh}}(\theta)|\right),
\]
which is exactly the stated collapse to the induced gain ratio.
\end{proof}

\begin{corollary}[Working structural meaning of the anti-phase corridor]
Under the same hypotheses, the anti-phase corridor around $\zeta_{\mathrm{sh}}(\theta)=-1$ is not a separate tuning problem between two independent shell defects. It is, to first order, the statement that the single observer-shell displacement variable $\xi_{\mathrm{sh}}(\theta)$ is seen by the phase and expansion readouts with opposite and nearly equal induced gains:
\[
  \Gamma_{\mathrm{exp}}(\theta)\approx -\Gamma_{\mathrm{ph}}(\theta).
\]
Thus the visible shell question reduces further to deciding whether the induced gain pair is
\begin{enumerate}[leftmargin=2em,label=(\roman*)]
\item sign-coherent,
\item anti-phase but unbalanced,
\item or anti-phase and nearly balanced.
\end{enumerate}
\end{corollary}

\begin{proof}
This is exactly the content of the theorem after replacing $\zeta_{\mathrm{sh}}(\theta)$ by its first-order approximation $\zeta_{\Gamma}(\theta)$.
\end{proof}

\begin{remark}[Working next decision after the shell-ratio corridor]
The visible shell problem has now moved from two apparent shell defects to one observer-shell displacement and one induced gain pair. This is a genuine structural compression. It means that the relevant question is no longer whether two unrelated corrections accidentally cancel, but whether one downstream shell displacement is read with opposite and nearly equal gains by the phase and expansion channels.
\end{remark}


\subsection*{Shell-conjugate observer geometry and induced gain balance}
The single-shell-coordinate reduction showed that the anti-phase corridor is natural exactly when the induced gains satisfy $\Gamma_{\mathrm{exp}}\approx-\Gamma_{\mathrm{ph}}$. The next structural step is to push that gain relation back to the observer-shell architecture itself. Rather than treating $\Gamma_{\mathrm{ph}}$ and $\Gamma_{\mathrm{exp}}$ as free coefficients, we ask whether they arise as two downstream readouts of one shell tangent under a shell involution inherited from the umkehrwalze. If so, then anti-phase is no longer a tuning accident between two unrelated gains, but the first-order signature of one shell-conjugate observer geometry.

\begin{definition}[Working shell-conjugate observer response]
Let $S_{\mathrm{obs}}$ be the downstream observer-shell tangent space. A \emph{shell-conjugate observer response} consists of
\begin{enumerate}[leftmargin=2em,label=(\roman*)]
\item an involution $R_{\mathrm{sh}}:S_{\mathrm{obs}}\to S_{\mathrm{obs}}$ with $R_{\mathrm{sh}}^2=\mathrm{id}$, interpreted as the observer-shell descendant of the umkehrwalze,
\item a shell tangent family $u_{\mathrm{sh}}(\theta)\in S_{\mathrm{obs}}$,
\item a base shell readout functional $\ell_{\mathrm{sh}}(\theta)\in S_{\mathrm{obs}}^*$,
\item and a shell-conjugacy defect $\varepsilon_{\mathrm{conj}}(\theta)$,
\end{enumerate}
with
\[
  a_{\mathrm{sh}}(\theta)=\ell_{\mathrm{sh}}(\theta)\bigl(u_{\mathrm{sh}}(\theta)\bigr),
\]
\[
  b_{\mathrm{sh}}(\theta)=-\ell_{\mathrm{sh}}(\theta)\bigl(R_{\mathrm{sh}}u_{\mathrm{sh}}(\theta)\bigr)+\varepsilon_{\mathrm{conj}}(\theta).
\]
If moreover
\[
  R_{\mathrm{sh}}u_{\mathrm{sh}}(\theta)=u_{\mathrm{sh}}(\theta),
\]
then the phase and expansion shell readouts are said to be read from the same shell tangent with conjugate sign, up to the defect $\varepsilon_{\mathrm{conj}}(\theta)$.
\end{definition}

\begin{definition}[Working transport-balance defect]
Define the downstream transport-balance defect by
\[
  \Delta_{AB}(\theta):=B_{\mathrm{sh}}(\theta)-A_{\mathrm{sh}}(\theta).
\]
Exact transport balance on the downstream shell means $\Delta_{AB}(\theta)=0$.
\end{definition}

\begin{theorem}[Working shell-conjugate observer geometry makes anti-phase natural]
Assume the single-shell-coordinate observer law and a shell-conjugate observer response with
\[
  R_{\mathrm{sh}}u_{\mathrm{sh}}(\theta)=u_{\mathrm{sh}}(\theta).
\]
Then
\[
  b_{\mathrm{sh}}(\theta)=-a_{\mathrm{sh}}(\theta)+\varepsilon_{\mathrm{conj}}(\theta),
\]
and hence the induced gains satisfy
\[
  \Gamma_{\mathrm{exp}}(\theta)+\Gamma_{\mathrm{ph}}(\theta)
  =A_{\mathrm{sh}}(\theta)\,\varepsilon_{\mathrm{conj}}(\theta)
   +\Delta_{AB}(\theta)\,b_{\mathrm{sh}}(\theta).
\]
In particular, if $\Gamma_{\mathrm{ph}}(\theta)\neq 0$, then
\[
  \zeta_{\Gamma}(\theta)+1
  =
  \frac{A_{\mathrm{sh}}(\theta)\,\varepsilon_{\mathrm{conj}}(\theta)
   +\Delta_{AB}(\theta)\,b_{\mathrm{sh}}(\theta)}{\Gamma_{\mathrm{ph}}(\theta)}.
\]
Thus the anti-phase corridor $\zeta_{\Gamma}(\theta)\approx -1$ is structurally natural whenever both the shell-conjugacy defect and the downstream transport-balance defect are small relative to $\Gamma_{\mathrm{ph}}(\theta)$.
\end{theorem}

\begin{proof}
By shell-conjugacy,
\[
  b_{\mathrm{sh}}(\theta)
  =-\ell_{\mathrm{sh}}(\theta)\bigl(R_{\mathrm{sh}}u_{\mathrm{sh}}(\theta)\bigr)+\varepsilon_{\mathrm{conj}}(\theta)
  =-\ell_{\mathrm{sh}}(\theta)\bigl(u_{\mathrm{sh}}(\theta)\bigr)+\varepsilon_{\mathrm{conj}}(\theta)
  =-a_{\mathrm{sh}}(\theta)+\varepsilon_{\mathrm{conj}}(\theta).
\]
Therefore
\[
  \Gamma_{\mathrm{exp}}(\theta)+\Gamma_{\mathrm{ph}}(\theta)
  =B_{\mathrm{sh}}(\theta)b_{\mathrm{sh}}(\theta)+A_{\mathrm{sh}}(\theta)a_{\mathrm{sh}}(\theta)
\]
\[
  =B_{\mathrm{sh}}(\theta)\bigl(-a_{\mathrm{sh}}(\theta)+\varepsilon_{\mathrm{conj}}(\theta)\bigr)+A_{\mathrm{sh}}(\theta)a_{\mathrm{sh}}(\theta)
\]
\[
  =(A_{\mathrm{sh}}(\theta)-B_{\mathrm{sh}}(\theta))a_{\mathrm{sh}}(\theta)+B_{\mathrm{sh}}(\theta)\varepsilon_{\mathrm{conj}}(\theta).
\]
Since $b_{\mathrm{sh}}=-a_{\mathrm{sh}}+\varepsilon_{\mathrm{conj}}$, one may rewrite this exactly as
\[
  \Gamma_{\mathrm{exp}}(\theta)+\Gamma_{\mathrm{ph}}(\theta)
  =A_{\mathrm{sh}}(\theta)\,\varepsilon_{\mathrm{conj}}(\theta)+\Delta_{AB}(\theta)b_{\mathrm{sh}}(\theta).
\]
Division by $\Gamma_{\mathrm{ph}}(\theta)$ gives the formula for $\zeta_{\Gamma}(\theta)+1$.
\end{proof}

\begin{corollary}[Working exact shell conjugacy and balanced transport force first-order anti-phase]
Under the same hypotheses, if
\[
  \varepsilon_{\mathrm{conj}}(\theta)=0,
  \qquad
  \Delta_{AB}(\theta)=0,
\]
then
\[
  \Gamma_{\mathrm{exp}}(\theta)=-\Gamma_{\mathrm{ph}}(\theta),
  \qquad
  \zeta_{\Gamma}(\theta)=-1.
\]
Hence the anti-phase corridor is then not merely admissible but exact at first order.
\end{corollary}

\begin{proof}
This is the immediate specialization of the theorem to vanishing shell-conjugacy and transport-balance defects.
\end{proof}

\begin{remark}[Working structural reading of the shell-gain pair]
The gain pair $(\Gamma_{\mathrm{ph}},\Gamma_{\mathrm{exp}})$ is no longer best read as two unrelated coefficients. In the shell-conjugate regime it is the image of one shell tangent under one observer-shell involution and one downstream transport imbalance. Therefore the real structural question is not ``why should two gains accidentally cancel?'', but rather ``how far is the shell-conjugate observer geometry from exact sign-conjugacy and balanced transport?''
\end{remark}


\subsection*{Shell anti-phase defect budget and epsilon-trap}
The shell-conjugate observer geometry showed that anti-phase is natural whenever the shell-conjugacy and transport-balance defects are small. The next step is to compress those two structural errors and the single-shell-coordinate remainder into one normalized anti-phase budget. This turns the visible shell question into an epsilon-trap statement: once the induced shell ratio enters a sufficiently small corridor around $-1$, the first-order visible mode cannot escape that corridor unless the higher-order shell remainder is of comparable size.

\begin{definition}[Working normalized shell anti-phase budget]
Assume $\Gamma_{\mathrm{ph}}(\theta)\neq 0$. Define the normalized shell-conjugacy, transport-balance, and single-shell linearization loads by
\[
  \rho_{\mathrm{conj}}(\theta)
  :=
  \frac{|A_{\mathrm{sh}}(\theta)\,\varepsilon_{\mathrm{conj}}(\theta)|}{|\Gamma_{\mathrm{ph}}(\theta)|},
\]
\[
  \rho_{AB}(\theta)
  :=
  \frac{|\Delta_{AB}(\theta)\,b_{\mathrm{sh}}(\theta)|}{|\Gamma_{\mathrm{ph}}(\theta)|},
\]
\[
  \rho_{\xi}(\theta)
  :=
  \left|\frac{\zeta_{\mathrm{sh}}(\theta)-\zeta_{\Gamma}(\theta)}{1+|\zeta_{\Gamma}(\theta)|}\right|.
\]
The total normalized shell anti-phase budget is
\[
  \rho_{\mathrm{trap}}(\theta)
  :=
  \rho_{\mathrm{conj}}(\theta)+\rho_{AB}(\theta)+\rho_{\xi}(\theta).
\]
\end{definition}

\begin{theorem}[Working shell anti-phase epsilon-trap]
\label{thm:working-shell-anti-phase-epsilon-trap}
Assume the single-shell-coordinate observer law and shell-conjugate observer geometry with
\[
  R_{\mathrm{sh}}u_{\mathrm{sh}}(\theta)=u_{\mathrm{sh}}(\theta),
  \qquad
  \Gamma_{\mathrm{ph}}(\theta)\neq 0.
\]
Then
\[
  |\zeta_{\Gamma}(\theta)+1|
  \le
  \rho_{\mathrm{conj}}(\theta)+\rho_{AB}(\theta),
\]
and hence
\[
  |\zeta_{\mathrm{sh}}(\theta)+1|
  \le
  \rho_{\mathrm{trap}}(\theta)\,\bigl(1+|\zeta_{\Gamma}(\theta)|\bigr).
\]
In particular, if for some $\varepsilon_{\mathrm{trap}}>0$ one has
\[
  \rho_{\mathrm{trap}}(\theta)
  \le
  \frac{\varepsilon_{\mathrm{trap}}}{1+|\zeta_{\Gamma}(\theta)|},
\]
then
\[
  |\zeta_{\mathrm{sh}}(\theta)+1|\le \varepsilon_{\mathrm{trap}}.
\]
Thus the shell ratio is forced into an anti-phase epsilon-tube around $-1$ by one normalized structural budget.
\end{theorem}

\begin{remark}[Reading of the shell epsilon-trap after OP3 localization]
After the OP3 shell/residual theorem, the anti-phase epsilon-trap is best read as a non-escape criterion for the visible shell ratio inside one OP3-local corridor. By itself it does not yet prove global OP closure or global branch admissibility; rather, it says that once the shell-side refinement stays inside a sufficiently small anti-phase budget, the visible first-order shell mode cannot leave that corridor without a comparably sized higher-order defect.
\end{remark}

\begin{proof}
The shell-conjugate observer theorem gave
\[
  \zeta_{\Gamma}(\theta)+1
  =
  \frac{A_{\mathrm{sh}}(\theta)\,\varepsilon_{\mathrm{conj}}(\theta)
   +\Delta_{AB}(\theta)\,b_{\mathrm{sh}}(\theta)}{\Gamma_{\mathrm{ph}}(\theta)}.
\]
Taking absolute values yields
\[
  |\zeta_{\Gamma}(\theta)+1|
  \le
  \frac{|A_{\mathrm{sh}}(\theta)\,\varepsilon_{\mathrm{conj}}(\theta)|}{|\Gamma_{\mathrm{ph}}(\theta)|}
  +
  \frac{|\Delta_{AB}(\theta)\,b_{\mathrm{sh}}(\theta)|}{|\Gamma_{\mathrm{ph}}(\theta)|}
  =
  \rho_{\mathrm{conj}}(\theta)+\rho_{AB}(\theta).
\]
Now
\[
  |\zeta_{\mathrm{sh}}(\theta)+1|
  \le
  |\zeta_{\mathrm{sh}}(\theta)-\zeta_{\Gamma}(\theta)|
  +
  |\zeta_{\Gamma}(\theta)+1|.
\]
By the definition of $\rho_{\xi}(\theta)$,
\[
  |\zeta_{\mathrm{sh}}(\theta)-\zeta_{\Gamma}(\theta)|
  =
  \rho_{\xi}(\theta)\,\bigl(1+|\zeta_{\Gamma}(\theta)|\bigr).
\]
Since $1\le 1+|\zeta_{\Gamma}(\theta)|$, the previous bound implies
\[
  |\zeta_{\Gamma}(\theta)+1|
  \le
  \bigl(\rho_{\mathrm{conj}}(\theta)+\rho_{AB}(\theta)\bigr)
  \bigl(1+|\zeta_{\Gamma}(\theta)|\bigr).
\]
Combining the two estimates gives
\[
  |\zeta_{\mathrm{sh}}(\theta)+1|
  \le
  \rho_{\mathrm{trap}}(\theta)\,\bigl(1+|\zeta_{\Gamma}(\theta)|\bigr).
\]
The epsilon-tube statement is immediate.
\end{proof}

\begin{corollary}[Working exact visible closure reduces to one shell-budget comparison]
Assume the anti-phase sign regime and $\Gamma_{\mathrm{ph}}(\theta)\neq 0$. If exact visible closure holds, then necessarily
\[
  \rho_{\mathrm{trap}}(\theta)
  \le
  \frac{\rho_{\mathrm{rem}}(\theta)
  \bigl(1+|\zeta_{\mathrm{sh}}(\theta)|\bigr)}{1+|\zeta_{\Gamma}(\theta)|}.
\]
Conversely, if
\[
  \rho_{\mathrm{trap}}(\theta)
  >
  \frac{\rho_{\mathrm{rem}}(\theta)
  \bigl(1+|\zeta_{\mathrm{sh}}(\theta)|\bigr)}{1+|\zeta_{\Gamma}(\theta)|},
\]
then exact visible closure is impossible and the shared shell mode is canonically nonzero.
\end{corollary}

\begin{proof}
The anti-phase corridor theorem for visible closure gave the necessary condition
\[
  |1+\zeta_{\mathrm{sh}}(\theta)|
  \le
  \rho_{\mathrm{rem}}(\theta)\,\bigl(1+|\zeta_{\mathrm{sh}}(\theta)|\bigr).
\]
The shell anti-phase epsilon-trap gives
\[
  |1+\zeta_{\mathrm{sh}}(\theta)|
  \le
  \rho_{\mathrm{trap}}(\theta)\,\bigl(1+|\zeta_{\Gamma}(\theta)|\bigr).
\]
If exact visible closure is to be possible, the first-order shell mismatch forced by the second estimate must fit inside the higher-order rescue corridor from the first. This yields the displayed necessary inequality. Its negation rules out exact visible closure, hence forces a canonically nonzero shared shell mode.
\end{proof}

\begin{remark}[Working structural reading of the shell anti-phase budget]
The anti-phase corridor is no longer best read as an unexplained near-accident between two signed shell loads. After the shell-conjugate reduction, the visible first-order defect is controlled by one budget only: shell-conjugacy error, transport-balance error, and the residual error from replacing the induced shell ratio by the gain ratio. Therefore the next structural question is not ``can phase and expansion accidentally cancel?'', but rather ``is the combined shell anti-phase budget genuinely smaller than the higher-order rescue scale?''
\end{remark}



\subsection*{3.10.31 Shell-balance driver as an architecture defect}
The remaining first-order shell driver should not be treated as an independent phenomenological parameter. It is more natural to regard it as the visible shadow of a joint architecture defect involving the active HC/MM/J fit on the admissible branch together with the observer-shell dressing of the same branch.

\begin{definition}[Working architecture-side shell driver loads]
Set
\[
  \lambda_{\mathrm{fit}}(\theta)
  :=
  \frac{\|[(\mathcal M,J)]|_{U}\|\,\|x(\theta)\|}{|\Gamma_{\mathrm{ph}}(\theta)|},
  \qquad
  \lambda_{\mathrm{obs}}(\theta):=\lambda_{\mathrm{sh}}(\theta).
\]
Here \(\lambda_{\mathrm{fit}}\) measures active architecture mismatch through the restricted commutator, while \(\lambda_{\mathrm{obs}}\) is the normalized observer-shell anti-phase load already extracted above.
\end{definition}

\begin{theorem}[Working shell-balance driver is downstream from HC/MM/J fit plus shell dressing]
Assume the active-branch setup from the previous section, with visible shared shell mode
\[
  r_{\mathrm{obs}}(\theta)
  =
  \Xi_{\mathrm{bal}}(\theta)\,\xi_{\mathrm{sh}}(\theta)
  + O\!\left(|\xi_{\mathrm{sh}}(\theta)|^2\right).
\]
Then there is a first-order architecture estimate of the form
\[
  \frac{|\Xi_{\mathrm{bal}}(\theta)|}{|\Gamma_{\mathrm{ph}}(\theta)|}
  \le
  \lambda_{\mathrm{fit}}(\theta)+\lambda_{\mathrm{obs}}(\theta).
\]
In particular, if both
\[
  [\mathcal M,J]\big|_{U}=0
  \qquad\text{and}\qquad
  \lambda_{\mathrm{sh}}(\theta)=0,
\]
then
\[
  \Xi_{\mathrm{bal}}(\theta)=0,
\]
and the shared visible shell mode has no first-order driver.
\end{theorem}

\begin{proof}
By the exact commutator identity on the admissible branch, the primitive
common-mode carrier is controlled by
\[
  E_{\mathrm{fit}}^{(U)}
  = \frac12[\mathcal M,J]\big|_{U}.
\]
Hence every first-order visible common-mode contribution inherited from the
active transport branch is bounded, up to the visible readout coefficient
entering \(\Gamma_{\mathrm{ph}}\), by a multiple of
\(\|[\mathcal M,J]|_{U}\|\,\|x(\theta)\|\). Normalizing by
\(|\Gamma_{\mathrm{ph}}(\theta)|\) gives the \(\lambda_{\mathrm{fit}}\) term.

The remaining first-order imbalance after exact admissible transport was already compressed into the shell-conjugacy and transport-balance loads, combined in \(\lambda_{\mathrm{sh}}=\lambda_{\mathrm{conj}}+\lambda_{\mathrm{bal}}\). These are exactly the observer-shell contributions entering the shell-balance driver \(\Xi_{\mathrm{bal}}\). Adding the architecture-side and observer-side contributions yields
\[
  \frac{|\Xi_{\mathrm{bal}}(\theta)|}{|\Gamma_{\mathrm{ph}}(\theta)|}
  \le
  \lambda_{\mathrm{fit}}(\theta)+\lambda_{\mathrm{sh}}(\theta).
\]
If both terms vanish, then the normalized driver vanishes and therefore \(\Xi_{\mathrm{bal}}(\theta)=0\).
\end{proof}

\begin{corollary}[Working first-order visible cancellation is architecture-driven]
If
\[
  \lambda_{\mathrm{fit}}(\theta)+\lambda_{\mathrm{sh}}(\theta)<1,
\]
then the shared shell mode is forced into a first-order small-driver regime. If
\[
  \lambda_{\mathrm{fit}}(\theta)+\lambda_{\mathrm{sh}}(\theta)=0,
\]
then any remaining visible residual is necessarily of strictly higher observer-shell order.
\end{corollary}

\begin{proof}
The theorem gives
\[
  |\Xi_{\mathrm{bal}}(\theta)|
  \le
  |\Gamma_{\mathrm{ph}}(\theta)|\bigl(\lambda_{\mathrm{fit}}(\theta)+\lambda_{\mathrm{sh}}(\theta)\bigr).
\]
The first claim is the corresponding small-driver statement. If the sum vanishes, then \(\Xi_{\mathrm{bal}}(\theta)=0\), and the first-order shell law reduces to
\[
  r_{\mathrm{obs}}(\theta)=O\!\left(|\xi_{\mathrm{sh}}(\theta)|^2\right)+R_{\ge 2}(\theta),
\]
so every visible remainder is of higher order.
\end{proof}

\begin{remark}[Working architectural reading of the remaining shell law]
The visible common-mode residual is now best read as the downstream image of one combined architectural budget. The active HC/MM/J fit supplies the primitive common-mode leakage budget, while the shell-conjugacy and transport-balance terms supply the observer-side dressing budget. In this reading, the next real question is no longer whether OP~1 and OP~5 are separate closures, but whether the combined architecture budget vanishes exactly or leaves a canonical higher-order observer-shell remainder.
\end{remark}


\subsection*{3.10.32 Final architecture-versus-rescue budget corridor}
After the shell-balance driver has been pushed back to the combined architecture budget, the remaining visible question should no longer be read as an independent shell law. It is now a final corridor comparison: can the higher-order rescue sector absorb the first-order architecture load, or does a canonically nonzero visible shell mode necessarily survive?

\begin{definition}[Working normalized architecture budget and quadratic shell truncation]
Set
\[
  \rho_{\mathrm{arch}}(\theta)
  :=
  \frac{|\Gamma_{\mathrm{ph}}(\theta)\,\xi_{\mathrm{sh}}(\theta)|}{\Lambda_{\mathrm{lin}}(\theta)}
  \bigl(\lambda_{\mathrm{fit}}(\theta)+\lambda_{\mathrm{obs}}(\theta)\bigr),
\]
and define the normalized quadratic shell truncation term by
\[
  \rho_{\xi^2}(\theta)
  :=
  \frac{|O(|\xi_{\mathrm{sh}}(\theta)|^2)|}{\Lambda_{\mathrm{lin}}(\theta)}.
\]
The first quantity measures the visible first-order architecture load inherited from the combined HC/MM/J-fit plus observer-shell budget; the second measures the normalized first neglected shell order.
\end{definition}

\begin{theorem}[Working final architecture-rescue corridor under sharp active loading]
Assume the sharp active-loading regime on a window where
\[
  |\Gamma_{\mathrm{ph}}(\theta)\,\xi_{\mathrm{sh}}(\theta)|
  \ge c_0\,\Lambda_{\mathrm{lin}}(\theta)
\]
for some fixed constant \(c_0>0\). Then the normalized first-order visible shell defect satisfies
\[
  |\rho_{\mathrm{sh}}(\theta)|
  \ge c_0\,\rho_{\mathrm{arch}}(\theta)-\rho_{\xi^2}(\theta).
\]
Consequently, exact visible closure can hold only if
\[
  c_0\,\rho_{\mathrm{arch}}(\theta)
  \le
  \rho_{\mathrm{rem}}(\theta)+\rho_{\xi^2}(\theta).
\]
\end{theorem}

\begin{proof}
By the architecture-driven shell law from the previous subsection,
\[
  r_{\mathrm{obs}}(\theta)
  =
  \Xi_{\mathrm{bal}}(\theta)\,\xi_{\mathrm{sh}}(\theta)
  + O\!\left(|\xi_{\mathrm{sh}}(\theta)|^2\right)
  + R_{\ge 2}(\theta),
\]
with
\[
  \frac{|\Xi_{\mathrm{bal}}(\theta)|}{|\Gamma_{\mathrm{ph}}(\theta)|}
  \le
  \lambda_{\mathrm{fit}}(\theta)+\lambda_{\mathrm{obs}}(\theta).
\]
Dividing the first-order term by \(\Lambda_{\mathrm{lin}}(\theta)\) yields the normalized architecture load \(\rho_{\mathrm{arch}}(\theta)\). Under the sharp loading hypothesis, the visible contribution of this term is bounded below by \(c_0\rho_{\mathrm{arch}}(\theta)\) up to the normalized quadratic truncation error \(\rho_{\xi^2}(\theta)\). This proves
\[
  |\rho_{\mathrm{sh}}(\theta)|
  \ge c_0\,\rho_{\mathrm{arch}}(\theta)-\rho_{\xi^2}(\theta).
\]
If exact visible closure is to be possible, then by the earlier normalized shell-defect theorem one must have
\[
  |\rho_{\mathrm{sh}}(\theta)|\le \rho_{\mathrm{rem}}(\theta).
\]
Combining both inequalities gives the displayed necessary corridor condition.
\end{proof}

\begin{corollary}[Working final reduced decision after exact transport]
Under exact admissible transport and sharp active loading, the visible OP~1/OP~5 question is reduced to a single final budget comparison:
\[
  c_0\,\rho_{\mathrm{arch}}(\theta)
  \stackrel{?}{\le}
  \rho_{\mathrm{rem}}(\theta)+\rho_{\xi^2}(\theta).
\]
If the inequality fails, then the common visible shell mode is canonically nonzero. If it holds, then exact visible closure remains possible only through higher-order rescue of a first-order architecture load.
\end{corollary}

\begin{remark}[Working interpretation of the final corridor]
The remaining visible problem is now compressed as far as the current structure allows. It is no longer a pair of unfinished closures, nor even two shell loads. It is one corridor question: can the higher-order rescue sector absorb the normalized combined architecture load after exact transport, or does a canonical observer-shell remainder necessarily survive?
\end{remark}



\subsection{QAM-over-ZF: operator families and lambda-like structural maps}

We now close the first internal QAM package by adding an operator-family and lambda-like structural layer on top of the already introduced transport, branch/shell, tensor, and translation strata. The point remains conservative: no higher-order foundation is introduced. Instead, families of coded operators and lambda-like structural maps are themselves represented as set-codes over ZF and are required to preserve closure content on the admissible fragment.

\begin{definition}[Operator-family code]
A set
\[
\mathfrak{o}=\langle 9,I,F\rangle
\]
is called an \emph{operator-family code} if $I$ is an index code and $F$ is a graph-code assigning to each $i\in I$ a coded operator $o_i$.
\end{definition}

\begin{definition}[Admissible operator family]
An operator-family code
\[
\mathfrak{o}=\langle 9,I,F\rangle
\]
is called \emph{admissible} if every member operator $o_i$ coded by $F$ is admissible on its intended coded state domain. Equivalently, for every admissible coded source state $s$ and every coded target state $t$ linked by some member $o_i$, one has
\[
\mathrm{QAdmTransport}(o_i,s,t)
\Longrightarrow
\mathrm{QClEq}(s,t).
\]
\end{definition}

\begin{remark}
Thus admissibility of an operator family means pointwise closure preservation across the family. The family layer does not weaken the closure-first discipline already imposed on single coded operators.
\end{remark}

\begin{definition}[Lambda-like structural map code]
A set
\[
\lambda=\langle 10,D,C,L\rangle
\]
is called a \emph{lambda-like structural map code} if $D$ codes an intended domain class of QAM-coded objects, $C$ codes an intended codomain class, and $L$ is a graph-code assigning to each coded object in $D$ a coded object in $C$.
\end{definition}

\begin{definition}[Closure-compatible lambda map]
A lambda-like structural map code $\lambda$ is called \emph{closure-compatible} on a coded state domain if for all coded states $s,t$ in that domain,
\[
\mathrm{QClEq}(s,t)
\Longrightarrow
\mathrm{QClEq}(\lambda(s),\lambda(t)),
\]
whenever the values $\lambda(s)$ and $\lambda(t)$ are coded states.
\end{definition}

\begin{remark}
This is the lambda-level analogue of admissible transport. A closure-compatible lambda map need not itself be read as a transport along one fixed route; it may instead be a structural selector, reindexer, collector, or support-transforming map, provided it preserves closure class on the admissible fragment.
\end{remark}

\begin{definition}[Support-selection lambda]
A closure-compatible lambda-like structural map is called a \emph{support-selection lambda} if it acts on admissible branch-shell support codes and returns a coded sub-support, preferred support, or selected routed support without leaving the admissible closure regime.
\end{definition}

\begin{proposition}[Pointwise closure preservation for admissible operator families]
Let
\[
\mathfrak{o}=\langle 9,I,F\rangle
\]
be an admissible operator family. If $o_i$ is a member of the family and
\[
\mathrm{QAdmTransport}(o_i,s,t),
\]
then
\[
\mathrm{QClEq}(s,t).
\]
Consequently, every closure-compatible readout identifies the closure-relevant content of $s$ and $t$ after factorization.
\end{proposition}

\begin{proof}
This is immediate from pointwise admissibility of the family members together with admissibility-implies-closure preservation and closure-compatible readout.
\end{proof}

\begin{proposition}[Closure preservation under closure-compatible lambda maps]
Let $\lambda$ be a closure-compatible lambda-like structural map on an admissible coded state domain. If
\[
\mathrm{QClEq}(s,t),
\]
then
\[
\mathrm{QClEq}(\lambda(s),\lambda(t)),
\]
whenever the values are coded states.
\end{proposition}

\begin{proof}
This is exactly the defining property of closure-compatible lambda maps.
\end{proof}

\begin{theorem}[First QAM operator-family/lambda theorem]
Assume:
\begin{enumerate}
\item $\mathfrak{o}=\langle 9,I,F\rangle$ is an admissible operator family;
\item $\lambda=\langle 10,D,C,L\rangle$ is a closure-compatible lambda-like structural map on the relevant admissible fragment;
\item $r$ is a closure-compatible coded readout.
\end{enumerate}
Then the following hold.
\begin{enumerate}
\item \textbf{Familywise transport preservation.}  
For every member operator $o_i$ of $\mathfrak{o}$ and every admissible transport step
\[
\mathrm{QAdmTransport}(o_i,s,t),
\]
one has
\[
\mathrm{QClEq}(s,t).
\]

\item \textbf{Lambda-level closure preservation.}  
If in addition
\[
\mathrm{QClEq}(s,t),
\]
then
\[
\mathrm{QClEq}(\lambda(s),\lambda(t)).
\]

\item \textbf{Readout stability after family transport and lambda action.}  
Whenever the composed expressions are defined on the admissible fragment, the closure-relevant readout content of
\[
s,
\quad t,
\quad \lambda(s),
\quad \lambda(t)
\]
agrees after factorization through the associated class-target maps.
\end{enumerate}
\end{theorem}

\begin{proof}
The first claim is pointwise admissibility of the operator family. The second is closure preservation under the lambda-like map. The third follows because closure-compatible readout factors through coded closure class, and all four coded states listed in the statement lie in the same closure class whenever the hypotheses apply.
\end{proof}

\begin{corollary}[Operator-family/lambda lift of the admissible OP fragment]
The admissible OP fragment translated into QAM-over-ZF remains stable not only under single admissible coded operators, but also under admissible operator families and closure-compatible lambda-like structural maps acting on the translated fragment.
\end{corollary}

\begin{remark}
This is the first point at which the QAM layer becomes more than a static coding scheme. It now supports controlled families of operations and closure-preserving structural maps on the admissible fragment while remaining conservative over ZF.
\end{remark}

\begin{remark}[Transition to support selection]
The preceding operator-family/lambda layer is the last part of the internal QAM package that still acts directly on admissible transport and closure data. The next layer is different in status: support selection and degeneracy handling do not generate admissibility, but only refine already admissible branch-shell support after the closure-first filters have been applied.
\end{remark}

\subsection{QAM-over-ZF: support selection and degeneracy handling}

We now add a first support-selection layer on top of the already translated admissible QAM fragment. The purpose of this layer is not to replace closure-admissibility or shell-admissibility, but to order admissible support after those primary filters have already been applied. Thus QAM selection is subordinate to closure preservation and only acts inside the admissible branch-shell sector.

\begin{definition}[Support candidate code]
A set
\[
\varsigma = \langle 11,\Omega,w\rangle
\]
is called a \emph{support candidate code} if $\Omega$ codes a family of admissible branch-shell pairs and $w$ codes a weight or preference assignment on $\Omega$.
\end{definition}

\begin{remark}
A support candidate code is the QAM analogue of a visible support sector restricted in advance to admissible branch-shell data. It is therefore a post-admissibility object rather than a primitive closure object.
\end{remark}

\begin{definition}[Support preference relation]
\[
\begin{aligned}
&\text{closure preservation}
\;>\;
\text{branch admissibility}
\;>\;
\text{shell admissibility}\\[0.2em]
&\;>\;
\text{support preference}
\;>\;
\text{visible routed weight}.
\end{aligned}
\]
\end{definition}

\begin{remark}
Thus support preference never overrides admissibility. It only compares already admissible support candidates.
\end{remark}

\begin{definition}[Closure-compatible support selector]
A lambda-like structural map
\[
\lambda_{\mathrm{sel}}=\langle 10,D,C,L\rangle
\]
is called a \emph{closure-compatible support selector} if:
\begin{enumerate}
\item its domain consists of admissible support candidate codes;
\item its values are coded selected supports or coded preferred subsupports;
\item whenever two input support candidates encode the same closure-bearing admissible content up to support degeneracy, their images under $\lambda_{\mathrm{sel}}$ are closure-equivalent or support-equivalent at coded level.
\end{enumerate}
\end{definition}

\begin{definition}[Support degeneracy]
Two admissible support candidate codes $\varsigma_1,\varsigma_2$ are called \emph{support-degenerate} if they differ only by admissible re-indexing, admissible redistribution of coded visible weight, or admissible closure-equivalent regrouping, while preserving the same coded closure-bearing support content.
\end{definition}

\begin{remark}
Support degeneracy is the QAM-over-ZF shadow of carrier degeneracy and readout multiplicity in the current Companion line. It records when several admissible support realizations should still be read as one structural support type.
\end{remark}

\begin{definition}[Selected routed support code]
If $\lambda_{\mathrm{sel}}$ is a closure-compatible support selector and $\varsigma$ is an admissible support candidate code, then
\[
\lambda_{\mathrm{sel}}(\varsigma)
\]
is called the corresponding \emph{selected routed support code}, whenever the value is defined.
\end{definition}

\begin{proposition}[Selection stays inside the admissible QAM fragment]
Let $\lambda_{\mathrm{sel}}$ be a closure-compatible support selector and let $\varsigma$ be an admissible support candidate code. Then every selected routed support code
\[
\lambda_{\mathrm{sel}}(\varsigma)
\]
remains inside the admissible QAM support fragment.
\end{proposition}

\begin{proof}
By definition, the domain of $\lambda_{\mathrm{sel}}$ consists only of admissible support candidate codes, and the selector is required to return coded selected supports or preferred subsupports without leaving the closure-compatible admissible regime. Hence the selected support remains inside the admissible fragment.
\end{proof}

\begin{proposition}[Closure preservation under support selection]
Let $\varsigma$ be an admissible support candidate code and let $\lambda_{\mathrm{sel}}$ be a closure-compatible support selector. Then the closure-relevant content of
\[
\lambda_{\mathrm{sel}}(\varsigma)
\]
agrees with the closure-relevant content of $\varsigma$ after passage through the associated closure-factor map.
\end{proposition}

\begin{proof}
Support selection acts only after admissibility has already been established and is required to be closure-compatible. Therefore the selector cannot change the coded closure-bearing content of the admissible support candidate; it may only refine, reduce, or prefer one admissible representation over another. Closure-compatible readout therefore identifies the closure-relevant content of the input and output supports after factorization.
\end{proof}

\begin{corollary}[Support selection does not create closure]
A selected routed support is closure-preserving only because it is selected from an already admissible support candidate. Selection itself is not an independent source of closure.
\end{corollary}

\begin{remark}
This is the coded QAM analogue of the structure-first rule that preferred support must be read only after closure and admissibility have already been fixed.
\end{remark}

\begin{proposition}[Degenerate supports may be identified after selection]
Let $\varsigma_1,\varsigma_2$ be support-degenerate admissible support candidate codes. If $\lambda_{\mathrm{sel}}$ is a closure-compatible support selector, then
\[
\lambda_{\mathrm{sel}}(\varsigma_1)
\qquad\text{and}\qquad
\lambda_{\mathrm{sel}}(\varsigma_2)
\]
may be treated as equivalent selected supports at the coded structural level.
\end{proposition}

\begin{proof}
Since $\varsigma_1$ and $\varsigma_2$ differ only by admissible re-indexing, admissible redistribution of visible weight, or closure-equivalent regrouping, they carry the same coded closure-bearing support content. Closure-compatibility of the selector ensures that this content is preserved under selection. Therefore their selected images may be identified at the coded structural level.
\end{proof}

\begin{theorem}[First QAM support-selection theorem]
Assume:
\begin{enumerate}
\item $\varsigma$ is an admissible support candidate code;
\item $\lambda_{\mathrm{sel}}$ is a closure-compatible support selector;
\item $r$ is a closure-compatible coded readout on the relevant support fragment.
\end{enumerate}
Then:
\begin{enumerate}
\item the selected routed support code $\lambda_{\mathrm{sel}}(\varsigma)$ remains admissible;
\item the closure-relevant readout content of $\lambda_{\mathrm{sel}}(\varsigma)$ agrees with that of $\varsigma$ after factorization;
\item support-degenerate admissible candidates yield equivalent selected supports at coded structural level.
\end{enumerate}
\end{theorem}

\begin{proof}
The first claim is selection-stays-inside-admissible-fragment. The second is closure preservation under support selection. The third is the degeneracy proposition above. Together they show that support selection in QAM is a post-admissibility refinement layer rather than an independent closure mechanism.
\end{proof}

\begin{remark}
This is the first QAM-over-ZF support-selection layer. It does not yet impose one final universal selector. It only provides the structural place where preferred support, selected support, and degeneracy handling may be coded without violating the closure-first architecture inherited from the current Companion line.
\end{remark}


\begin{remark}[Current closure of the internal QAM package]
At the present internal stage, the QAM-over-ZF line has reached a first connected package state: coded objects, primitive predicates, admissible transport, coded closure-class preservation, closure-compatible readout, branch/shell/routed admissibility, tensor closure rules, faithful translation of the admissible OP fragment, operator-family/lambda-like maps, and support-selection/degeneracy handling are now all present inside one conservative ZF-grounded coding line. The next natural work should therefore focus on normalization, compression, and later public extraction rather than immediate uncontrolled enlargement.
\end{remark}


\clearpage
